% concatenated from main.tex
\documentclass[11pt]{amsart}
\usepackage[
  left=2.5cm,
  right=2.5cm,
  top=3cm,
  bottom=3cm
]{geometry}

\usepackage{amsmath,amssymb,amsthm,graphicx,mathrsfs,url,bbm}
\usepackage[usenames,dvipsnames]{color}
\usepackage[colorlinks=true,linkcolor=Red,citecolor=Green]{hyperref}
\usepackage{amsxtra}
\usepackage{wasysym} % for "\varhexagon" macro
\usepackage{graphicx}% for "\scalebox" macro
\usepackage{subcaption}

\usepackage{framed}


\usepackage{braket}

% arXiv bibliography macro
\def\arXiv#1{\href{http://arxiv.org/abs/#1}{arXiv:#1}}


\def\blue#1{\textcolor{blue}{#1}}
\def\red#1{\textcolor{red}{#1}}
\def\green#1{\textcolor{halfgreen}{#1}}
\def\cyan#1{\textcolor{cyan}{#1}}

\DeclareMathOperator{\const}{const}

\setlength{\marginparwidth}{0.6in}


%\renewcommand{\marginpar}[2][]{}
 

\def\?[#1]{\textbf{[#1]}\marginpar{\Large{\textbf{??}}}}

\def\smallsection#1{\smallskip\noindent\textbf{#1}.}
\let\epsilon=\varepsilon % sorry Knuth

\setlength{\textheight}{8.50in} \setlength{\oddsidemargin}{0.00in}
\setlength{\evensidemargin}{0.00in} \setlength{\textwidth}{6.08in}
\setlength{\topmargin}{0.00in} \setlength{\headheight}{0.18in}
\setlength{\marginparwidth}{1.0in}
\setlength{\abovedisplayskip}{0.2in}
\setlength{\belowdisplayskip}{0.2in}
\setlength{\parskip}{0.05in}
\renewcommand{\baselinestretch}{1.10}
\newcommand{\indentalign}{\hspace{0.3in}&\hspace{-0.3in}}
\newcommand{\RR}{{\mathbb R}}
\newcommand{\NN}{{\mathbb N}}
\newcommand{\CC}{{\mathbb C}}
\newcommand{\TT}{{\mathbb T}}
\newcommand{\ZZ}{{\mathbb Z}}
\newcommand{\Sp}{{\operatorname{Sp}}}
\newtheorem{conditionA}{Condition}
\renewcommand{\theconditionA}{1}
% =================================================
%                 TITLE STYLE
% ==================================

% optional: slightly tighter title bloc

% =================================================
%                 MACROS
% =================================================
\newcommand{\nn}{\nonumber}
\newcommand{\1}{\mathbbm{1}}
\newcommand{\Tr}{\operatorname{Tr}}
\newcommand{\tr}[1]{\operatorname{Tr}\left(#1\right)}
\newcommand{\norm}[1]{\left\lVert#1\right\rVert}
\newcommand{\ceil}[1]{\left\lceil {#1} \right\rceil}

\newcommand{\ketbra}[2]{|#1\rangle\langle#2|}
\newcommand{\kb}[1]{|#1\rangle\!\langle#1|}
\newcommand{\scalar}[2]{\langle#1|#2\rangle}

\setlength{\itemsep}{0pt}
\setlength{\parskip}{0pt}
\setlength{\parsep}{0pt}
\usepackage{amsthm}

% three special one-off environments


\newtheorem{conditionB}{Condition}
\renewcommand{\theconditionB}{B}

\newtheorem{conditionC}{Condition}
\renewcommand{\theconditionC}{C}

\newcommand{\id}{\mathbb{I}}
\newcommand{\R}{\mathbb{R}}
\newcommand{\C}{\mathbb{C}}
\newcommand{\N}{\mathbb{N}}
\newcommand{\Z}{\mathbb{Z}}

\newcommand{\cH}{\mathcal{H}}
\newcommand{\cB}{\mathcal{B}}
\newcommand{\cL}{\mathcal{L}}
\newcommand{\cE}{\mathcal{E}}
\newcommand{\Var}{\mathrm{Var}}
\newcommand{\E}{\mathbb{E}}
\newcommand{\Inf}{\operatorname{Inf}}
\newcommand{\spec}{\operatorname{Sp}}

\newcommand{\eps}{\varepsilon}
\newcommand{\NA}{N_{\mathcal{A}}}
\renewcommand{\epsilon}{\varepsilon}
\renewcommand{\phi}{\varphi}
\renewcommand{\d}{d}
\renewcommand{\L}{\mathcal{L}}
\newcommand{\cA}{\mathcal{A}}
\newcommand{\hatfM}{\widehat{f}_{\scalebox{0.6}{$\mathscr{M}$}}}
\newcommand{\hatfMdelta}{\widehat{f}_{\scalebox{0.6}{$\mathscr{M}$}_\delta}}
\newcommand{\fM}{f_{\scalebox{0.6}{$\mathscr{M}$}}}
\newcommand{\fMdelta}{f_{\scalebox{0.6}{$\mathscr{M}$}_\delta}}
\newcommand{\fMdeltakappa}{f_{\scalebox{0.6}{$\mathscr{M}$}_{\delta,\kappa}}}

\newcommand{\Ntot}{N_{\operatorname{tot}}}
\newcommand{\bin}[1]{}
\newcommand{\rob}[1]{\textcolor{cyan}{[rob: #1]}}
\newcommand{\robi}[1]{\textcolor{cyan}{#1}}

% =================================================
%                 THEOREMS
% =================================================
\theoremstyle{plain}
\newtheorem{theorem}{Theorem}[section]
\newtheorem{lemma}[theorem]{Lemma}
\newtheorem{proposition}[theorem]{Proposition}
\newtheorem{corollary}[theorem]{Corollary}
\newtheorem{conjecture}[theorem]{Conjecture}
\newtheorem{hypothesis}[theorem]{Hypothesis}
\newtheorem{remark}[theorem]{Remark}
\newcommand{\gap}{\mathrm{gap}}


\theoremstyle{definition}
\newtheorem{definition}[theorem]{Definition}
\newtheorem{assumption}[theorem]{Assumption}
\newtheorem{notation}[theorem]{Notation}
\newtheorem{condition}[theorem]{Condition}
\newtheorem{problem}[theorem]{Problem}

\theoremstyle{remark}
\newtheorem{example}[theorem]{Example}
\newtheorem{claim}{Claim}[theorem]

\numberwithin{equation}{section}

\begin{document}
% =================================================
%                 TITLE DATA
% =================================================
\title[Gibbs sampling for Coulomb quantum gases]{Computing the free energy of quantum Coulomb gases and molecules via quantum Gibbs sampling}

\author{Simon Becker}
\email{simon.becker@unibocconi.it}
\address{Uni Bocconi, 20136, Milan, Italy}



\author{Cambyse Rouz\'e}
\email{cambyse.rouze@inria.fr}
\address{Inria, T\'el\'ecom Paris -- LTCI, Institut Polytechnique de Paris, 91120 Palaiseau, France}

\author{Robert Salzmann}
\email{robert.salzmann@rwth-aachen.de}
\address{
Institute for Quantum Information, RWTH Aachen University, Germany}




% your text starts here



\begin{abstract}
We develop a quantum algorithm for estimating the free energy as well as the total Gibbs state of interacting quantum Coulomb gases and molecular systems in dimensions $d \in \{2,3\}$ at finite temperature. These systems lie beyond the reach of existing methods due to their singular interactions and infinite-dimensional Hilbert space structure. First, we show that the free energy of the full many-body Hamiltonian can be approximated by that of the same Hamiltonian with a finite-rank low-energy truncation of the interaction, with an explicit error bound polynomial in the particle number. This reduces the problem to a controlled finite-rank perturbation problem. Second, we introduce a quantum Gibbs sampling scheme tailored to this truncated system, based on a class of quantum Markov semigroups. Our main analytical result establishes that the associated generator has a strictly positive spectral gap for every truncation, implying exponential convergence to the target Gibbs state. This provides, to our knowledge, the first rigorous mixing-time guarantee for Gibbs sampling in a Coulomb interacting continuous-variable quantum system. Finally, we give an explicit quantum circuit implementation of the dynamics and derive an end-to-end complexity bound for approximating the free energy and the Gibbs state itself. Our results provide a mathematically rigorous route to quantum algorithms for free energy estimation in interacting quantum systems, without relying on classical approximations such as the Born–Oppenheimer reduction.

\end{abstract}


\
\maketitle

\vspace{-1,2cm}



 \tableofcontents




% =========================
% Suggested Introduction
% =========================
\section{Introduction}


\noindent Free energy plays a foundational role in thermodynamics and has become an essential tool for describing complex chemical, biological, and physical phenomena.
In practice, the free-energy landscape connecting an initial configuration to a final one governs the likelihood and timescale of transitions between states: large free-energy barriers strongly suppress such transitions, whereas lower barriers make them more accessible. A familiar example is provided by chemical reactions in solution, including processes occurring inside living cells, where reaction rates are often controlled by the need to overcome intermediate free-energy obstacles. Similar ideas also arise in quantum settings, where free-energy landscapes help characterize thermal states of many-body systems, metastability, or the difficulty of steering a quantum system from one phase or configuration to another.



\subsection{General setting} This paper is concerned with the estimation of the free energy of a class of Schr\"{o}dinger operators modeling a system of $n$ $d$-dimensional quantum particles, $d\in\{2,3\}$, interacting via Coulomb interactions and each experiencing a trapping potential. The Hamiltonian of the system takes the form:
\begin{align*}
H_n=H_{0,n}+W_n,\qquad \text{ with }\qquad  H_{0,n}:=-\Delta+|x|^2,
\end{align*}
where $\Delta$ stands for $d\cdot n$-dimensional Laplacian, $|x|^2=\sum_{i\in[n]}\sum_{j\in [d]}x_{ij}^2$ and 
\[
W_n:=\sum_{1\le i<j\le n}\alpha_{n,i,j}\,W_{ij}\]
is defined as a form on $L^2((\mathbb{R}^d)^n)$, with $\alpha_{n,i,j}\in\R$   and 
\begin{align*}
W_{ij}[\psi]:=\int_{\R^{nd}}w_d(x_i-x_j)|\psi(x)|^2 \ dx,\qquad
w_d(y):=
\begin{cases}
-\log|y|, & d=2,\\[0.3em]
|y|^{-1}, & d=3.
\end{cases}
\end{align*}
We allow for the possibility that the interaction strengths \(\alpha_{n,i,j}\) scale with the number of particles \(n\). In principle, however, one may also choose \(\alpha_{n,i,j}\) to be independent of \(n\) and account only for the particle charges.

The quadratic trap provides a coercive external potential that grows at infinity and suppresses
the escape of particles. At finite temperature, the latter guaranties the Gibbs hypothesis, namely that the Gibbs operator $\mathrm e^{-\beta H_n}$ is trace class for \(\beta>0\), so that
\begin{align*}
\sigma_\beta(H_n):=\frac{e^{-\beta H_n}}{\mathcal{Z}_\beta(H_n)}\qquad \text{ with }\qquad \mathcal{Z}_\beta(H_n):=\Tr(e^{-\beta H_n})
\end{align*}
is a well-defined Gibbs state on $L^2((\mathbb{R}^d)^n)$. Notably, the Hamiltonian $H_n$
is bounded from below\footnote{In this setting, \(h_0\) increases polynomially with the particle number \(n\). For simplicity, we nevertheless keep the notation \(h_0\).} $H_n\ge -h_0 1$, $h_0\ge 0$. The free energy of the system is described by the functional 
\begin{align*}
F_\beta(H_n):=-\frac{1}{\beta}\log\mathcal{Z}_\beta(H_n).
\end{align*}

\subsection{Hamiltonian truncation}
\label{sec:HamiltonianTruncation}
Our first result establishes that the free energy of \(H_n\) is well approximated by a finite-rank low-energy truncation. This reduction is the starting point for the Gibbs sampling scheme constructed later in the paper to estimate the free energy. Concretely, for \(M\in\mathbb N\), let \(P_M\) be the spectral projector onto the low-energy one-body subspace spanned by the eigenfunctions corresponding to the first \(M\) eigenvalues \((\lambda_k)_{k=1}^M\) of $h:=-\Delta+|x|^2
\quad\text{on }L^2(\mathbb R^d)$, and define the corresponding \(n\)-body projector $\Pi_M:=P_M^{\otimes n}$. We then consider the truncated form 
\begin{align*}
W_{n,M}[\psi]:=W_n[\Pi_M\psi],\qquad \text{ and }\qquad H_{n,M}:=H_{0,n}+W_{n,M}\,.
\end{align*}
The Hamiltonian $H_{n,M}$ admits a finite free energy $F_{\beta}(H_{n,M})$ and Gibbs state $\sigma_\beta(H_{n,M})$. In Section~\ref{perturbationpartitionfunction}, we prove that this free energy efficiently approximates that of the untruncated model:
\begin{theorem}[Free energy perturbation bound]\label{thmfreenregyblablaIntro}
Fix $d\in\{2,3\}$, $\beta>0$. There exists a constant $C_\beta<\infty$, independent of $n,M\ge 1$, such that
\[
|F_\beta(H_n)-F_{\beta}(H_{n,M})|
\le C_\beta\,\left(n^{3}\alpha^{\max}_n + n^5 \left(\alpha^{\max}_n\right)^3\right)M^{-\frac{1}{4d}},
\]
where we defined $\alpha^{\max}_{n} :=\max_{1\le i<j\le n} |\alpha_{n,i,j}|.$
\end{theorem}
\noindent Therefore, given $\alpha^{\max}_n=\operatorname{poly}(n),$ it is enough to estimate the free energy $F_\beta(H_{n,M})$ of the truncated Hamiltonian $H_{n,M}$ up to accuracy $\eps>0$ for $M=\operatorname{poly}(n,\varepsilon^{-1})$. To do so, we consider the free energy difference
\begin{align*}
\Delta F_\beta:=F_\beta(H_{n,M})-F_\beta(H_{0,n}).
\end{align*}
Given the path $H(s):=(1-s)H_{0,n}+sH_{n,M}$, $s\in[0,1]$, we have
\begin{align}
\Delta F_\beta=\int_0^1 \frac{d}{ds} F_\beta(H(s))\,ds= \int_0^1\,\Tr\Big(\sigma_\beta(H(s))\Pi_M W_n \Pi_M \Big)\, ds.\nonumber
\end{align}
Therefore, it suffices to estimate the thermal averages of $\Pi_M W_n\Pi_M$ with respect to the Gibbs state $\sigma_\beta(H_{n,M})$ over enough $s\in[0,1]$ in order to recover a good approximation of $\Delta F_\beta$. Since $F_\beta(H_{0,n})={dn}\beta^{-1}\log(2\sinh\beta)$ is known, this directly yields an additive approximation of $F_{\beta}(H_{n})$.

In addition to the free energies, we establish in Section~\ref{sec:EntropyTraceDistanceConv} that the Gibbs state themselves is also efficiently approximated in trace distance.
\begin{theorem}[Trace norm perturbation bound]
\label{thm:trace-norm-perturbation-boundIntro}
Fix \(d\in\{2,3\}\), \(\beta>0\). Then there exists \(C_\beta<\infty\), independent of \(n,M\), such that
\begin{equation}
\|\sigma_\beta(H_n)-\sigma_\beta(H_{n,M})\|_1 \le C_\beta \sqrt{n^3\alpha^{\max}_n + n^5(\alpha^{\max}_n)^3}\ M^{-\frac{1}{8d}}.
\end{equation}
\end{theorem}




\subsection{Quantum Gibbs sampling} 
To sample from the Gibbs state $\sigma_\beta(H_{n,M})$, we introduce and analyse a new quantum algorithm based on recent advances in quantum Gibbs sampling via quantum Markov semigroups \cite{RallWangWocjan2023,chen2023quantum,gilyen2024quantum,chen2023efficient,ding2025efficient,GalkowskiZworskiDavies}. In particular, in our previous work \cite{BeckerRouzeSalzmannToAppearcmp}, we extended this framework to infinite-dimensional quantum systems, which is the relevant setting here. As a separate application, we used it in \cite{BeckerRouzeSalzmannBose} to establish efficient quantum Gibbs sampling and free energy estimation for Bose-Hubbard models.

To introduce the generator of the semigroup, we make use of quantum Sobolev spaces \cite{gondolf2024energy}: given $\delta_1,\delta_2\ge 0$ and $\widetilde{H}_{n,M}:=H_{n,M}+h_0+1$,
\begin{align*}
D(\mathcal{W}_{H_{n,M}}^{\delta_1,\delta_2}):=\Big\{\widetilde{H}_{n,M}^{-\delta_1}a \widetilde{H}_{n,M}^{-\delta_2}\,\Big|\,a\in\mathscr{T}_1(\mathcal{H}_n)\Big\}\text{ with norms }
\|x\|_{\mathcal{W}_{H_{n,M}}^{\delta_1,\delta_2}}:=\|\mathcal{W}_{H_{n,M}}^{\delta_1,\delta_2}(x)\|_1,
\end{align*}
where $\mathcal{W}^{\delta_1,\delta_2}_{H_{n,M}}(x):=\widetilde{H}_{n,M}^{\delta_1} x \widetilde{H}_{n,M}^{\delta_2}$ and $\mathscr{T}_1(\mathcal{H}_n)$ denotes the Banach space of trace-class operators on $\mathcal{H}_n:=L^2((\mathbb{R}^d)^n)$. Given $h_{ij}:=-\Delta_{x_{ij}}+|x_{ij}|^2$, we consider the ladder operators on the form domain $Q(h_{ij})\equiv D((h_{ij}+1)^{1/2})$
\begin{align*}
a_{ij}:=\frac{x_{ij}+\partial_{x_{ij}}}{\sqrt{2}}\qquad \text{ and }\qquad a_{ij}^\dagger:=\frac{x_{ij}-\partial_{x_{ij}}}{\sqrt{2}}\qquad i\in[n],\,j\in[d].
\end{align*}

 We show in \cite[Proposition 4.3]{BeckerRouzeSalzmannToAppearcmp} that these operators induce filtered energy jumps
\begin{align}\label{jumpLalpha}
L^\alpha{\psi}:=\int_{-\infty}^\infty f(t) \,e^{itH_{n,M}}a^\sigma_{ij} e^{-itH_{n,M}}{\psi}\,dt\,,\quad
{\psi}\in Q(\widetilde{H}_{n,M}), ~~\alpha=(i,j,\sigma)\in [n]\times [d]\times \{\emptyset,\dagger\}.
\end{align}
for some Schwartz function $f\in\mathcal{S}(\mathbb{R},\mathbb{C})$, later referred to as the filter function of the sampler, with Fourier transform  satisfying the following detailed balance and boundedness conditions
\begin{align}
\label{eq:KMSFilterFunction}
\overline{\widehat{f}(\nu)}=\widehat{f}(-\nu)\,e^{-\beta\nu/2}\qquad \text{ and }\qquad \|\widehat{f}\|_{L^\infty},\quad  \|\nu\mapsto e^{\frac{\beta\nu}{2}}\widehat{f}(\nu)\|_{L^\infty} <\infty.
\end{align}
Next, given $\sigma_E\ge 0$, we introduce the operator  
\begin{align}\label{intrepGGG}
 G\psi:
&=-\sum_{\alpha\in\mathcal{A}}\int_{-\infty}^\infty g(t)\,e^{itH_{n,M}}(L^\alpha)^\dagger L^\alpha e^{-itH_{n,M}}{\psi}\,dt\,,\qquad {\psi}\in D(\widetilde{H}_{n,M}),
\end{align}
where $g(t)=\frac{1}{2\pi}\int_{\mathbb R}
\frac{e^{-\nu^2/8\sigma_E^2}}{1+e^{\beta\nu/2}}
e^{-i\nu t}\,d\nu$, and the sum is over tuples $\alpha=(i,j,\sigma)\in \mathcal{A}:=[n]\times[d]\times\{\emptyset,\dagger\}$. We also consider the map
\begin{align}\label{Phimapintrep}
\Phi(\rho):=\,
\sigma_E\,\sqrt{\frac{2}{\pi}}\sum_{\alpha\in\mathcal{A}}\int_{\mathbb{R}}e^{-2\sigma_E^2 s^2}\, \, X^{\alpha}_s \cdot \rho \cdot  (X_s^\alpha)^\dagger \,ds\,\qquad \rho\in D(\mathcal{W}_{H_{n,M}}^{\frac{1}{2},\frac{1}{2}}).
\end{align}
where $X^{\alpha}_s:=e^{isH_{n,M}}L^\alpha e^{-isH_{n,M}}$, and the dot products refer to the multiplication rule 
$X_s^\alpha\cdot \rho\cdot (X_s^\alpha)^\dagger \equiv X_s^\alpha\widetilde{H}_{n,M}^{-1}(X^\alpha_s\widetilde{H}_{n,M}^{-1}\mathcal{W}_{H_{n,M}}^{1,1}(\rho))^\dagger)^\dagger$. The following generation theorem is from \cite{BeckerRouzeSalzmannToAppearcmp}:



\begin{proposition}\cite[Theorem 2.1 and Corollary 4.4]{BeckerRouzeSalzmannToAppearcmp}
\label{condpropdefsampler}
 The following map defines the generator of a strongly continuous semigroup of completely positive, trace preserving maps over $\mathscr{T}_1(\cH_n)$, with a domain containing $D(\mathcal W_{H_{n,M}}^{\frac{1}{2},\frac{1}{2}}) \cap D(\mathcal W_{H_{n,M}}^{1,0}) \cap D(\mathcal W_{H_{n,M}}^{0,1})$:  
\begin{align}
\mathcal L_{\sigma_E,H_{n,M}}(\rho):=
\sum_{\alpha\in\mathcal{A}}\, G\cdot \rho+\rho\cdot G^\dagger +\Phi(\rho)\,.
\nonumber
\end{align}
Moreover, $\operatorname{ker}(\cL_{\sigma_E,H_{n,M}})=\mathbb{C}\, \sigma_\beta(H_{n,M})$.
\end{proposition}


\subsection{Mixing time via spectral gap analysis}

\noindent Next, as proved in \cite{BeckerRouzeSalzmannToAppearcmp}, the generator $\mathcal L_{\sigma_E,H_{n,M}}$ gives rise to the self-adjoint generator $L_{\sigma_E,H_{n,M}}$ of a strongly continuous, symmetric semigroup of contractions on the Hilbert space $\mathscr{T}_2(\cH_n)$ of Hilbert-Schmidt operators on $\cH_n$ with the associated norm $\|.\|_2$: introducing the embedding
\[
\iota_2:\mathscr T_2(\mathcal H_n)\to \mathscr T_1(\mathcal H_n),
\qquad
\iota_2(x)=\sigma_\beta(H_{n,M})^{1/4}x\sigma_\beta(H_{n,M})^{1/4},
\]
the corresponding semigroups are intertwined according to
\begin{align}
e^{t\mathcal L_{\sigma_E,H_{n,M}}}\circ \iota_2(x)
=
\iota_2\circ e^{tL_{\sigma_E,H_{n,M}}}(x).
\end{align}
 Here, we utilize the spectral properties of the operator $L_{\sigma_E,H_{n,M}}$ in order to control the mixing time: given a well-chosen initial state $\rho_{\operatorname{ini}}$ of the system with $\rho_{\operatorname{ini}}\le \mathfrak{c}\,\sigma_\beta(H_{n,M})$ and $\varepsilon\in (0,1)$,
\begin{align}
t_{\operatorname{mix}}(\varepsilon)&:=\inf\big\{t\ge 0:\,\big\|e^{t\mathcal{L}_{\sigma_E,H_{n,M}}}(\,\rho_{\operatorname{ini}}\,)-\sigma_\beta(H_{n,M})\big\|_1\le \varepsilon\big\}\le  \frac{2\log\big({\mathfrak{c}}/{\varepsilon}\big)}{\operatorname{gap}(L_{\sigma_E,H_{n,M}})}~~,\label{mixingtimeintro}
\end{align}
where $\operatorname{gap}(L_{\sigma_E,H_{n,M}})$ denotes the spectral gap of the generator $L_{\sigma_E,H_{n,M}}$. 
In particular, in \cite[Proposition 4.2]{BeckerRouzeSalzmannToAppearcmp} (see also \cite{slezak2026polynomial}), it was shown that the larger $\sigma_E$, the smaller the gap. Thus, it appears to be enough to control the gap of the generator $L_{H_{n,M}}\equiv L_{\infty,H_{n,M}}$. The latter takes the following form on $\mathscr F:=\operatorname{span}\{|E\rangle\langle E'|:E,E'\in\operatorname{Sp}(H_{n,M})\}$: denoting $\Gamma_\tau(X):=e^{\tau H}Xe^{-\tau H}$ and $A^\alpha=a_{ij}^{\sigma}$ for any $\alpha=(i,j,\sigma)\in\mathcal{A}$,
\begin{align}\label{LHnmgenerator}
L_{H_{n,M}}(X)
=
-i(B_+X-XB_-)
+\sum_{\alpha\in\mathcal A}
\Big(L^\alpha_+\,X\,(L^\alpha_-)^\dagger-\tfrac12 K^\alpha_+X-\tfrac12 XK^\alpha_-\Big),
\qquad X\in\mathscr F.
\end{align}
with 
\[
B:=\frac{i}{2}\sum_{\alpha\in\mathcal A}\sum_{E,F,G\in\operatorname{Sp}(H)}
\tanh\!\Big(\frac{\beta(F-E)}{4}\Big)\,
\overline{\widehat f(G-F)}\,\widehat f(G-E)\,
P_{F}(A^\alpha)^\dagger P_GA^\alpha P_E
\]
and
\[
L^\alpha_\pm:=\Gamma_{\pm\beta/4}(L^\alpha),
\qquad
K^\alpha_\pm:=\Gamma_{\pm\beta/4}\big((L^\alpha)^\dagger L^\alpha\big),\qquad B_\pm:=\Gamma_{\pm\beta/4}(B).
\] 
We refer to \cite{BeckerRouzeSalzmannToAppearcmp} for a full treatment of these generators in the case of unbounded Hamiltonians on separable Hilbert spaces. Our next main result, proven in Section~\ref{sec:GapForFiniteRankPerturb}, establishes that the spectral gap of $L_{H_{n,M}}$ is always positive:
\begin{theorem}\label{mainthmpositivegap}
For any $\sigma_E\in (0,\infty]$, $\operatorname{gap}(L_{\sigma_E,H_{n,M}})>0$.
\end{theorem}

A weak coupling scaling of the Coulomb energies that allows for the presence of a spectral gap that is uniform in the particle number is presented in Theorem \ref{thm:uniform_gap}.





\subsection{Circuit implementation}


\noindent In order to obtain an efficient qubit-based circuit implementation of the Gibbs sampler, we consider a finite-dimensional approximation of the generator $\cL_{\sigma_E,H_{n,M}}$, in which both the bare jump operators and the Hamiltonian are replaced by finite-dimensional truncations. As established in the abstract framework of \cite{BeckerRouzeSalzmannToAppearcmp}, such truncated generators and their dynamics approximate their untruncated counterparts well on suitably energy-constrained inputs, such as $\rho_{\operatorname{ini}}$, provided suitable assumptions on the Hamiltonian and the bare jump operators are satisfied. In Section~\ref{sec:GibbscircuitCoulomb}, we verify that these assumptions hold for the present model.

In the finite-dimensional setting, the integral representations \eqref{jumpLalpha}, \eqref{intrepGGG}, and \eqref{Phimapintrep} are well suited for implementation via block-encodings of the Heisenberg evolution and the truncated bare jump operators, combined with an appropriate time discretization, as established in \cite{chen2023efficient,ding2025efficient}. Combining these results, we show that the Gibbs states of $H_{n,M}$ and $H_n$ can be prepared efficiently on a qubit based quantum computer:

\begin{theorem}[Gibbs state preparation via quantum circuits]
\label{thm:GibbsPrepareCoulombIntro}
Let $d\in\{2,3\}$ and \\ $\eps,\beta,\sigma_E>0$. Consider  
\begin{align*}
H_{n,M}\equiv H_{0,n} + W_{n,M} \qquad\text{with}\qquad M = \operatorname{poly}(n,q(1/\eps))\,, \ \, \alpha^{\max}_{n} \lesssim 1,
\end{align*} 
 where $q:\R_+\to \R_+$ denotes some non-decreasing function which is lower bounded as $q(x)\gtrsim \log(x)$.
Denote $\lambda_2 := \operatorname{gap}(L_{\sigma_E,H_{n,M}})>0,$ where positivity of the gap follows by Theorem~\ref{mainthmpositivegap}.

Then we have that Gibbs state $\sigma_\beta(H_{n,M})$ can be prepared within $\eps$-trace distance on a quantum computer with $\mathcal{O}\left(n\log (n\,q(1/\eps))\,\log\log(1/\lambda_2)))\right)$ many qubits with circuit depth\footnote{Here, the $\mathcal{O},$ $\widetilde{\mathcal{O}}$ and $\Omega$ notations hide constants independent of the displayed parameters and $\widetilde{\mathcal{O}}$ additionally suppresses subdominant $\operatorname{poly}\log$ factors in $1/\lambda_2$.}
\begin{align*}
  \widetilde{ \mathcal{O}}\left(\frac{1}{\lambda_2}\operatorname{poly}\left(n, q(1/\eps)\right)\right).
\end{align*}
Hence, using Theorem~\ref{thm:trace-norm-perturbation-boundIntro} and $M = \Theta\left(\left(\frac{n^{5/2}}{\eps}\right)^{8d}\right),$ this provides a preparation procedure for $\sigma_\beta(H_n)$ with the same complexities as above corresponding to the specific choice $q(x) = x^{8d}.$ 
\end{theorem}

In our next main result, proven in Section~\ref{sec:eff_impl}, we combine Theorems~\ref{thmfreenregyblablaIntro} and~\ref{thm:GibbsPrepareCoulombIntro} to provide an efficient quantum algorithm for estimating the free energy of $H_n:$
\begin{theorem}[Quantum algorithm for free energy estimation]
The free energy $F_\beta(H_n)$ can be estimated with accuracy $\eps>0$ and probability of failure bounded by $\delta>0$ on a quantum computer with $\mathcal{O}\left(n\log( n/\eps)\,\log\log(1/\lambda^{\min}_{2})\right)$ many qubits with total runtime of order 
\begin{align*}
    \widetilde{ \mathcal{O}}\left(\frac{1}{\lambda^{\min}_{2}}\log\left(1/\delta\right)\operatorname{poly}\left(n\,,\,1/\eps\right)\right)
\end{align*}
where $\lambda^{\min}_{2}:= \min_{s\in[0,1]} \lambda_2(s)>0$ and $\lambda_2(s)\equiv\operatorname{gap}(L_{\sigma_E,H(s)})>0$ with   $H(s):=(1-s)H_{0,n}+sH_{n,M}$ for $\alpha^{\max}_n\lesssim 1$ and $M=\Theta\left(\left(\frac{n^{5}}{\eps}\right)^{4d}\right).$ 
\end{theorem}







\subsection{Related work}


\paragraph{\textbf{Classical systems}}
Given potential $V_n$ which associates to each $n$-particle configuration $(x,p)\in\mathbb{R}^{3n\times 3n}$ an energy $V_n(x,p)$, the goal of classical molecular dynamics is to compute averages and free energies with respect to the Boltzmann-Gibbs measure $d\mu_\beta(x,p):=e^{-\beta V_n(x,p)}/\mathcal{Z}_\beta(V) dx$. Markov chain Monte Carlo methods have been employed to estimate $F_\beta(V_n)$ \cite{rousset2010free,chafai2019simulating,chafai2021coulomb}. A typical example is the over-damped Langevin dynamics, whose generator on $L^2(\mu_\beta(V_n))$ takes the form $\beta^{-1}\Delta -\nabla V_n\cdot \nabla$. Our generators $\cL_{\sigma_E,H_{n,M}}$ can be understood as quantum extensions of such methods for the preparation of $\sigma_\beta(H_n)$, where one promotes the momentum $p$ to the operator $-i\nabla$. In fact, both generators are known to coincide with that of the Ornstein Uhlenbeck semigroup in the case of $V_n=H_n+\Delta=|x|^2$ when restricted e.g.~to the commutative algebra $\mathcal{S}(\mathbb{R}^n)$ \cite{OU}. For both the classical Langevin dynamics and our quantum dynamics, the rate of convergence is controlled by the spectral gap of the corresponding generator. When considering $V_n(x)=W_n(x)$ the potential associated to Coulomb interactions, the positivity of the spectral gap was proved in \cite{Bolley2018} in the two-dimensional setting, using local Poincaré inequalities and comparison to the uniform distribution on compact domains (see also \cite{chafai2021aspects,Bourgade2019,Akemann2019,Chafa2020,lu2020geometric,Duong2024} for more recent results on such dynamics). However, to the best of our knowledge, the dependency of the gap on the number of particles for a given inverse temperature $\beta$ remains open. This contrasts with non-singular interaction settings which have extensively been studied in relation with McKean--Vlasov equations \cite{Malrieu2001,Mlard1996,Sznitman1991}.


\smallskip

\paragraph{\textbf{Hybrid classical-quantum descriptions}}
Recent works have been concerned with the computation of free energy differences for the description of chemical systems within the Born-Oppenheimer approximation where the wavefunctions of the nuclei 
are considered independent of the wavefunctions of the electrons \cite{xu2023slow,ries2022relative,santagati2024drug}. Molecular dynamics simulations introduce a further approximation by treating the nuclei as classical particles while retaining the quantum description of the electrons. Quantum simulation algorithms within these approximations and beyond were recently explored in \cite{kassal2011simulating,ollitrault2020nonadiabatic,fedorov2021ab,sokolov2021microcanonical,o2019calculating,o2022efficient,steudtner2023fault}. For many practical applications, we need to compute the
properties of the chemical systems in the canonical ensemble \cite{xu2023slow,santagati2024drug,tuckerman2000understanding,manathunga2022computer}, which involves preparing the classical  nuclei according to the Boltzmann distribution at a certain $\beta$. 

Inspired by  \cite{joseph2020koopman,jin2023time},
in \cite{simon2024improved,huang2025fullqubit,GuntherEtAl2025}, the authors proposed a method to simulate the classical evolution of the nuclei on a quantum computer based on Koopman and von Neumann 
Liouvillian formalism \cite{koopman1931hamiltonian}, whose generator depends explicitly on the ground state energy of the electrons treated quantumly using quantized simulation methods  \cite{babbush2019quantum,su2021fault}.
While these quantum algorithms scale favourably with the inverse precision and number of particles, reaching e.g.~a $\mathcal{O}(n^7/\varepsilon)$-runtime, they also rely on the so-called equilibration time. While the latter can scale as $\mathcal{O}(1/\varepsilon)$ with additional weighting over time averages for some completely integrable systems \cite{cances2005long}, this parameter is often unknown and may potentially be unbounded.

In contrast to these prior works, our algorithm is based on the quantum Gibbs sampling of a controllable truncation of the Hamiltonian, which can be regarded as a direct quantum analogue of Monte Carlo methods for classical free energies. In particular, our method does not rely on the Born-Oppenheimer approximation, and replaces heuristic equilibration assumptions with a guaranteed convergence to the target Gibbs state for any fixed number of particles, due to the positivity of the gap proved in Theorem \ref{mainthmpositivegap}. It remains nonetheless an important open question for which parameters the spectral gap can be shown to scale inverse polynomially with the system size, hence yielding an efficient end-to-end method for the computation of free energy differences for interacting quantum particles. 









































































% \section{Finite rank perturbation of partition functions}\label{perturbationpartitionfunction}


% \noindent In this section we construct a general perturbative framework for the truncation of the interaction to low-energy one-body modes. 
% We start with the following abstract result: we recall that, given a Hamiltonian $H$ satisfying the Gibbs Hypothesis, we denote $\mathcal{Z}_\beta(H):=\Tr(e^{-\beta H})$,  $F_\beta(H):=-\beta^{-1}\log\mathcal{Z}_\beta(H)$ and $\sigma_\beta(H)=e^{-\beta H}/\mathcal{Z}_\beta(H)$.
% We also recall that the form domain of a self-adjoint operator $H\ge -h_0$ corresponds to the domain $D((H+h_0+1)^{1/2})$.


% \begin{lemma}[Perturbation bound for free energies under relative form convergence]
% \label{lem:free-energy-relative-form-convergence}
%  Let \(H_{0}\) be a self-adjoint operator on \(\mathcal H\), bounded from below, $H_0\ge -h_0$, with compact resolvent. Let \(W\) and \(W'\) be symmetric quadratic forms on the form domain \(Q(H_{0})\), and define the form sums
% \[
% H:=H_{0}+ W,
% \qquad
% H':=H_{0}+ W'.
% \]
% Assume further that the forms \(W\) and \(W'\) are \(H_{0}\)-form bounded with relative bound strictly smaller than \(1\). Let
% \begin{equation}
% \label{eq:DeltaNM}
% \Delta:=H-H'=W-W'
% \end{equation}
% in the sense of quadratic forms, and assume that there exist numbers \(\varepsilon\ge0\) and \(c\ge0\), such that
% \begin{equation}
% \label{eq:relative-form-smallness-delta}
% |\Delta[\psi]|
% \le
% \varepsilon \,\langle \psi,H_{0}\psi\rangle + c\|\psi\|^2
% \qquad\text{ for all }\qquad \psi\in Q(H_{0}).
% \end{equation}
% Then, the Gibbs states associated to $H$, resp.~$H'$, as $\sigma_\beta(H)$, resp.~$\sigma_\beta(H')$ are well-defined as trace-class operators, and
% \begin{equation}
% \label{eq:sharper-max-estimate}
% |F_\beta(H)-F_\beta(H')|
% \le
% {\varepsilon}
% \max\bigl\{\Tr(\sigma_\beta(H) H_{0}),\Tr(\sigma_\beta(H')H_{0})\bigr\}
% +{c}.
% \end{equation}
% \end{lemma}

% \begin{proof}

% Since the forms \(W\) and \(W'\) are \(H_{0}\)-form bounded with a relative bound strictly smaller than \(1\), \(H\) and \(H'\) are self-adjoint and bounded from below on \(Q(H_{0})\) with a compact resolvent. Hence, the Gibbs operators \(e^{-\beta H}\) and \(e^{-\beta H'}\) are trace class. Next, by the Bogoliubov inequality \cite{Dereziski2003}
% \begin{equation}
% \label{eq:PB}
% \log \frac{\Tr(e^{X+Y})}{\Tr(e^{X})}
% \ge
% \frac{\Tr(e^{X}Y)}{\Tr(e^{X})}
% \end{equation}
% for self-adjoint \(X,Y\) with $X:=-\beta H$, $Y:=-\beta(H'-H)$ as well as for $X:=-\beta H'$, $Y:=-\beta(H-H')$, we get 
% \[
% \Tr(\sigma_\beta(H)\Delta)
% \le \frac1\beta\log\frac{\mathcal{Z}_\beta(H')}{\mathcal{Z}_\beta(H)}
% \le \Tr(\sigma_\beta(H')\Delta).
% \]
% Hence
% \begin{equation}
% \label{eq:free-energy-max-step}
% |F_\beta(H)-F_{\beta}(H')|
% \le
% \max\Bigl\{
% |\Tr(\sigma_\beta(H)\Delta)|,
% |\Tr(\sigma_\beta(H')\Delta)|
% \Bigr\}.
% \end{equation}
% Next, applying \eqref{eq:relative-form-smallness-delta} in the spectral decompositions of \(\sigma_\beta(H)\) and \(\sigma_\beta(H')\) gives
% \[
% |\Tr(\sigma_\beta(H)\Delta)|
% \le \varepsilon\,\Tr(\sigma_\beta(H)H_{0})+c\qquad \text{ and }\qquad 
% |\Tr(\sigma_\beta(H')\Delta)|
% \le \varepsilon\,\Tr(\sigma_\beta(H')H_{0})+c.
% \]
% Inserting these estimates into \eqref{eq:free-energy-max-step}, we obtain
% \[
% |F_\beta(H)-F_\beta(H')|
% \le
% {\varepsilon}
% \max\bigl\{\Tr(\sigma_\beta(H)H_{0}),\Tr(\sigma_\beta(H')H_{0})\bigr\}
% +c.
% \]
% \end{proof}


% \noindent For now on, we consider the $n$-particle Hamiltonian
% \[
% H_n:=H_{0,n}+W_n
% \]
% on \(L^2((\R^d)^n)\), where $h:=-\Delta+|x|^2$ over $L^2(\mathbb{R}^d)$, $H_{0,n}:=\sum_{i\in[n]}h_i$ and for $\psi\in Q(H_{0,n})$, denote $W_n[\psi]=\alpha_n\sum_{1\le i<j\le n}W_{ij}[\psi]$, where
% \[
% W_{ij}[\psi]:=\int_{\R^{nd}}w_d(x_i-x_j)|\psi(x)|^2 \ dx,\qquad
% w_d(y):=
% \begin{cases}
% -\log|y|, & d=2,\\[0.3em]
% |y|^{-1}, & d=3.
% \end{cases}
% \]
% Let further $P_M$ be the spectral projector onto the first $M$ eigenvalues \((\lambda_k)_{k=1}^M\) of \(h\). Define \(\Pi_M:=P_M^{\otimes n}\), \(\Pi_M^\perp:=1-\Pi_M\), and consider the form
% \[W_{n,M}[\psi]:=W_n[\Pi_M\psi]\]
% and the associated Hamiltonian $H_{n,M}$. In the next result, we establish relative form boundedness of $\Delta_{n,M}:=W_n-W_{n,M}$ with respect to $H_{0,n}$:


% \begin{lemma}[Product cutoff estimate without scaling assumption]
% \label{lem:product-cutoff-riesz-no-scaling}
% Given $d\in\{1,2\}$, for all \(n,M\) and \(\psi\in Q(H_{0,n})\),
% \begin{align}
% |W_n[\psi]|&\lesssim \,|\alpha_n|n(n-1)\langle\psi,(H_{0,n}+1)\psi\rangle,\label{eq111}\\
% |W_n[\Pi_M^\perp\psi]|&\lesssim\,|\alpha_n|n(n-1)(\lambda_{M+1}+1)^{-1/2}\langle\psi,(H_{0,n}+1)\psi\rangle,\label{eq222}\\
% |\Delta_{n,M}[\psi]|&\lesssim\,|\alpha_n|n(n-1)(\lambda_{M+1}+1)^{-1/4}\langle\psi,(H_{0,n}+1)\psi\rangle.\label{eq333}
% \end{align}
% \end{lemma}

% \begin{proof}
%  Fix $1\le i<j\le n$ and $d=3$. Introduce the relative variable $r:=2^{-1/2}
% {(x_i-x_j)}$, so that ${|x_i-x_j|^{-1}}=2^{-1/2}{|r|^{-1}}$. By Hardy's inequality in $\mathbb R^3$,
% \[
% \left\|\frac{u}{|r|}\right\|^2\le 4\|\nabla_r u\|^2.
% \]
% Hence, by Cauchy--Schwarz and Young's inequalities,
% \[
% \left\langle u,\frac{1}{|x_i-x_j|}u\right\rangle
% =
% 2^{-1/2}\left\langle u,\frac{1}{|r|}u\right\rangle
% \le
% \sqrt2\,\|u\|\,\|\nabla_r u\|
% \le
% \frac12\|\nabla_r u\|^2+\|u\|^2.
% \]
% Since
% \[
% \|\nabla_r u\|^2\le \langle u,(-\Delta_{x_i}-\Delta_{x_j})u\rangle
% \le
% \langle u,\bigl(-\Delta_{x_i}-\Delta_{x_j}+|x_i|^2+|x_j|^2\bigr)u\rangle,
% \]
% we obtain
% \[
% \frac{1}{|x_i-x_j|}
% \le
% \frac12\bigl(-\Delta_{x_i}-\Delta_{x_j}+|x_i|^2+|x_j|^2\bigr)+1
% \]
% as quadratic forms. Summing over all pairs $(i,j)$ yields
% \[
% W_n
% \le
% \frac{|\alpha_n|}2\sum_{1\le i<j\le n}\bigl(-\Delta_{x_i}-\Delta_{x_j}+|x_i|^2+|x_j|^2\bigr)+|\alpha_n|\,\binom{n}{2}=\frac{|\alpha_n|}{2}H_{0,n}+|\alpha_n|\,\binom{n}{2}.
% \]
% In dimension \(d=2\), let \(u\in Q(h_1+h_2)\). Fix \(0<\varepsilon<1\). Since
% \[
% |\log r|\le C_\varepsilon\bigl(1+r^{-\varepsilon}+r^\varepsilon\bigr),\qquad r>0,
% \]
% it suffices to estimate these three terms separately. First,
% \[
% \int_{\R^2\times\R^2}|u(x,y)|^2\,dx\,dy=\|u\|^2
% \le
% \langle u,(h_1+h_2+1)u\rangle.
% \]
% Next, using the change of variables \(r=x-y\), \(s=x+y\), we obtain
% \[
% \int_{\R^2\times\R^2}|x-y|^{-\varepsilon}|u(x,y)|^2\,dx\,dy
% =
% \frac{1}{4}\int_{\R^2\times\R^2}|r|^{-\varepsilon}|u(r,s)|^2\,dr\,ds.
% \]
% Set
% \[
% F(r):=\|u(r,\cdot)\|_{L^2(\R^2_s)}.
% \]
% Then
% \[
% \int_{\R^2\times\R^2}|r|^{-\varepsilon}|u(r,s)|^2\,dr\,ds
% =\int_{\R^2}|r|^{-\varepsilon}F(r)^2\,dr.
% \]
% Choose \(p>1\) such that \(\varepsilon p<2\), and let \(q=\frac p{p-1}\) its Hölder conjugate. Since
% \(|r|^{-\varepsilon}\in L^p(B_1(\R^2))\) for $B_1(\mathbb{R}^2)$ the unit ball in $\mathbb{R}^2$, Hölder's inequality in the \(r\)-variable yields
% \[
% \int_{|r|\le1}|r|^{-\varepsilon}F(r)^2\,dr
% \le
% \bigl\||r|^{-\varepsilon}\bigr\|_{L^p(B_1(\mathbb{R}^2))}\|F^2\|_{L^q(\R^2)}
% =
% C\|F\|_{L^{2q}(\R^2)}^2.
% \]
% Now \(2q<\infty\), so by the Sobolev embedding
% \[
% H^1(\R^2_r;L^2(\R^2_s))\hookrightarrow L^{2q}(\R^2_r;L^2(\R^2_s))
% \]
% we obtain
% \[
% \|F\|_{L^{2q}(\R^2)}=\|u\|_{L^{2q}(\mathbb{R}_r^2,L^2(\mathbb{R}_s^2))}
% \lesssim\|u\|_{H^1(\R^2_r;L^2(\R^2_s))}.
% \]
% Hence
% \[
% \int_{|r|\le1}\int_{\R^2}|r|^{-\varepsilon}|u(r,s)|^2\,ds\,dr
% \lesssim
% \|u\|_{H^1(\R^2_r;L^2(\R^2_s))}^2.
% \]
% Since \(|r|^{-\varepsilon}\le1\) on \(\R^2\setminus B_1(\mathbb{R}^2)\), we also control
% \[
% \int_{|r|>1}\int_{\R^2}|r|^{-\varepsilon}|u(r,s)|^2\,ds\,dr
% \le
% \|u\|_{L^2(\R^4)}^2.
% \]
% Therefore
% \[
% \int_{\R^2\times\R^2}|r|^{-\varepsilon}|u(r,s)|^2\,dr\,ds
% \lesssim
% \Bigl(\|u\|_{H^1(\R^2_r;L^2(\R^2_s))}^2+\|u\|_{L^2(\R^4)}^2\Bigr).
% \]
% By the Sobolev embedding \(H^1(\R^4)\hookrightarrow L^4(\R^4)\),
% \[
% \|u\|_{L^4(\R^4)}^2
% \lesssim\|u\|_{H^1(\R^4)}^2
% \lesssim
% \langle u,(h_1+h_2+1)u\rangle.
% \]
% Finally, since \(0<\varepsilon<2\), there exist a constants $C_\varepsilon,C_\varepsilon'>0$ such that
% \[
% |x-y|^\varepsilon\le C_{\varepsilon}\bigl(1+|x|^\varepsilon+|y|^\varepsilon\bigr)
% \le C'_{\varepsilon}\bigl(1+|x|^2+|y|^2\bigr),
% \]
% and therefore
% \[
% \int_{\R^2\times\R^2}|x-y|^\varepsilon |u(x,y)|^2\,dx\,dy
% \le
% C_\varepsilon'\langle u,(h_1+h_2+1)u\rangle.
% \]
% Putting the three estimates together yields
% \[
% \int_{\R^2\times\R^2}|\log|x-y||\,|u(x,y)|^2\,dx\,dy
% \lesssim
% \langle u,(h_1+h_2+1)u\rangle
% \]
%  Summing over all combinations of $x=x_i,y=x_j$ yields then 
% \[
% |W_n[\psi]|\lesssim |\alpha_n|n(n-1)\langle\psi,(H_{0,n}+1)\psi\rangle.
% \]
% This ends the proof of \eqref{eq111}. For \eqref{eq222}, we may proceed for both $d=2,3$ as follows:
% since \((H_{0,n}+1)\Pi_M^\perp\ge(\lambda_{M+1}+1)\Pi_M^\perp\), we have \[ \|\Pi_M^\perp\psi\|^2\le(\lambda_{M+1}+1)^{-1}\langle\psi,(H_{0,n}+1)\psi\rangle,\] hence by \eqref{eq111}, the commutation $[\Pi_M^\perp,H_{0,n}]=0$ and the Cauchy Schwarz inequality,
% \begin{align}
% |W_n[\Pi_M^\perp\psi]|&\lesssim |\alpha_n|n(n-1)\langle \Pi_M^\perp\psi,(H_{0,n}+1)\psi\rangle\nonumber\\
% &\le (\lambda_{M+1}+1)^{-1/2}|\alpha_n|n(n-1)\langle \psi,(H_{0,n}+1)\psi\rangle,\label{eq2222}
% \end{align}
% which ends the proof of \eqref{eq222}. For \eqref{eq333}, we associate to \(W_n\) the sesquilinear form
% \[
% W_n(\varphi,\chi)
% :=
% \alpha_n\sum_{1\le i<j\le n}
% \int_{(\R^d)^n}w_d(x_i-x_j)\,\varphi(x)\,\overline{\chi(x)}\,dx,
% \]
% so that \(W_n[\psi]=W_n(\psi,\psi)\). Hence, writing \(P:=\Pi_M\) and \(Q:=\Pi_M^\perp\),
% \[
% W_n[\psi]-W_{n,M}[\psi]
% =
% W_n[P\psi+Q\psi]-W_n[P\psi]
% =
% W_n[Q\psi]+2\Re W_n(P\psi,Q\psi).
% \]
% Similarly, if \(\widetilde W_n\) denotes the positive sesquilinear form with kernel
% \(|w_d|\), then combining the Cauchy-Schwarz inequality, \eqref{eq111}, the commutation of $H_{0,n}$ with $P$ and $Q$, and another use of Cauchy-Schwarz as done previously in \eqref{eq2222}, we get
% \begin{align}
% |W_n(P\psi,Q\psi)|&\le \widetilde{W}_n[P\psi]^{1/2}\widetilde{W}_n[Q\psi]^{1/2}\nonumber\\
% &\lesssim |\alpha_n|n(n-1)\langle Q\psi,(H_{0,n}+1)Q\psi\rangle^{\frac{1}{2}}\langle P\psi,(H_{0,n}+1)P\psi\rangle^{\frac{1}{2}}\nonumber\\
% &\lesssim|\alpha_n|n(n-1)\langle Q\psi,(H_{0,n}+1)Q\psi\rangle^{\frac{1}{2}}(\lambda_M+1)^{-\frac{1}{4}}\langle \psi,(H_{0,n}+1)\psi\rangle^{\frac{1}{2}}\nonumber\\
% &\le |\alpha_n|n(n-1)(\lambda_M+1)^{-\frac{1}{4}}\langle \psi,(H_{0,n}+1)\psi\rangle\nonumber.
% \end{align}
% Finally, combining the last bound with the estimate for \(W_n[Q\psi]\) derived in \eqref{eq222} yields the claimed bound \eqref{eq333}.
% \end{proof}


% \noindent \noindent Combining Lemmas \ref{lem:free-energy-relative-form-convergence} and \ref{lem:product-cutoff-riesz-no-scaling}, we get that the difference of free energies is controlled by
% \begin{align}\label{almostthereyes}
% \frac{|F_\beta(H_n)-F_\beta(H_{n,M})|}{|\alpha_n|n^2}\lesssim (\lambda_{M+1}+1)^{-\frac{1}{4}} \max \left\{\Tr(\sigma_\beta(H_n)H_{0,n}),\Tr(\sigma_\beta(H_{n,M})H_{0,n})\right\}.
% \end{align}
% It remains to control each of the average energies appearing in the previous bound. We first bound them, via Lemma \ref{lem:thermodynamic-inequality}, by the difference between the free energies of the perturbed and unperturbed models. We then show in Lemmas \ref{lem:lower-bound-partition-gibbs} and \ref{lowerboundsnow} how to control this free-energy difference. The final energy bounds are stated in Lemma \ref{lem:uniform-H0-bound-abstract-detailed}.

% \begin{lemma}
% \label{lem:thermodynamic-inequality}
% Let \(K\) be a self-adjoint operator on a Hilbert space, bounded from below with smooth partition function. Then, for every \(t>0\),
% \begin{equation}
% \label{eq:UK-logZ}
% \Tr(\sigma_{t}(K)K)\le 2F_t(K)-F_{t/2}(K).
% \end{equation}
% \end{lemma}

% \begin{proof}
% Writing $U_K(t)=\Tr(\sigma_t(K)K)$,
% we immediately get
% \[
% -\partial_t\log \mathcal Z_t(K)
% =
% -\frac{\partial_t\mathcal Z_t(K)}{\mathcal Z_t(K)}
% =
% \frac{\Tr(Ke^{-tK})}{\Tr(e^{-tK})}
% =
% U_K(t).
% \]
% Moreover, by the quotient rule,
% \begin{align}
% U_K'(t)
% &=
% \frac{\Tr(Ke^{-tK})^2-\Tr(K^2e^{-tK})\Tr(e^{-tK})}{\Tr(e^{-tK})^2}\nonumber \\
% &=
% -\Tr(\sigma_{t}(K)K^2)+\bigl(\Tr(\sigma_{t}(K)K)\bigr)^2.
% \end{align}
%  Expanding \((K-U_K(t))^2\) yields $\Tr\bigl(\sigma_{t}(K)(K-U_K(t))^2\bigr)
% =
% \Tr(\sigma_{t}(K)K^2)-U_K(t)^2$. Thus
% \[
% U_K'(t)
% =
% -\Tr\bigl(\sigma_{t}(K)(K-U_K(t))^2\bigr)\le 0,
% \]
% This proves that \(t\mapsto U_K(t)\) is decreasing. Finally, fix \(t>0\). Since \(U_K\) is decreasing on \((0,\infty)\), we have
% \[
% U_K(t)\le \frac{2}{t}\int_{t/2}^t U_K(s)\,ds.
% \]
% Using \(U_K(s)=-\partial_s\log \mathcal Z_s(K)\), we obtain
% \[
% \frac{2}{t}\int_{t/2}^t U_K(s)\,ds
% =
% -\frac{2}{t}\int_{t/2}^t \partial_s\log \mathcal Z_s(K)\,ds
% =
% \frac{2}{t}\Bigl(\log \mathcal Z_{t/2}(K)-\log \mathcal Z_t(K)\Bigr),
% \]
% which is exactly \eqref{eq:UK-logZ}.
% \end{proof}
% \noindent Next, we control the free energy difference obtained in the previous Lemma. In what follows, we
% denote the one-body $\rho_t(x):=\langle x|\sigma_t(h)|x\rangle$. Since \(h=-\Delta+|x|^2\), the kernel $\rho_t(x):=\langle x|\sigma_t(h)|x\rangle$ of \(e^{-th}\) is given by the Mehler formula. In particular, \(\rho_t(x)\) is smooth, strictly positive, and has Gaussian decay as \(|x|\to\infty\):
% \[
% \rho_t(x)\le C_t e^{-c_t|x|^2}
% \]
% for suitable constants \(c_t,C_t>0\). Therefore, for any $d\in\{2,3\}$,
% \[
% \widetilde I_{d,t}
% :=
% \int_{\mathbb R^{2d}}|w_d(x-y)|\,\rho_t(x)\rho_t(y)\,dx\,dy<\infty.
% \]


% \begin{lemma}[Upper bounds on the free energy via free Gibbs states]
% \label{lem:lower-bound-partition-gibbs}Given $d\in\{2,3\}$, fix $t>0$. Then
% \begin{equation}
% \label{eq:lower-ZN-lemma}
% F_t(H_n)\le nF_t(h)+{\alpha_n}\binom{n}{2}\widetilde I_{d,t},
% \qquad \text{ and }\qquad 
% F_t(H_{n,M})\le nF_t(h)+{\alpha_n }\binom{n}{2}\widetilde I_{d,t}.
% \end{equation}
% \end{lemma}



% \begin{proof}

% By a use of the Bogoliubov inequality \eqref{eq:PB}, with $X=-tH_{0,n}$ and $Y=-t(H_n-H_{0,n})$,
% \[
% F_t(H_{n})
% \le
% {n}F_t(h)+\Tr(\sigma_t(H_{0,n}) W_n).
% \]
% Next, we estimate the interaction term: for each pair \(i<j\) one has
% \[
% \Tr\bigl(\sigma_t(H_{0,n})\, w_d(x_i-x_j)\bigr)
% =
% \int_{\R^{2d}} w_d(x-y)\rho_t(x)\rho_t(y)\,dx\,dy,
% \]
% independently of $i<j$. Hence
% \begin{align}
% \Tr(\sigma_t(H_{0,n}) W_n)
% &=
% \alpha_n\binom{n}{2}
% \int_{\R^{2d}} w_d(x-y)\rho_t(x)\rho_t(y)\,dx\,dy\nonumber\\
% &\le
% \alpha_n\binom{n}{2}
% \int_{\R^{2d}} |w_d(x-y)|\,\rho_t(x)\rho_t(y)\,dx\,dy
% \nonumber\\
% &=
% \alpha_n\binom{n}{2}\widetilde I_{d,t}.\nonumber
% \end{align}
% The proof for \(H_{n,M}\) is almost identical: by \eqref{eq:PB},
% \[
% F_t(H_{n,M})\ge nF_t(h)+\Tr(\sigma_t(H_{0,n}) W_{n,M}).
% \]
% Then, since \(\Pi_M\) commutes with \(H_{0,n}\), \(\Pi_M\) also commutes with \(\sigma_t(H_{0,n})\) and
% \[
% \Tr(\sigma_t(H_{0,n}) W_{n,M})
% =
% \alpha_n\sum_{1\le i<j\le n}\Tr\bigl(\Pi_M\sigma_t(H_{0,n})\Pi_M\, w_d(x_i-x_j)\bigr).
% \]
% Since \(\Pi_M\sigma_t(H_{0,n})\Pi_M\ge0\) is trace class, we may use $\bigl|\Tr(AB)\bigr|\le \Tr(A|B|)$
%  for $A\ge0$ trace class to get
% \[
% \bigl|\Tr\bigl(\Pi_M\sigma_t(H_{0,n})\Pi_M\, w_d(x_i-x_j)\bigr)\bigr|
% \le
% \Tr\bigl(\Pi_M\sigma_t\Pi_M\, |w_d(x_i-x_j)|\bigr).
% \]
% Again, since \(\Pi_M\) commutes with \(\sigma_t(H_{0,n})\), one has
% \[
% 0\le \Pi_M\sigma_t(H_{0,n})\Pi_M=\sigma_t(H_{0,n})\Pi_M\le \sigma_t(H_{0,n}).
% \]
% Hence
% \[
% \Tr\bigl(\Pi_M\sigma_t(H_{0,n})\Pi_M\, |w_d(x_i-x_j)|\bigr)
% \le
% \Tr\bigl(\sigma_t(H_{0,n})\, |w_d(x_i-x_j)|\bigr)
% =
% \widetilde I_{d,t}
% \]
% which proves the second bound.
% \end{proof}
% \noindent Similarly, we derive lower bounds on the free energies $F_t(H_n)$ and $F_t(H_{n,M})$:
% \begin{lemma}[Lower bounds on the free energy via free Gibbs states]\label{lowerboundsnow}
% Fix $t>0$ and assume that $\alpha_n =\mathcal{O}(1/n)$ Then, there exist constants $c,C>0$ such that
% \begin{align}
% F_t(H_n),\, F_t(H_{n,M})\ge\left\{
% \begin{aligned}
% & nF_t(h)&d=3\\
% & nF_{tc}(h)-C \alpha_n\binom{n}{2}     &d=2
% \end{aligned}
% \right.~~.
% \end{align}


% \end{lemma}

% \begin{proof}

% We start with the case \(d=3\), so $w_3(x)=|x|^{-1}\ge0$. Therefore both interaction forms are nonnegative:
% $W_n\ge0$ and $W_{n,M}\ge0$. Consequently
% \[
% H_n\ge H_{0,n}
% \qquad\text{ and }\qquad 
% H_{n,M}\ge H_{0,n}
% \]
% in the sense of quadratic forms. From the Golden-Thompson inequality, it follows that
% \[
% \mathcal Z_t(H_n)\le \Tr(e^{-tH_{0,n}})=\mathcal Z_t(h)^n
% \qquad\text{ and }\qquad 
% \mathcal Z_{n,M}(t)\le \mathcal Z_t(h)^n.
% \]
% We directly obtain 
% \[
% F_t(H_n)
% \ge
% nF_t(h)
% \qquad \text{ and }\qquad  F_t(H_{n,M})
% \ge
% nF_t(h).
% \]

% \noindent
% Now assume \(d=2\), so
% $w_2(x)=-\log|x|$. This kernel is not pointwise nonnegative, so the previous argument must be modified. We use the elementary estimate: for every \(\eta>0\) there exists \(C_\eta<\infty\) such that
% \begin{equation}
% \label{eq:log-bound-pointwise}
% -\log r\ge -\eta r^2-C_\eta
% \qquad (r>0).
% \end{equation}
% Applying \eqref{eq:log-bound-pointwise} with \(r=|x-y|\), and using $|x-y|^2\le 2|x|^2+2|y|^2$, we infer for any $x\ne y$
% \[
% -\log|x-y|
% \ge
% -2\eta(|x|^2+|y|^2)-C_\eta.
% \]
% Since \(h=-\Delta+|x|^2\ge |x|^2\) as quadratic forms, it follows that
% \begin{equation}
% \label{eq:log-form-bound}
% -\log|x-y|
% \ge
% -2\eta(h_x+h_y)-C_\eta
% \end{equation}
% as quadratic forms on \(L^2(\R^2\times\R^2)\). Summing \eqref{eq:log-form-bound} over all pairs \(1\le i<j\le N\), we obtain
% \[
% W_n
% \ge
% -2\eta\alpha_n\sum_{i<j}(h_i+h_j)-C_\eta\alpha_n\binom{n}{2}.
% \]
% Now each \(h_k\) appears exactly \(n-1\) times in the sum \(\sum_{1\le i<j\le n}(h_i+h_j)\), so $\sum_{i<j}(h_i+h_j)=(n-1)H_{0,n}$. Therefore
% \begin{equation}
% \label{eq:HN-lower-d2}
% H_n=H_{0,n}+W_n \ge
% \big(1-2\eta\alpha_n(n-1)\big)H_{0,n}-C_\eta\alpha_n\binom{n}{2}.
% \end{equation}
% Choose \(\eta>0\) so small that $1-2\eta\alpha_n(n-1)>0$, and set $c_\eta:=1-2\eta\Gamma$. Then \eqref{eq:HN-lower-d2} becomes
% \[
% H_n\ge c_\eta H_{0,n}-C_\eta \alpha_n\binom{n}{2}.
% \]
% Exactly the same estimate holds for \(H_{n,M}\). Indeed, since \(\Pi_M\) commutes with \(H_{0,n}\), from \eqref{eq:log-form-bound} we obtain
% \[
% \Pi_M\,(-\log|x_i-x_j|)\,\Pi_M
% \ge
% -2\eta\,\Pi_M(h_i+h_j)\Pi_M-C_\eta\Pi_M
% \ge
% -2\eta(h_i+h_j)-C_\eta,
% \]
% because \(0\le \Pi_M\le 1\) and \(\Pi_M(h_i+h_j)\Pi_M\le h_i+h_j\) on the form domain. Summing again over pairs yields
% \begin{equation}
% \label{eq:HNM-lower-d2}
% H_{n,M}\ge c_\eta H_{0,n}-C_\eta \alpha_n\binom{n}{2}.
% \end{equation}
% Now \eqref{eq:HN-lower-d2} together with the monotonicity of $A\mapsto \Tr(e^A)$ implies
% \[
% \Tr(e^{-tH_n}),\,\Tr(e^{-tH_{n,M}})\le e^{t C_\eta \alpha_n\binom{n}{2}}\Tr(e^{-tc_\eta H_{0,n}}),
% \]
% which directly yields
% \begin{equation}
% \label{eq:ZN-upper-d2-detailed}
% F_t(H_n),\,F_t(H_{n,M})\ge -C_\eta \alpha_n\binom{n}{2}    + nF_{tc_\eta}(h).
% \end{equation}


% \end{proof}


% \noindent We are ready to prove the required bounds on the average energy appearing in \eqref{almostthereyes}:

% \begin{lemma}[Bounding the average oscillator energy]
% \label{lem:uniform-H0-bound-abstract-detailed}
% Fix \(d\in\{2,3\}\), \(\beta>0\) and assume that $\alpha_n=\mathcal{O}(1/n)$. Then there exists a constant \(C_\beta<\infty\) and independent of \(n\) and \(M\), such that
% \[
% \Tr(\rho_\beta(H_{n}) H_{0,n})\le C_\beta n
% \qquad\text{ and }\qquad
% \Tr(\rho_\beta(H_{n,M}) H_{0,n})\le C_\beta n.
% \]
% \end{lemma}

% \begin{proof}
% We start with $d=3$. Applying Lemma \ref{lem:thermodynamic-inequality} with \(K=H_n\) and \(t=\beta\), we find
% \[
% \Tr(\sigma_\beta(H_n) H_n)
% \le
% 2F_\beta(H_n)-F_{\beta/2}(H_n).
% \]
% Using the upper bound for \(F_\beta(H_n)\) derived in Lemma \ref{lem:lower-bound-partition-gibbs}, the lower bound for \(F_\beta(H_n)\) from Lemma \ref{lowerboundsnow}, and that \(H_n\ge H_{0,n}\), we get
% \[
% \Tr(\sigma_\beta(H_n) H_{0,n})
% \le
% 2nF_\beta(h)+2\alpha_n n (n-1)\widetilde{I}_{3,\beta}-nF_{\beta/2}(h).
% \]
% The same argument yields an identical bound for \(\sigma_\beta(H_{n,M})\).
% \noindent
% Next, we consider $d=2$: again, using the upper bound for \(F_\beta(H_n)\) derived in Lemma \ref{lem:lower-bound-partition-gibbs} and the lower bound for \(F_\beta(H_n)\) from Lemma \ref{lowerboundsnow}, combined with Lemma \ref{lem:thermodynamic-inequality} with \(K=H_n\), we get
% \[
% \Tr(\rho_n H_n)
% \le
% 2nF_\beta(h)+2\alpha_n n (n-1)\widetilde{I}_{2,\beta}-nF_{\beta c/2}(h)+C\alpha_n\binom{n}{2}.
% \]
% Finally, from \eqref{eq:HN-lower-d2} we have
% \begin{equation}
% \label{eq:HN-lower-d3}
% H_n=H_{0,n}+W_n \ge
% c_\eta H_{0,n}-C_\eta\alpha_n\binom{n}{2}
% \end{equation}
% which does not change the scaling of the thermal average energy. Once again, the same proof, yields an identical bound for $\sigma_\beta(H_{n,M})$, and ends the proof. 
% \end{proof}

% \noindent We are now ready to state and prove the main theorem of this section:

% \begin{theorem}[Free energy perturbation bound]\label{thmfreenregyblabla}
% Fix $d\in\{2,3\}$, $\beta>0$, and let $\alpha_n:=n^{-\gamma}$ for $\gamma\ge 1$. There exists a constant $C_\beta<\infty$, independent of $n,M\ge 1$, such that
% \[
% |F_\beta(H_n)-F_{\beta}(H_{n,M})|
% \le C_\beta\,n^{3-\gamma}M^{-\frac{1}{4d}}.
% \]
% \end{theorem}

% \begin{proof}
% The bound is a direct consequence of \eqref{almostthereyes}, Lemma \ref{lem:uniform-H0-bound-abstract-detailed} and the fact that $\lambda_m=\Omega(m^{1/d})$.
% \end{proof}


\section{Finite rank perturbation analysis}

\subsection{Partition functions}\label{perturbationpartitionfunction}

\noindent In this section we construct a general perturbative framework for the truncation of the interaction to low-energy one-body modes. 
We start with the following abstract result: we recall that, given a Hamiltonian $H$ satisfying the Gibbs Hypothesis, we denote $\mathcal{Z}_\beta(H):=\Tr(e^{-\beta H})$,  $F_\beta(H):=-\beta^{-1}\log\mathcal{Z}_\beta(H)$ and $\sigma_\beta(H)=e^{-\beta H}/\mathcal{Z}_\beta(H)$.
We also recall that the form domain of a self-adjoint operator $H\ge -h_0$ corresponds to the domain $D((H+h_0+1)^{1/2})$.

\begin{lemma}[Perturbation bound for free energies under relative form convergence]
\label{lem:free-energy-relative-form-convergence}
 Let \(H_{0}\) be a self-adjoint operator on \(\mathcal H\), bounded from below, $H_0\ge -h_0$, with compact resolvent. Let \(W\) and \(W'\) be symmetric quadratic forms on the form domain \(Q(H_{0})\), and define the form sums
\[
H:=H_{0}+ W,
\qquad
H':=H_{0}+ W'.
\]
Assume further that the forms \(W\) and \(W'\) are \(H_{0}\)-form bounded with relative bound strictly smaller than \(1\). Let
\begin{equation}
\label{eq:DeltaNM}
\Delta:=H-H'=W-W'
\end{equation}
in the sense of quadratic forms, and assume that there exist numbers \(\varepsilon\ge0\) and \(c\ge0\), such that
\begin{equation}
\label{eq:relative-form-smallness-delta}
|\Delta[\psi]|
\le
\varepsilon \,\langle \psi,H_{0}\psi\rangle + c\|\psi\|^2
\qquad\text{ for all }\qquad \psi\in Q(H_{0}).
\end{equation}
Then, the Gibbs states associated to $H$, resp.~$H'$, as $\sigma_\beta(H)$, resp.~$\sigma_\beta(H')$ are well-defined as trace-class operators, and
\begin{equation}
\label{eq:sharper-max-estimate}
|F_\beta(H)-F_\beta(H')|
\le
{\varepsilon}
\max\bigl\{\Tr(\sigma_\beta(H) H_{0}),\Tr(\sigma_\beta(H')H_{0})\bigr\}
+{c}.
\end{equation}
\end{lemma}

\begin{proof}

Since the forms \(W\) and \(W'\) are \(H_{0}\)-form bounded with a relative bound strictly smaller than \(1\), \(H\) and \(H'\) are self-adjoint and bounded from below on \(Q(H_{0})\) with a compact resolvent. Hence, the Gibbs operators \(e^{-\beta H}\) and \(e^{-\beta H'}\) are trace class. Next, by the Bogoliubov inequality \cite{Dereziski2003}
\begin{equation}
\label{eq:PB}
\log \frac{\Tr(e^{X+Y})}{\Tr(e^{X})}
\ge
\frac{\Tr(e^{X}Y)}{\Tr(e^{X})}
\end{equation}
for self-adjoint \(X,Y\) with $X:=-\beta H$, $Y:=-\beta(H'-H)$ as well as for $X:=-\beta H'$, $Y:=-\beta(H-H')$, we get 
\[
\Tr(\sigma_\beta(H)\Delta)
\le \frac1\beta\log\frac{\mathcal{Z}_\beta(H')}{\mathcal{Z}_\beta(H)}
\le \Tr(\sigma_\beta(H')\Delta).
\]
Hence
\begin{equation}
\label{eq:free-energy-max-step}
|F_\beta(H)-F_{\beta}(H')|
\le
\max\Bigl\{
|\Tr(\sigma_\beta(H)\Delta)|,
|\Tr(\sigma_\beta(H')\Delta)|
\Bigr\}.
\end{equation}
Next, applying \eqref{eq:relative-form-smallness-delta} in the spectral decompositions of \(\sigma_\beta(H)\) and \(\sigma_\beta(H')\) gives
\[
|\Tr(\sigma_\beta(H)\Delta)|
\le \varepsilon\,\Tr(\sigma_\beta(H)H_{0})+c\qquad \text{ and }\qquad 
|\Tr(\sigma_\beta(H')\Delta)|
\le \varepsilon\,\Tr(\sigma_\beta(H')H_{0})+c.
\]
Inserting these estimates into \eqref{eq:free-energy-max-step}, we obtain
\[
|F_\beta(H)-F_\beta(H')|
\le
{\varepsilon}
\max\bigl\{\Tr(\sigma_\beta(H)H_{0}),\Tr(\sigma_\beta(H')H_{0})\bigr\}
+c.
\]
\end{proof}

\noindent For now on, we consider the $n$-particle Hamiltonian
\[
H_n:=H_{0,n}+W_n
\]
on \(L^2((\R^d)^n)\), where $h:=-\Delta+|x|^2$ over $L^2(\mathbb{R}^d)$, $H_{0,n}:=\sum_{i\in[n]}h_i$ and for $\psi\in Q(H_{0,n})$, denote
\[
W_n[\psi]=\sum_{1\le i<j\le n}\alpha_{n,i,j}W_{ij}[\psi],
\]
where
\[
W_{ij}[\psi]:=\int_{\R^{nd}}w_d(x_i-x_j)|\psi(x)|^2 \ dx,\qquad
w_d(y):=
\begin{cases}
-\log|y|, & d=2,\\[0.3em]
|y|^{-1}, & d=3.
\end{cases}
\]
With the convention \(\alpha_{n,j,i}:=\alpha_{n,i,j}\) for \(j>i\), let further
\[
A_n:=\max_{i\in[n]}\sum_{j\in[n]\setminus\{i\}}|\alpha_{n,i,j}|,
\qquad
B_n:=\sum_{1\le i<j\le n}|\alpha_{n,i,j}|.
\]
Let further $P_M$ be the spectral projector onto the first $M$ eigenvalues \((\lambda_k)_{k=1}^M\) of \(h\). Define \(\Pi_M:=P_M^{\otimes n}\), \(\Pi_M^\perp:=1-\Pi_M\), and consider the form
\[W_{n,M}[\psi]:=W_n[\Pi_M\psi]\]
and the associated Hamiltonian $H_{n,M}$. In the next result, we establish relative form boundedness of $\Delta_{n,M}:=W_n-W_{n,M}$ with respect to $H_{0,n}$:

\begin{lemma}[Product cutoff estimate without scaling assumption]
\label{lem:product-cutoff-riesz-no-scaling}
Given $d\in\{2,3\}$, for all \(n,M\) and \(\psi\in Q(H_{0,n})\),
\begin{align}
|W_n[\psi]|&\lesssim \,B_n\langle\psi,(H_{0,n}+1)\psi\rangle,\label{eq111}\\
|W_n[\Pi_M^\perp\psi]|&\lesssim\,B_n(\lambda_{M+1}+1)^{-1/2}\langle\psi,(H_{0,n}+1)\psi\rangle,\label{eq222}\\
|\Delta_{n,M}[\psi]|&\lesssim\,B_n(\lambda_{M+1}+1)^{-1/4}\langle\psi,(H_{0,n}+1)\psi\rangle.\label{eq333}
\end{align}
\end{lemma}

\begin{proof}
 Fix $1\le i<j\le n$ and $d=3$. Introduce the relative variable $r:=2^{-1/2}{(x_i-x_j)}$, so that ${|x_i-x_j|^{-1}}=2^{-1/2}{|r|^{-1}}$. By Hardy's inequality in $\mathbb R^3$,
\[
\left\|\frac{u}{|r|}\right\|^2\le 4\|\nabla_r u\|^2.
\]
Hence, by Cauchy--Schwarz and Young's inequalities, for every \(\eta>0\),
\begin{align*}
\left\langle u,\frac{1}{|x_i-x_j|}u\right\rangle
=
2^{-1/2}\left\langle u,\frac{1}{|r|}u\right\rangle
\le
\sqrt2\,\|u\|\,\|\nabla_r u\|
&\le
\eta\|\nabla_r u\|^2+ (2\eta)^{-1}\|u\|^2\\
&= \eta\|\nabla_r u\|^2+ C_\eta\|u\|^2,
\end{align*}
where we defined $C_\eta = (2\eta)^{-1}$ in the last equality.
Since
\[
\|\nabla_r u\|^2\le \langle u,(-\Delta_{x_i}-\Delta_{x_j})u\rangle
\le
\langle u,\bigl(-\Delta_{x_i}-\Delta_{x_j}+|x_i|^2+|x_j|^2\bigr)u\rangle,
\]
we obtain
\begin{equation}
\frac{1}{|x_i-x_j|}
\le
\eta\bigl(-\Delta_{x_i}-\Delta_{x_j}+|x_i|^2+|x_j|^2\bigr)+(2\eta)^{-1} \label{usefulestimate111}
\end{equation}
as quadratic forms. Taking \(\eta=1\), and using \(h_i+h_j+1\le 2(H_{0,n}+1)\), we get
\[
|W_n[\psi]|
\lesssim
\sum_{1\le i<j\le n}|\alpha_{n,i,j}|\langle\psi,(h_i+h_j+1)\psi\rangle
\lesssim
B_n\langle\psi,(H_{0,n}+1)\psi\rangle.
\]
In dimension \(d=2\), let \(u\in Q(h_1+h_2)\). Fix \(0<\varepsilon<1\). Since
\[
|\log r|\le C_\varepsilon\bigl(1+r^{-\varepsilon}+r^\varepsilon\bigr),\qquad r>0,
\]
it suffices to estimate these three terms separately. First,
\[
\int_{\R^2\times\R^2}|u(x,y)|^2\,dx\,dy=\|u\|^2
\le
\langle u,(h_1+h_2+1)u\rangle.
\]
Next, using the change of variables \(r=x-y\), \(s=x+y\), we obtain
\[
\int_{\R^2\times\R^2}|x-y|^{-\varepsilon}|u(x,y)|^2\,dx\,dy
=
\frac{1}{4}\int_{\R^2\times\R^2}|r|^{-\varepsilon}|u(r,s)|^2\,dr\,ds.
\]
Set
\[
F(r):=\|u(r,\cdot)\|_{L^2(\R^2_s)}.
\]
Then
\[
\int_{\R^2\times\R^2}|r|^{-\varepsilon}|u(r,s)|^2\,dr\,ds
=\int_{\R^2}|r|^{-\varepsilon}F(r)^2\,dr.
\]
Choose \(p>1\) such that \(\varepsilon p<2\), and let \(q=\frac p{p-1}\) its Hölder conjugate. Since
\(|r|^{-\varepsilon}\in L^p(B_1(\R^2))\) for $B_1(\mathbb{R}^2)$ the unit ball in $\mathbb{R}^2$, Hölder's inequality in the \(r\)-variable yields
\[
\int_{|r|\le1}|r|^{-\varepsilon}F(r)^2\,dr
\le
\bigl\||r|^{-\varepsilon}\bigr\|_{L^p(B_1(\mathbb{R}^2))}\|F^2\|_{L^q(\R^2)}
=
C\|F\|_{L^{2q}(\R^2)}^2.
\]
Now \(2q<\infty\), so by the Sobolev embedding
\[
H^1(\R^2_r;L^2(\R^2_s))\hookrightarrow L^{2q}(\R^2_r;L^2(\R^2_s))
\]
we obtain
\[
\|F\|_{L^{2q}(\R^2)}=\|u\|_{L^{2q}(\mathbb{R}_r^2,L^2(\mathbb{R}_s^2))}
\lesssim\|u\|_{H^1(\R^2_r;L^2(\R^2_s))}.
\]
Hence
\[
\int_{|r|\le1}\int_{\R^2}|r|^{-\varepsilon}|u(r,s)|^2\,ds\,dr
\lesssim
\|u\|_{H^1(\R^2_r;L^2(\R^2_s))}^2.
\]
Since \(|r|^{-\varepsilon}\le1\) on \(\R^2\setminus B_1(\mathbb{R}^2)\), we also control
\[
\int_{|r|>1}\int_{\R^2}|r|^{-\varepsilon}|u(r,s)|^2\,ds\,dr
\le
\|u\|_{L^2(\R^4)}^2.
\]
Therefore
\[
\int_{\R^2\times\R^2}|r|^{-\varepsilon}|u(r,s)|^2\,dr\,ds
\lesssim
\Bigl(\|u\|_{H^1(\R^2_r;L^2(\R^2_s))}^2+\|u\|_{L^2(\R^4)}^2\Bigr).
\]
By the Sobolev embedding \(H^1(\R^4)\hookrightarrow L^4(\R^4)\),
\[
\|u\|_{L^4(\R^4)}^2
\lesssim\|u\|_{H^1(\R^4)}^2
\lesssim
\langle u,(h_1+h_2+1)u\rangle.
\]
Finally, since \(0<\varepsilon<2\), there exist a constants $C_\varepsilon,C_\varepsilon'>0$ such that
\[
|x-y|^\varepsilon\le C_{\varepsilon}\bigl(1+|x|^\varepsilon+|y|^\varepsilon\bigr)
\le C'_{\varepsilon}\bigl(1+|x|^2+|y|^2\bigr),
\]
and therefore
\[
\int_{\R^2\times\R^2}|x-y|^\varepsilon |u(x,y)|^2\,dx\,dy
\le
C_\varepsilon'\langle u,(h_1+h_2+1)u\rangle.
\]
Putting the three estimates together yields
\[
\int_{\R^2\times\R^2}|\log|x-y||\,|u(x,y)|^2\,dx\,dy
\lesssim
\langle u,(h_1+h_2+1)u\rangle.
\]
Summing over all combinations of $x=x_i,y=x_j$ yields then 
\[
|W_n[\psi]|\lesssim B_n\langle\psi,(H_{0,n}+1)\psi\rangle.
\]
This ends the proof of \eqref{eq111}. For \eqref{eq222}, we may proceed for both $d=2,3$ as follows:
instead of bounding the interaction form entirely by the energy right away, we first apply the Cauchy--Schwarz inequality to the interaction potential. For $d=3$, using Hardy's inequality $\| |x_i-x_j|^{-1}\phi \| \lesssim \langle \phi, (h_i+h_j)\phi \rangle^{1/2}$, and for $d=2$ using the analogous interpolation from the Gagliardo--Nirenberg inequality derived above, we obtain the fractional form bound for any $\phi \in Q(H_{0,n})$:
\[
|W_n[\phi]| \lesssim B_n \|\phi\| \langle \phi, (H_{0,n}+1)\phi \rangle^{1/2}.
\]
We evaluate this bound on the high-energy projection $\phi = \Pi_M^\perp\psi$. Since $(H_{0,n}+1)\Pi_M^\perp\ge(\lambda_{M+1}+1)\Pi_M^\perp$, the norm decays as:
\[ 
\|\Pi_M^\perp\psi\| \le (\lambda_{M+1}+1)^{-1/2}\langle\psi,(H_{0,n}+1)\psi\rangle^{1/2}.
\]
Furthermore, since $[\Pi_M^\perp,H_{0,n}]=0$ and $\Pi_M^\perp \le 1$, we have $\langle \Pi_M^\perp\psi, (H_{0,n}+1)\Pi_M^\perp\psi\rangle \le \langle \psi, (H_{0,n}+1)\psi\rangle$. Substituting these into our fractional form bound yields
\begin{align}
|W_n[\Pi_M^\perp\psi]| &\lesssim B_n \|\Pi_M^\perp\psi\| \langle \Pi_M^\perp\psi, (H_{0,n}+1)\Pi_M^\perp\psi \rangle^{1/2}\nonumber\\
&\le B_n \left( (\lambda_{M+1}+1)^{-1/2} \langle \psi, (H_{0,n}+1)\psi \rangle^{1/2} \right) \langle \psi, (H_{0,n}+1)\psi \rangle^{1/2}\nonumber\\
&= (\lambda_{M+1}+1)^{-1/2}B_n\langle \psi,(H_{0,n}+1)\psi\rangle,\label{eq2222}
\end{align}
which ends the proof of \eqref{eq222}.


For \eqref{eq333}, we associate to \(W_n\) the sesquilinear form
\[
W_n(\varphi,\chi)
:=
\sum_{1\le i<j\le n}
\alpha_{n,i,j}\int_{(\R^d)^n}w_d(x_i-x_j)\,\varphi(x)\,\overline{\chi(x)}\,dx,
\]
so that \(W_n[\psi]=W_n(\psi,\psi)\). Hence, writing \(P:=\Pi_M\) and \(Q:=\Pi_M^\perp\),
\[
W_n[\psi]-W_{n,M}[\psi]
=
W_n[P\psi+Q\psi]-W_n[P\psi]
=
W_n[Q\psi]+2\Re W_n(P\psi,Q\psi).
\]
Similarly, if \(\widetilde W_n\) denotes the positive sesquilinear form with kernel
\(|w_d|\) and coefficients \(|\alpha_{n,i,j}|\), then combining the Cauchy-Schwarz inequality, \eqref{eq111} applied to \(P\psi\), and \eqref{eq222} applied to \(Q\psi\), we get
\begin{align}
|W_n(P\psi,Q\psi)|&\le \widetilde{W}_n[P\psi]^{1/2}\widetilde{W}_n[Q\psi]^{1/2}\nonumber\\
&\lesssim B_n(\lambda_{M+1}+1)^{-1/4}\langle \psi,(H_{0,n}+1)\psi\rangle.\nonumber
\end{align}
Finally, combining the last bound with the estimate for \(W_n[Q\psi]\) derived in \eqref{eq222} yields the claimed bound \eqref{eq333}.
\end{proof}

\noindent \noindent Combining Lemmas \ref{lem:free-energy-relative-form-convergence} and \ref{lem:product-cutoff-riesz-no-scaling}, we get that the difference of free energies is controlled by
\begin{align}\label{almostthereyes}
\frac{|F_\beta(H_n)-F_\beta(H_{n,M})|}{B_n}\lesssim (\lambda_{M+1}+1)^{-\frac{1}{4}} \max \left\{\Tr(\sigma_\beta(H_n)H_{0,n}),\Tr(\sigma_\beta(H_{n,M})H_{0,n})\right\}.
\end{align}
It remains to control each of the average energies appearing in the previous bound. We first bound them, via Lemma \ref{lem:thermodynamic-inequality}, by the difference between the free energies of the perturbed and unperturbed models. We then show in Lemmas \ref{lem:lower-bound-partition-gibbs} and \ref{lowerboundsnow} how to control this free-energy difference. The final energy bounds are stated in Lemma \ref{lem:uniform-H0-bound-abstract-detailed}.

\begin{lemma}
\label{lem:thermodynamic-inequality}
Let \(K\) be a self-adjoint operator on a Hilbert space, bounded from below with smooth partition function. Then, for every \(t>0\),
\begin{equation}
\label{eq:UK-logZ}
\Tr(\sigma_{t}(K)K)\le 2F_t(K)-F_{t/2}(K).
\end{equation}
\end{lemma}

\begin{proof}
Writing $U_K(t)=\Tr(\sigma_t(K)K)$,
we immediately get
\[
-\partial_t\log \mathcal Z_t(K)
=
-\frac{\partial_t\mathcal Z_t(K)}{\mathcal Z_t(K)}
=
\frac{\Tr(Ke^{-tK})}{\Tr(e^{-tK})}
=
U_K(t).
\]
Moreover, by the quotient rule,
\begin{align}
U_K'(t)
&=
\frac{\Tr(Ke^{-tK})^2-\Tr(K^2e^{-tK})\Tr(e^{-tK})}{\Tr(e^{-tK})^2}=
-\Tr(\sigma_{t}(K)K^2)+\bigl(\Tr(\sigma_{t}(K)K)\bigr)^2.
\end{align}
 Expanding \((K-U_K(t))^2\) yields $\Tr\bigl(\sigma_{t}(K)(K-U_K(t))^2\bigr)
=
\Tr(\sigma_{t}(K)K^2)-U_K(t)^2$. Thus
\[
U_K'(t)
=
-\Tr\bigl(\sigma_{t}(K)(K-U_K(t))^2\bigr)\le 0,
\]
This proves that \(t\mapsto U_K(t)\) is decreasing. Finally, fix \(t>0\). Since \(U_K\) is decreasing on \((0,\infty)\), we have
\[
U_K(t)\le \frac{2}{t}\int_{t/2}^t U_K(s)\,ds.
\]
Using \(U_K(s)=-\partial_s\log \mathcal Z_s(K)\), we obtain
\[
\frac{2}{t}\int_{t/2}^t U_K(s)\,ds
=
-\frac{2}{t}\int_{t/2}^t \partial_s\log \mathcal Z_s(K)\,ds
=
\frac{2}{t}\Bigl(\log \mathcal Z_{t/2}(K)-\log \mathcal Z_t(K)\Bigr),
\]
which is exactly \eqref{eq:UK-logZ}.
\end{proof}
\noindent Next, we control the free energy difference obtained in the previous Lemma. In what follows, we
denote the one-body $\rho_t(x):=\langle x|\sigma_t(h)|x\rangle$. Since \(h=-\Delta+|x|^2\), the kernel $\rho_t(x):=\langle x|\sigma_t(h)|x\rangle$ of \(e^{-th}\) is given by the Mehler formula. In particular, \(\rho_t(x)\) is smooth, strictly positive, and has Gaussian decay as \(|x|\to\infty\):
\[
\rho_t(x)\le C_t e^{-c_t|x|^2}
\]
for suitable constants \(c_t,C_t>0\). Therefore, for any $d\in\{2,3\}$,
\[
\widetilde I_{d,t}
:=
\int_{\mathbb R^{2d}}|w_d(x-y)|\,\rho_t(x)\rho_t(y)\,dx\,dy<\infty.
\]

\begin{lemma}[Upper bounds on the free energy via free Gibbs states]
\label{lem:lower-bound-partition-gibbs}Given $d\in\{2,3\}$, fix $t>0$. Then
\begin{equation}
\label{eq:lower-ZN-lemma}
F_t(H_n)\le nF_t(h)+B_n\widetilde I_{d,t},
\qquad \text{ and }\qquad 
F_t(H_{n,M})\le nF_t(h)+B_n\widetilde I_{d,t}.
\end{equation}
\end{lemma}

\begin{proof}

By a use of the Bogoliubov inequality \eqref{eq:PB}, with $X=-tH_{0,n}$ and $Y=-t(H_n-H_{0,n})$,
\[
F_t(H_{n})
\le
{n}F_t(h)+\Tr(\sigma_t(H_{0,n}) W_n).
\]
Next, we estimate the interaction term: for each pair \(i<j\) one has
\[
\Tr\bigl(\sigma_t(H_{0,n})\, w_d(x_i-x_j)\bigr)
=
\int_{\R^{2d}} w_d(x-y)\rho_t(x)\rho_t(y)\,dx\,dy,
\]
independently of $i<j$. Hence
\begin{align}
\Tr(\sigma_t(H_{0,n}) W_n)
&=
\sum_{1\le i<j\le n}\alpha_{n,i,j}
\int_{\R^{2d}} w_d(x-y)\rho_t(x)\rho_t(y)\,dx\,dy\nonumber\\
&\le
\sum_{1\le i<j\le n}|\alpha_{n,i,j}|
\int_{\R^{2d}} |w_d(x-y)|\,\rho_t(x)\rho_t(y)\,dx\,dy
\nonumber\\
&=
B_n\widetilde I_{d,t}.\nonumber
\end{align}
The proof for \(H_{n,M}\) is almost identical: by \eqref{eq:PB},
\[
F_t(H_{n,M})\le nF_t(h)+\Tr(\sigma_t(H_{0,n}) W_{n,M}).
\]
Then, since \(\Pi_M\) commutes with \(H_{0,n}\), \(\Pi_M\) also commutes with \(\sigma_t(H_{0,n})\) and
\[
\Tr(\sigma_t(H_{0,n}) W_{n,M})
=
\sum_{1\le i<j\le n}\alpha_{n,i,j}\Tr\bigl(\Pi_M\sigma_t(H_{0,n})\Pi_M\, w_d(x_i-x_j)\bigr).
\]
Since \(\Pi_M\sigma_t(H_{0,n})\Pi_M\ge0\) is trace class, we may use $\bigl|\Tr(AB)\bigr|\le \Tr(A|B|)$
 for $A\ge0$ trace class to get
\[
\bigl|\Tr\bigl(\Pi_M\sigma_t(H_{0,n})\Pi_M\, w_d(x_i-x_j)\bigr)\bigr|
\le
\Tr\bigl(\Pi_M\sigma_t(H_{0,n})\Pi_M\, |w_d(x_i-x_j)|\bigr).
\]
Again, since \(\Pi_M\) commutes with \(\sigma_t(H_{0,n})\), one has
\[
0\le \Pi_M\sigma_t(H_{0,n})\Pi_M=\sigma_t(H_{0,n})\Pi_M\le \sigma_t(H_{0,n}).
\]
Hence
\[
\Tr\bigl(\Pi_M\sigma_t(H_{0,n})\Pi_M\, |w_d(x_i-x_j)|\bigr)
\le
\Tr\bigl(\sigma_t(H_{0,n})\, |w_d(x_i-x_j)|\bigr)
=
\widetilde I_{d,t}
\]
which proves the second bound.
\end{proof}
\noindent Similarly, we derive lower bounds on the free energies $F_t(H_n)$ and $F_t(H_{n,M})$:
\begin{lemma}[Lower bounds on the free energy via free Gibbs states]\label{lowerboundsnow}
For $t>0$ and $d\in\{2,3\}$ we have
\begin{align}
F_t(H_n),\, F_t(H_{n,M})\ge 
 nF_{t/2}(h)- CA_n B_n\,.
\end{align}
For some universal constant $C\ge 0.$
\end{lemma}

\begin{proof}

We start with the case \(d=3\), so $w_3(x)=|x|^{-1}\ge0$. For every \(\eta>0\), the estimate proved in \eqref{usefulestimate111} gives
\[
\frac{1}{|x_i-x_j|}
\le
\eta\bigl(-\Delta_{x_i}-\Delta_{x_j}+|x_i|^2+|x_j|^2\bigr)+(2\eta)^{-1}
\]
as quadratic forms. Therefore
\[
\alpha_{n,i,j}|x_i-x_j|^{-1}
\ge
-|\alpha_{n,i,j}|\Bigl(\eta(h_i+h_j)+(2\eta)^{-1}\Bigr).
\]
Summing over all pairs \(1\le i<j\le n\), we obtain
\[
W_n
\ge
-\eta\sum_{1\le i<j\le n}|\alpha_{n,i,j}|(h_i+h_j)-(2\eta)^{-1} B_n.
\]
Now
\[
\sum_{1\le i<j\le n}|\alpha_{n,i,j}|(h_i+h_j)
=
\sum_{i=1}^n\Bigl(\sum_{j\in[n]\setminus\{i\}}|\alpha_{n,i,j}|\Bigr)h_i
\le
A_n H_{0,n}.
\]
Choosing \(\eta = \frac{1}{2A_n},\) we get
\begin{equation}
\label{eq:HN-lower-d3-general}
H_n=H_{0,n}+W_n \ge (1- \eta A_n)H_{0,n} - (2\eta)^{-1}B_n = \frac{1}{2}H_{0,n} - A_nB_n
\end{equation}
Exactly the same estimate holds for \(H_{n,M}\). Indeed, since \(\Pi_M\) commutes with \(H_{0,n}\), from the same two-body bound we obtain
\[
\Pi_M\,|x_i-x_j|^{-1}\,\Pi_M
\le
\eta\,\Pi_M(h_i+h_j)\Pi_M+(2\eta)^{-1}\Pi_M
\le
\eta(h_i+h_j)+(2\eta)^{-1},
\]
because \(0\le \Pi_M\le 1\) and \(\Pi_M(h_i+h_j)\Pi_M\le h_i+h_j\) on the form domain. Summing again over pairs yields and taking the same choice for $\eta$ yields 
\begin{equation}
\label{eq:HNM-lower-d3-general}
H_{n,M}\ge \frac{1}{2}H_{0,n}- A_n B_n.
\end{equation}
Now \eqref{eq:HN-lower-d3-general} together with the monotonicity of $A\mapsto \Tr(e^A)$ implies
\[
\Tr(e^{-tH_n}),\,\Tr(e^{-tH_{n,M}})\le e^{t A_n B_n}\Tr(e^{-t H_{0,n}/2}),
\]
which directly yields
\[
F_t(H_n),\,F_t(H_{n,M})\ge - A_n B_n+ nF_{t/2}(h).
\]

\noindent
Now assume \(d=2\), so
$w_2(x)=-\log|x|$. This kernel is not pointwise nonnegative, so the previous argument must be modified. We use the following bound, whose proof we postpone to Lemma \ref{lem:two-sided-log-form-bound} below: for any $\eta>0$ we have that 
\begin{equation}
|\,\log|x-y|\,|
\le
\eta(h_x+h_y)+C \,\eta^{-1}
\end{equation}
as quadratic forms on \(L^2(\R^2\times\R^2)\) and for some universal constant $C\ge 0.$
\noindent As a consequence, if
\[
W_n=\sum_{1\le i<j\le n}\alpha_{n,i,j}\,(-\log|x_i-x_j|),
\qquad
B_n:=\sum_{1\le i<j\le n}|\alpha_{n,i,j}|,
\]
then for every \(\eta>0\),
\begin{equation}
\label{eq:signed-log-lower-bound}
W_n
\ge
-\eta\sum_{1\le i<j\le n}|\alpha_{n,i,j}|(h_i+h_j)-C \,\eta^{-1} B_n\ge -\eta A_n H_{0,n}-C \,\eta^{-1}B_n
\end{equation}
as quadratic forms. Thus, we obtain
\begin{equation}
\label{eq:HN-lower-d2}
H_n=H_{0,n}+W_n \ge
\big(1-\eta A_n\big)H_{0,n}-C \,\eta^{-1} B_n.
\end{equation}
Choosing once again $\eta = \frac{1}{2A_n}$ we see that \eqref{eq:HN-lower-d2} becomes
\begin{align}\label{fordequals2222}
H_n\ge  \frac{1}{2}H_{0,n}- 2C A_nB_n.
\end{align}
Again, exactly the same estimate holds for \(H_{n,M}\), using that \(\Pi_M\) commutes with \(H_{0,n}\). Thus, \eqref{eq:HN-lower-d2} together with the monotonicity of $A\mapsto \Tr(e^A)$ implies
\[
\Tr(e^{-tH_n}),\,\Tr(e^{-tH_{n,M}})\le e^{t 2C A_nB_n}\Tr(e^{-t H_{0,n}/2}),
\]
which directly yields
\begin{equation}
\label{eq:ZN-upper-d2-detailed}
F_t(H_n),\,F_t(H_{n,M})\ge -2CA_n B_n    + nF_{t/2}(h).
\end{equation}

\end{proof}








\begin{lemma}[Two-sided logarithmic form bound in dimension \(2\)]
\label{lem:two-sided-log-form-bound}
For \(\eta>0\) we have 
\begin{equation}
\label{eq:two-sided-log-form-bound}
\bigl|\langle u,(-\log|x-y|)u\rangle\bigr|
\le
\eta\,\langle u,(h_x+h_y)u\rangle + C \,\eta^{-1}\|u\|^2
\qquad\text{ for all }u\in Q(h_x+h_y),
\end{equation}
where \(h=-\Delta+|x|^2\) on \(L^2(\R^2)\) and for some universal constant $C\ge 0.$  Equivalently,
\begin{equation}
\label{eq:two-sided-log-operator-bound}
-|\,\log|x-y|\,|
\le
-\log|x-y|
\le
|\,\log|x-y|\,|
\le
\eta(h_x+h_y)+C \,\eta^{-1}
\end{equation}
as quadratic forms on \(L^2(\R^2\times\R^2)\).
\end{lemma}

\begin{proof}
Let \(u\in Q(h_x+h_y)\) and introduce the orthogonal change of variables
$r:=2^{-1/2}(x-y)$, $s:=2^{-1/2}(x+y)$, and define $v(r,s):=u\big(\frac{r+s}{\sqrt2},\frac{s-r}{\sqrt2}\big)$. Moreover, the map \((x,y)\mapsto(r,s)\) satisfies
\[
\|v\|_{L^2(\R^4)}=\|u\|_{L^2(\R^4)},
\qquad
-\Delta_x-\Delta_y=-\Delta_r-\Delta_s,
\qquad
|x|^2+|y|^2=|r|^2+|s|^2.
\]
Hence
\begin{equation}
\label{eq:energy-rs}
\langle u,(h_x+h_y)u\rangle
=
\int_{\R^4}\Bigl(|\nabla_r v|^2+|\nabla_s v|^2+(|r|^2+|s|^2)|v|^2\Bigr)\,dr\,ds.
\end{equation}
Also, since \(|x-y|=\sqrt2\,|r|\),
\[
|\log|x-y||
\le
C+|\log|r||.
\]
Therefore it suffices to prove that for every \(\eta>0\),
\begin{equation}
\label{eq:goal-r}
\int_{\R^4}|\log|r||\,|v(r,s)|^2\,dr\,ds
\le
\eta\,\langle u,(h_x+h_y)u\rangle + C \eta^{-1}\|u\|^2.
\end{equation}
We split the integral into the regions \(|r|\le 1\) and \(|r|>1\). For the region \(|r|>1\), we use that for every \(\eta>0\) we can bound
\[
\log t= \frac{1}{2}\log\left(\eta t^2/2\right) + \frac{1}{2}\log(2/\eta) \le \eta t^2 +  \eta^{-1},
\]
where for the last inequality we have used the elementary bound $\log x \le x.$
Hence
\begin{align}
\int_{|r|>1}|\log|r||\,|v|^2=
\int_{|r|>1}\log|r|\,|v|^2\le
\eta\int_{\R^4}|r|^2|v|^2\,dr\,ds
+
\eta^{-1}\|v\|^2.
\label{eq:large-r-log}
\end{align}
For the region \(|r|\le1\), fix \(0<\delta<2\). Since
\[
-\log t\le C_\delta\, t^{-\delta}
\qquad (0<t\le1),
\]
we obtain
\[
\int_{|r|\le1}|\log|r||\,|v(r,s)|^2\,dr\,ds
\le
C_\delta\int_{\R^2_s}\int_{|r|\le1}|r|^{-\delta}|v(r,s)|^2\,dr\,ds.
\]
Fix \(s\in\R^2\), and set \(f_s(r):=v(r,s)\). Choose \(p>1\) such that \(\delta p<2\), and let \(q=\frac{p}{p-1}\). Since \(|r|^{-\delta}\in L^p(B_1(\R^2))\), Hölder's inequality gives
\[
\int_{|r|\le1}|r|^{-\delta}|f_s(r)|^2\,dr
\le
\||r|^{-\delta}\|_{L^p(B_1(\R^2))}\|f_s\|_{L^{2q}(\R^2)}^2.
\]
Now by the Gagliardo--Nirenberg inequality in dimension \(2\),
\[
\|f_s\|_{L^{2q}(\R^2)}^2
\le
C_q\|\nabla_r f_s\|_{L^2(\R^2)}^{2/p}\|f_s\|_{L^2(\R^2)}^{2(1-1/p)}.
\]
Using Young's inequality, for every \(\eta>0\),
\[
\|\nabla_r f_s\|_{L^2(\mathbb{R}^2)}^{2/p}\|f_s\|_{L^2(\mathbb{R}^2)}^{2(1-1/p)}
\le
\frac{\eta}{C_\delta C_q\||r|^{-\delta}\|_{L^p(B_1(\R^2))}}\,\|\nabla_r f_s\|_{L^2(\mathbb{R}^2)}^2 + c_{p,\delta} \,\eta^{-1/(p-1)}\|f_s\|_{L^2(\mathbb{R}^2)}^2.
\]
Therefore, choosing $\delta = 1/2$ and $p= 2,$ yields after integrating in \(s\),
\begin{equation}
\label{eq:small-r-log}
\int_{|r|\le1}|\log|r||\,|v(r,s)|^2\,dr\,ds
\le
\eta\|\nabla_r v\|_{L^2(\R^4)}^2 + C'\, \eta^{-1}\|v\|_{L^2(\R^4)}^2.
\end{equation}
Combining \eqref{eq:large-r-log} and \eqref{eq:small-r-log}, and using \eqref{eq:energy-rs}, we get
\begin{align*}
\int_{\R^4}|\log|r||\,|v(r,s)|^2\,dr\,ds
&\le
\eta\Bigl(\|\nabla_r v\|_{L^2(\R^4)}^2+\||r|v\|_{L^2(\R^4)}^2\Bigr)+ C \eta^{-1}\|v\|^2\\
&\le
\eta\,\langle u,(h_x+h_y)u\rangle + C \eta^{-1}\|u\|^2,
\end{align*}
which proves \eqref{eq:goal-r}, hence \eqref{eq:two-sided-log-form-bound}. The operator inequalities in \eqref{eq:two-sided-log-operator-bound} follow immediately.
\end{proof}






















\noindent We are ready to prove the required bounds on the average energy appearing in \eqref{almostthereyes}:

\begin{lemma}[Bounding the average oscillator energy]
\label{lem:uniform-H0-bound-abstract-detailed}
Fix \(d\in\{2,3\}\). Then there exists a constants \(C_\beta<\infty\) and independent of \(n\) and \(M\), such that
\begin{align*}
\Tr(\sigma_\beta(H_{n}) H_{0,n})&\le C_\beta\,\left(n + (A_n+1)B_n\right), \\
\Tr(\sigma_\beta(H_{n,M}) H_{0,n})&\le C_\beta\,\left(n + (A_n+1)B_n\right).
\end{align*}
\end{lemma}

\begin{proof}
 Using the upper bound for \(F_\beta(H_n)\) derived in Lemma \ref{lem:lower-bound-partition-gibbs} and the lower bound for \(F_\beta(H_n)\) from Lemma \ref{lowerboundsnow}, combined with Lemma \ref{lem:thermodynamic-inequality} with \(K=H_n\), we get 
\[
\Tr(\sigma_\beta(H_n) H_n)
\le
2nF_\beta(h)+2B_n\widetilde{I}_{d,\beta}-nF_{\beta/4}(h)+ CA_n B_n.
\]
Finally, from \eqref{fordequals2222} and \eqref{eq:HN-lower-d3-general} we have
\begin{equation}
\label{eq:HN-lower-d3}
H_n=H_{0,n}+W_n \ge
\frac{1}{2} H_{0,n}- CA_nB_n
\end{equation}
from which we can conclude 
\begin{align*}
    \Tr(\sigma_\beta(H_n)H_{0,n})\le 2\Tr(\sigma_\beta(H_n)H_{n}) + 2CA_nB_n \le4nF_\beta(h)+4B_n\widetilde{I}_{d,\beta}-2nF_{\beta/4}(h)+ 4CA_n B_n. 
\end{align*}
 Once again, the same proof yields an identical bound for $\sigma_\beta(H_{n,M})$ and concludes the proof. 
\end{proof}

\noindent We are now ready to state and prove the main theorem of this section:

\begin{theorem}[Free energy perturbation bound]\label{thmfreenregyblabla}
Fix $d\in\{2,3\}$, $\beta>0$. There exists a constant $C_\beta<\infty$, independent of $n,M\ge 1$, such that
\[
|F_\beta(H_n)-F_{\beta}(H_{n,M})|
\le C_\beta\,\left(n^{3}\alpha^{\max}_n + n^5 \left(\alpha^{\max}_n\right)^3\right)M^{-\frac{1}{4d}},
\]
where we defined $\alpha^{\max}_{n} :=\max_{1\le i<j\le n} |\alpha_{n,i,j}|.$
\end{theorem}

\begin{proof}
From \eqref{almostthereyes}, Lemma \ref{lem:uniform-H0-bound-abstract-detailed} and the fact that $\lambda_M=\Omega(M^{1/d}),$ we see 
\begin{align*}
  |F_\beta(H_n)-F_{\beta}(H_{n,M})|
\le C_\beta \left(nB_n +(A_n+1)B^2_n \right)M^{-\frac{1}{4d}}.  
\end{align*}
Using that \(B_n\le n^{2} \alpha^{\max}_{n}\) and $A_n\le n \alpha^{\max}_{n}$ finishes the proof. 
\end{proof}


\subsection{Relative entropy and trace-distance convergence}
\label{sec:EntropyTraceDistanceConv}

\noindent In this section, we use the relative entropy to obtain a trace-distance bound on the distance between the two full and approximated Gibbs state. 
Given two density operators \(\rho,\sigma\) on a Hilbert space \(\mathcal H\), the (quantum) relative entropy of \(\rho\) with respect to \(\sigma\) is defined by
\[
D(\rho\|\sigma):=
\begin{cases}
\Tr\bigl(\rho(\log\rho-\log\sigma)\bigr), & \text{if }\mathrm{supp}(\rho)\subset \mathrm{supp}(\sigma),\\
+\infty, & \text{otherwise}.
\end{cases}
\]
It is nonnegative and vanishes if and only if \(\rho=\sigma\). Moreover, it controls the trace norm via Pinsker's inequality,
\[
\|\rho-\sigma\|_1^2 \le 2 D(\rho\|\sigma).
\]
\begin{lemma}[Relative entropy bound for Gibbs states under relative form perturbations]
\label{lem:relative-entropy-relative-form-convergence}
Let \(H_{0}\) be self-adjoint on \(\mathcal H\), bounded from below with compact resolvent. Let \(W,W'\) be symmetric quadratic forms on \(Q(H_0)\) with relative form bound strictly smaller than \(1\), and define \(H=H_0+W\), \(H'=H_0+W'\) such that the respective Gibbs states \(\sigma_\beta(H),\sigma_\beta(H')\) exist. Set \(\Delta:=H-H'\) and assume that for some \(\varepsilon,c\ge0\),
\begin{equation}
\label{eq:relative-form-smallness-delta-entropy}
|\Delta[\psi]|\le \varepsilon\langle\psi,H_0\psi\rangle+c\|\psi\|^2\qquad\text{for all }\psi\in Q(H_0).
\end{equation}
Then the Gibbs states \(\sigma_\beta(H),\sigma_\beta(H')\) satisfy the relative entropy bounds
\begin{equation}
D(\sigma_\beta(H)\|\sigma_\beta(H')),\,D(\sigma_\beta(H')\|\sigma_\beta(H))
\le \beta\Bigl(\varepsilon \Tr(\sigma_\beta(H)H_0)+\varepsilon \Tr(\sigma_\beta(H')H_0)+2c\Bigr),
\end{equation}
and consequently
\begin{equation}
\|\sigma_\beta(H)-\sigma_\beta(H')\|_1
\le \sqrt{2\beta\Bigl(\varepsilon \Tr(\sigma_\beta(H)H_0)+\varepsilon \Tr(\sigma_\beta(H')H_0)+2c\Bigr)}.
\end{equation}
\end{lemma}

\begin{proof}
Self-adjointness, lower bounds, and compact resolvent follow from the form assumptions; hence, the Gibbs operators are trace class. Using \(\log\sigma_\beta(H)=-\beta H-\log\mathcal Z_\beta(H)\), we compute
\[\begin{split}
D(\sigma_\beta(H)\|\sigma_\beta(H'))&=\beta\Tr(\sigma_\beta(H)(H'-H))+\log\mathcal Z_\beta(H')-\log\mathcal Z_\beta(H)\\
&=\beta\bigl(F_\beta(H)-F_\beta(H')-\Tr(\sigma_\beta(H)\Delta)\bigr),
\end{split}
\]
and similarly
\[
D(\sigma_\beta(H')\|\sigma_\beta(H))=\beta\bigl(\Tr(\sigma_\beta(H')\Delta)-(F_\beta(H)-F_\beta(H'))\bigr).
\]
By the Bogoliubov inequality as in Lemma~\ref{lem:free-energy-relative-form-convergence},
\[
\Tr(\sigma_\beta(H)\Delta)\le F_\beta(H)-F_\beta(H')\le \Tr(\sigma_\beta(H')\Delta),
\]
hence both relative entropies are bounded by \(\beta(\Tr(\sigma_\beta(H')\Delta)-\Tr(\sigma_\beta(H)\Delta))\). Using the spectral decomposition \(\sigma_\beta(H)=\sum_k p_k|\psi_k\rangle\langle\psi_k|\) and using \eqref{eq:relative-form-smallness-delta-entropy},
\[
|\Tr(\sigma_\beta(H)\Delta)|=\Bigl|\sum_k p_k\Delta[\psi_k]\Bigr|\le \varepsilon\Tr(\sigma_\beta(H)H_0)+c,
\]
and analogously for \(\sigma_\beta(H')\). Therefore
\[
\Tr(\sigma_\beta(H')\Delta)-\Tr(\sigma_\beta(H)\Delta)\le \varepsilon\Tr(\sigma_\beta(H)H_0)+\varepsilon\Tr(\sigma_\beta(H')H_0)+2c,
\]
which yields the entropy bounds. The trace norm estimate follows from Pinsker \(\|\rho-\sigma\|_1^2\le2D(\rho\|\sigma)\).
\end{proof}
\noindent Specifying this to the case of Coulomb interactions yields then:
\begin{theorem}[Trace norm perturbation bound]
\label{thm:trace-norm-perturbation-bound}
Fix \(d\in\{2,3\}\), \(\beta>0\). Then there exists \(C_\beta<\infty\), independent of \(n,M\), such that
\begin{equation}
\|\sigma_\beta(H_n)-\sigma_\beta(H_{n,M})\|_1 \le C_\beta \sqrt{n^3\alpha^{\max}_n + n^5(\alpha^{\max}_n)^3}\ M^{-\frac{1}{8d}},
\end{equation}
where we defined $\alpha^{\max}_{n} :=\max_{1\le i<j\le n} |\alpha_{n,i,j}|.$
\end{theorem}

\begin{proof}
By Lemma~\ref{lem:product-cutoff-riesz-no-scaling}, we have
\[
|\Delta_{n,M}[\psi]|\lesssim B_n(\lambda_{M+1}+1)^{-1/4}\langle\psi,(H_{0,n}+1)\psi\rangle,
\]
hence there exists \(\varepsilon_{n,M},c_{n,M}>0\) such that
\[
|\Delta_{n,M}[\psi]|\le \varepsilon_{n,M}\langle\psi,H_{0,n}\psi\rangle + c_{n,M}\|\psi\|^2,
\quad
\varepsilon_{n,M}\sim c_{n,M}\sim B_n(\lambda_{M+1}+1)^{-1/4}.
\]
Applying Lemma~\ref{lem:relative-entropy-relative-form-convergence} with \(H=H_n\), \(H'=H_{n,M}\), \(H_0=H_{0,n}\), we obtain
\[
\|\sigma_\beta(H_n)-\sigma_\beta(H_{n,M})\|_1 \le \sqrt{2\beta\Bigl(\varepsilon_{n,M}\Tr(\sigma_\beta(H_n)H_{0,n})+\varepsilon_{n,M}\Tr(\sigma_\beta(H_{n,M})H_{0,n})+2c_{n,M}\Bigr)}.
\]
By Lemma~\ref{lem:uniform-H0-bound-abstract-detailed}, and the fact that $B_n \le n^2\alpha^{\max}_n$ and $A_n\le n \alpha^{\max}_n$ we see 
$$\Tr(\sigma_\beta(H_n)H_{0,n})\, , \,\Tr(\sigma_\beta(H_{n,M})H_{0,n}) \le C_\beta (n + n^3 (\alpha^{\max}_n)^2)\quad$$  and hence
\[
\|\sigma_\beta(H_n)-\sigma_\beta(H_{n,M})\|_1 \lesssim_\beta \Bigl((n^3\alpha^{\max}_n + n^5(\alpha^{\max}_n)^3)(\lambda_{M+1}+1)^{-1/4}\Bigr)^{1/2}.
\]
Using \(\lambda_{M+1}\gtrsim M^{1/d}\), we conclude
\[
\|\sigma_\beta(H_n)-\sigma_\beta(H_{n,M})\|_1 \le C_\beta \sqrt{n^3\alpha^{\max}_n + n^5(\alpha^{\max}_n)^3}\ M^{-\frac{1}{8d}}.
\]
\end{proof}


\section{Spectral gap analysis}
\label{section:Pert_Theory}


\noindent In this section, we study the spectral gap of the generators $L_{H_{n,M}}$. The proof of Theorem \ref{mainthmpositivegap} is based on a perturbative analysis of the reference generator $L_{H_{0,n}}$.

\subsection{Quadratic Hamiltonians}



\noindent In order to study the generator of the dynamics corresponding to $H_{0,n}$, we consider the quantum Ornstein–Uhlenbeck semigroup converging to the Gibbs state of the number operator $N:=a^\dagger a\equiv \frac{1}{2}(-\frac{d^2}{dx^2}+x^2-1)$ on $L^2(\mathbb{R})$, with associated ladder operators $a$ and $a^\dagger$, whose generator on $\mathscr{T}_1(L^2(\mathbb{R}))$ is given by \cite{OU}
\begin{align}
\label{eq:qOUgen}
\mathcal{L}_{N}(\rho)
=&|\widehat{f}(1)|^2\,\Big(a^\dagger \rho a-\frac{1}{2}\{aa^\dagger,\rho\}\Big)
+|\widehat{f}(-1)|^2\,\Big(a\rho a^\dagger-\frac{1}{2}\{a^\dagger a,\rho\}\Big).
\end{align}

\noindent This generator coincides with $\cL_{(H_{0,n=1}-1)/2}$ in the case of a single one-dimensional particle, i.e.~$d=1$. Denoting the birth rate $\nu_+:=|\widehat{f}(1)|^2$ and death rate $\nu_-:=|\widehat{f}(-1)|^2$, the associated self-adjoint generator on $\mathscr{T}_2(\cH)$ takes the form \cite{OU}:
\begin{align}
\label{eq:HSgenerator}
L_{N}(x)
=&
-\left( \frac{\nu_- + \nu_+}{2}(Nx + xN) + \nu_+ x \right)
+ \sqrt{\nu_+ \nu_-}( a x a^{\dagger} +  a^{\dagger} x a).
\end{align}
In later parts of the paper, we shall also use the lower bound
\begin{align}\label{ladderblockboundqOU}
L_{N} \;\ge\; L_{\operatorname{LB}},
\end{align}
in the sense of quadratic forms, as shown in \cite[Proof of Thm.~6.3]{OU}. Here $L_{\operatorname{LB}}$ is the so-called \textit{ladder-block} operator defined by
\begin{equation}
\label{eq:ladderblock}
L_{\operatorname{LB}}(x)
:= \frac{1}{2}(\nu_+ - \nu_-)^2 (N x + x N) \;-\; \nu_+(\nu_- - \nu_+)\, x.
\end{equation}
The operator $L_{\operatorname{LB}}$ admits the spectral decomposition
\begin{equation}
\label{eq:LBdecomp}
L_{\operatorname{LB}} = \sum_{k \ge 0} \kappa_k\, R_k\qquad \text{ where }
\kappa_k := \frac{1}{2}(\nu_+ - \nu_-)^2\, k \;-\; \nu_+(\nu_- - \nu_+),
\end{equation}
and $R_k$ is the orthogonal projection onto the subspace $\operatorname{span}\{\, |n\rangle\langle m| : n+m = k \,\}$, where $\ket{n}$ denotes the Fock state with $N\ket{n}=n\ket{n}$. This decomposition provides a simpler spectral structure than that of $L_{N}$, making it a useful comparison operator for later perturbative spectral analysis.





\medskip 

\noindent Next, to motivate our unperturbed Hamiltonian $H_{0,n}$ in the context of Coulomb interacting systems, we consider a model of $d=2$ dimensional particles in a constant magnetic field with a quadratic trapping potential
\begin{align}\label{H0unperturedH}
   H_0\equiv  h_i = (-i\nabla_{x_i} - \vec{A}(x_i))^2 + \langle x_i,Mx_i\rangle.
\end{align}
Furthermore, we set \( {A}(x_i) = \frac{1}{2}(-Bx_{i2},Bx_{i1})^{\mathsf T} \)  the vector potential for a constant magnetic field \(B\) and coordinates $x_i:=(x_{i1},x_{i2})$ in the orthogonal direction and assume that \( M \in \mathbb R^{2\times 2} \) is a positive-definite matrix. 



\noindent Next, we study the spectral gap of $L_{H_0}$ for the Hamiltonian $H_0$ introduced in \eqref{H0unperturedH}. First, the spectral decomposition $H_0$ follows the method introduced in \cite{lin2002anisotropic}: we note that $H_0 = X^{\dagger} \mathcal H X$, where $X = (x_1,D_{x_1},x_2,D_{x_2})$, where $D_{x_i}:=-i \frac{\partial}{\partial x_j}$ and 
\[ \mathcal H = \begin{pmatrix} M_{11} + \frac{B^2}{4} &0 & M_{12} & -\frac{B}{2}  \\ 0 & 1&  \frac{B}{2} & 0\\  M_{12} &  \frac{B}{2}&  M_{22}+ \frac{B^2}{4} & 0 \\ -\frac{B}{2}& 0& 0& 1 \end{pmatrix}.\]
We then notice that $[X_{\alpha},X_{\beta}]= -\Sigma_{\alpha \beta}$, where $\Sigma = \operatorname{diag}(\sigma_2,\sigma_2).$ Next, we compute 
\[\begin{split} [X,H_0]&=\sum_{k} X_{k} \sum_{i,j} X_i \mathcal H_{ij} X_j e_k- \sum_{i,j} X_i \mathcal H_{ij} X_jX_{k}e_k \\
&= \sum_{i,j,k} (X_i  \mathcal H_{ij}X_{k} X_j +[X_k,X_i] \mathcal H_{ij}X_j)  e_k 
- \sum_{i,j} X_i \mathcal H_{ij} X_jX_{k}e_k \\ 
&= \sum_{i,j,k} X_i  \mathcal H_{ij}[X_{k}, X_j]  e_k -\Sigma_{ki} \mathcal H_{ij}X_j  e_k =- \sum_{i,j,k} X_i  \mathcal H_{ij}\Sigma_{kj}  e_k + \Sigma_{ki} \mathcal H_{ij}X_j  e_k \\
&=- \sum_{i,j,k}   e_k\Sigma_{kj}\mathcal H_{ji}X_i   + e_k\Sigma_{ki} \mathcal H_{ij}X_j = -2\Sigma \mathcal H X.\end{split} \]
Thus, we have $[iH_0,X]=2\Omega X$, where 
\[ \Omega:=i\Sigma \mathcal H=
\begin{pmatrix} 0 & 1 & -\frac{B}{2}& 0\\
 -M_{11}-\frac{B^2}{4} & 0 & -M_{12}& -\frac{B}{2} \\
 \frac{B}{2} & 0 & 0 & 1\\
-M_{12} & \frac{B}{2} & -M_{22}-\frac{B^2}{4} & 0
\end{pmatrix},\]
the characteristic polynomial of $\Omega$ reads
\[ \det(\Omega-\lambda I)=\lambda^4 + \lambda^2 ( \Tr(M) +B^2) + \operatorname{det}(M).\]

\noindent Since the characteristic polynomial is invariant under flips $\lambda \to -\lambda$, the eigenvalues of $\Omega$ are of the form 
\[ \operatorname{Sp}(\Omega) = \operatorname{diag}(\pm i \sigma_1 ,\pm i \sigma_2)\]
with
\[\begin{split}\sigma_1 &= \sqrt{\frac{\Tr(M)+B^2 + \sqrt{(\Tr(M)+B^2)^2-4\det(M)}}{2}} \,~, \\
\sigma_2 &=\sqrt{\frac{\Tr(M)+B^2 - \sqrt{(\Tr(M)+B^2)^2-4\det(M)}}{2}}\end{split}\]
which are both real and strictly positive since $\Tr(M),\det(M)>0.$ We then have left and right eigenvectors 
\[ \Omega^{\top}u_i = -i\sigma_i u_i \text{ and }\Omega v_i = -i\sigma_i v_i .\]

\noindent Taking the complex conjugate of this equation and using the fact that all entries in $\Omega$ are real, we find 
\[ \Omega^{\top}\overline{u_i} = i\sigma_i \overline{u_i} \text{ and } \Omega \overline{v_i} = i\sigma_i \overline{v_i},\]
We find that 
\[ i\sigma_i \langle v_j,u_i \rangle = \langle v_j,\Omega^{\top}u_i\rangle = \langle \Omega v_j,u_i\rangle=-i\sigma_j \langle v_j,u_i \rangle.  \]
Thus, $\langle v_j,u_i\rangle =0$, and by the same computation $\langle \overline{v_j},u_i\rangle =\delta_{ij}.$
In the following, we write $\Sigma_y = \Sigma$ and introduce $\Sigma_{z}:=\operatorname{diag}(\sigma_3,\sigma_3).$ We thus find that using $\Sigma_y^2 = 1$ and $\Sigma_y^{-1} = \Sigma_y =- \Sigma_y^{\top}$
\[\Sigma_y v_i = \frac{i\Sigma_y\Omega v_i}{\sigma_i} =-\frac{\mathcal H v_i}{\sigma_i} = -\frac{\mathcal H \Sigma_y (\Sigma_y v_i)}{\sigma_i} = -\frac{i\Omega^{\top} (\Sigma_y v_i)}{\sigma_i}.  \]
This shows that $\Sigma_y v_i$ is also an eigenvector with eigenvalue $i \sigma_i$ of $\Omega^{\top}$, and we conclude that $\overline{u_i} \propto \Sigma_y v_i.$ Thus, we may choose $\overline{u_i}:=-\Sigma_y v_i$, and therefore $\Sigma_y \overline{u_i} = -\Sigma_y^2 v_i = -v_i$, and normalize $u_i$ such that $\langle u_i,\overline{v_i} \rangle =-u_i^{\top}\Sigma_y \overline{u_i}=1.$
Defining $Q=(v_1,\overline{v_1},v_2,\overline{v_2})$ and $Q^{-1}= (u_1,\overline{u_1},u_2,\overline{u_2})^{\top}$, we find 
\[ Q^{-1}\Omega Q = \operatorname{diag}(-i\sigma_1,i\sigma_1,-i\sigma_2,i\sigma_2).\]
We also observe that
\[\begin{split} Q^{\dagger} &= (\overline{v_1}, v_1,\overline{v_2},v_2)^{\top} = (\Sigma_y \overline{u_1},-\Sigma_y u_1,\Sigma_y \overline{u_2},-\Sigma_y u_2)^{\top} =( -\overline{u_1}, u_1, -\overline{u_2},\Sigma_y u_2)^{\top}\Sigma_y= -\Sigma_z Q^{-1} \Sigma_y  .
\end{split} \]
We then define $a_i :=u_i^{\top} X$ and $a_i^{\dagger}=  u_i^{\dagger} X $ and observe that \[\begin{split} [a_i,a_j^{\dagger}]&=\sum_{n,m} (u_i)_n (\overline{u_j})_m [X_n,X_m] =- \sum_{n,m} (u_i)_n \Sigma_{nm}(\overline{u_j})_m =\delta_{ij} 
\end{split}\] and
\[\begin{split}  [a_i,a_j]&=\sum_{n,m} (u_i)_n (u_j)_m [X_n,X_m]=-\langle \Sigma_y\overline{u_i}, u_j \rangle =\langle v_i, u_j \rangle =0.  \end{split}\] 
Defining the operator-valued vector $A=(a_1,a_1^{\dagger},a_2,a_2^{\dagger})^{\top}$, we have, by construction, $A= Q^{-1}X$, and thus $X=QA.$ We may thus write the Hamiltonian as: 
\[ \begin{split} H_0 
&= X^{\dagger} \mathcal H X = A^{\dagger} Q^{\dagger} \mathcal H QA =-A^{\dagger}\Sigma_z Q^{-1}\Sigma_y \mathcal H QA = i A^{\dagger} \Sigma_z Q^{-1} \Omega Q A \\
& =i A^{\dagger} \Sigma_z \operatorname{diag}(-i\sigma_1,i\sigma_1,-i\sigma_2,i\sigma_2) A = A^{\dagger} \Sigma_z \operatorname{diag}(\sigma_1,-\sigma_1,\sigma_2,-\sigma_2)A \\
& = A^{\dagger} \operatorname{diag}(\sigma_1,\sigma_1,\sigma_2,\sigma_2)A = 2\sigma_1 (a_1^{\dagger}a_1+1) + 2\sigma_2 (a_2^{\dagger}a_2+1).\end{split} \]
Thus, we have found
\[ \operatorname{Sp}(H_0) = \sigma_1(2\mathbb N_0+1)+\sigma_2(2\mathbb N_0+1).\]
We thus see that 
\[H_{0,n} =\sum_{j\in[n]} h_j= \sum_{j\in[n]} \sum_{k\in[2]} h_{jk}\] with $h_{jk}=2 \sigma_k (a_{jk}^{\dagger}a_{jk}+1)$ satisfying 
$[h_{jk},h_{j'k'}]=\delta_{k,k'}\delta_{j,j'}$ where $j$ labels the particle. Hence, the generator $L_{H_{0,n}}$ defined by $H_{0,n}$, the Hamiltonian of the non-interacting particle, is a quantum Ornstein-Uhlenbeck process \cite{OU}. Combining with the spectral analysis carried in \cite{OU}, we find
\begin{theorem}
\label{thm:DFT2}
   The associated generator of the Gibbs sampler in $\mathscr{T}_2(L^2(\mathbb{R}^{2n}))$ is a family of quantum Ornstein-Uhlenbeck generators with coefficients $\nu_{-,k}:=\vert  \widehat{f}(-2\sigma_k)\vert^2 $, $\nu_{+,k}:=\vert  \widehat{f}(2\sigma_k)\vert^2$ and
with spectrum 
\[ \operatorname{Sp}(L_{H_{0,n}}) =- \sum_{k,j}\left\{ \ell\left(\frac{\nu_{-,k}-\nu_{+,k}}{2}\right) ; \ell \in \mathbb N_0\right\}. \]
\end{theorem}

\noindent From here onwards, we assume for simplicity that the unperturbed Hamiltonian $H_0$ is already in diagonal form, i.e. $H_{0,n}:=\sum_{j\in [n]}\sum_{k\in [d]}h_{jk}$ with $\sigma_k=1$ for all $k$.















% \subsection{Perturbative gap analysis}



































\subsection{Finite rank perturbative stability}

 
\noindent In order to establish the spectral gap of $L_{H_{n,M}}$, we now show that finite-rank perturbations of the $H_{0,n}$ only produce finite-rank corrections in the structure of the $\mathscr{T}_2$ generator. In the next section, we will show that this is sufficient to preserve discreteness of the spectrum and hence a spectral gap. Our first Lemma below shows that perturbing the Hamiltonian in a finite-dimensional sector does not affect the jump operators on the infinite-dimensional bulk, and only creates a finite-rank correction.



\begin{lemma}[Finite-rank stability under exponential conjugation]
\label{lem:conjugated-L-finite-rank-remainder}
Let \(H_0\) be a self-adjoint operator on a Hilbert space \(\mathcal H\), bounded from below, with discrete spectrum. Let \(P\) be a finite-rank orthogonal projection such that $[P,H_0]=0,$
and set \(Q:=1-P\). Let \(R=PRP\) be bounded and self-adjoint, and define
\[
H:=H_0+R.
\]
Then \(H\) is reduced by the decomposition
\[
\mathcal H=Q\mathcal H\oplus P\mathcal H, \text{ such that }
H=QH_0Q\oplus (PH_0P+R).
\]
Let \(A^\alpha\) be a possibly unbounded operator such that
\[
P\mathcal H\subset D(A^\alpha)
\qquad\text{and}\qquad
PA^\alpha \text{ extends to a bounded operator on }\mathcal H.
\]
Fix \(s\in \mathbb R\). Assume that the operators
\[\begin{split}
L_s^\alpha(H)
&:=
\sum_{E,E'\in \spec(H)}
e^{s(E-E')}\widehat f(E-E')\,P_E(H)A^\alpha P_{E'}(H) \text{ and }\\
L_s^\alpha(H_0)
&:=
\sum_{E,E'\in \spec(H_0)}
e^{s(E-E')}\widehat f(E-E')\,P_E(H_0)A^\alpha P_{E'}(H_0)
\end{split}
\]
are well-defined, and denote $L^\alpha(H):=L^\alpha_0(H)$, resp.~$L^\alpha(H_0):=L^\alpha_0(H_0)$. Then
\[
L_s^\alpha(H)=L_s^\alpha(H_0)+\mathcal R_s^\alpha,
\qquad
Q\mathcal R_s^\alpha Q=0,
\]
and \(\mathcal R_s^\alpha\) is finite rank. In particular, if the conjugated operators
\[
e^{sH}L^\alpha(H)e^{-sH},
\qquad
e^{sH_0}L^\alpha(H_0)e^{-sH_0}
\]
are well-defined and coincide with the above series expansions, then with \(\mathcal R_s^\alpha\) finite rank
\[
e^{sH}L^\alpha(H)e^{-sH}
=
e^{sH_0}L^\alpha(H_0)e^{-sH_0}
+\mathcal R_s^\alpha,
\qquad
Q\mathcal R_s^\alpha Q=0.
\]

\end{lemma}

\begin{proof}
Since \([P,H_0]=0\), the subspaces \(P\mathcal H\) and \(Q\mathcal H\) reduce \(H_0\). Since \(R=PRP\), one has $RQ=QR=0.$
Hence \(H=H_0+R\) is block diagonal with respect to the orthogonal decomposition
$\mathcal H=Q\mathcal H\oplus P\mathcal H,$
namely
\[
H=QH_0Q\oplus (PH_0P+R).
\]
We now describe the spectral projections of \(H_0\) and \(H\) relative to this decomposition. Since \(Q\mathcal H\) reduces \(H_0\), we may write
\[
QH_0Q=\sum_{\mu\in \spec(QH_0Q)} \mu\,Q_\mu,
\]
where the \(Q_\mu\) are pairwise orthogonal spectral projections satisfying \(Q_\mu\le Q\). On the finite-dimensional space \(P\mathcal H\), the operator \(PH_0P\) has a spectral decomposition
\[
PH_0P=\sum_{\ell=1}^r \nu_\ell \widetilde\Pi_\ell,
\]
where \(\widetilde\Pi_\ell\le P\), the projections \(\widetilde\Pi_\ell\) are pairwise orthogonal, and
\[
\sum_{\ell=1}^r \widetilde\Pi_\ell=P.
\]
Likewise, since \(PH_0P+R\) acts on the finite-dimensional space \(P\mathcal H\), it has a spectral decomposition
\[
PH_0P+R=\sum_{j=1}^m \lambda_j \Pi_j,
\]
with \(\Pi_j\le P\), pairwise orthogonal, and
\[
\sum_{j=1}^m \Pi_j=P.
\]
It follows that the spectral projections of \(H_0\) are precisely $\{Q_\mu\}_{\mu\in \spec(QH_0Q)}\cup \{\widetilde\Pi_\ell\}_{\ell=1}^r,$
while the spectral projections of \(H\) are precisely $\{Q_\mu\}_{\mu\in \spec(QH_0Q)}\cup \{\Pi_j\}_{j=1}^m.$
Indeed, on \(Q\mathcal H\), the operators \(H\) and \(H_0\) coincide, whereas on \(P\mathcal H\) the operator \(H_0\) is replaced by \(PH_0P+R\).
Now expand \(L_s^\alpha(H)\) according to this decomposition. Since the spectral projections of \(H\) are \(\{Q_\mu\}\cup\{\Pi_j\}\), we obtain
\[
\begin{aligned}
L_s^\alpha(H)
&=
\sum_{\mu,\nu} e^{s(\mu-\nu)}\widehat f(\mu-\nu)\,Q_\mu A^\alpha Q_\nu 
+\sum_{\mu,j} e^{s(\mu-\lambda_j)}\widehat f(\mu-\lambda_j)\,Q_\mu A^\alpha \Pi_j 
\\&\quad  +\sum_{j,\mu} e^{s(\lambda_j-\mu)}\widehat f(\lambda_j-\mu)\,\Pi_j A^\alpha Q_\mu \qquad +\sum_{j,k} e^{s(\lambda_j-\lambda_k)}\widehat f(\lambda_j-\lambda_k)\,\Pi_j A^\alpha \Pi_k .
\end{aligned}
\]
Similarly,
\[
\begin{aligned}
L_s^\alpha(H_0)
&=
\sum_{\mu,\nu} e^{s(\mu-\nu)}\widehat f(\mu-\nu)\,Q_\mu A^\alpha Q_\nu 
+\sum_{\mu,\ell} e^{s(\mu-\nu_\ell)}\widehat f(\mu-\nu_\ell)\,Q_\mu A^\alpha \widetilde\Pi_\ell \\
&\quad +\sum_{\ell,\mu} e^{s(\nu_\ell-\mu)}\widehat f(\nu_\ell-\mu)\,\widetilde\Pi_\ell A^\alpha Q_\mu
+\sum_{\ell,\ell'} e^{s(\nu_\ell-\nu_{\ell'})}\widehat f(\nu_\ell-\nu_{\ell'})\,\widetilde\Pi_\ell A^\alpha \widetilde\Pi_{\ell'} .
\end{aligned}
\]
Subtracting these two expressions yields
\[
L_s^\alpha(H)-L_s^\alpha(H_0)=\mathcal R_s^\alpha,
\]
where
\[
\begin{aligned}
\mathcal R_s^\alpha
&=
\sum_{\mu,j} e^{s(\mu-\lambda_j)}\widehat f(\mu-\lambda_j)\,Q_\mu A^\alpha \Pi_j
-\sum_{\mu,\ell} e^{s(\mu-\nu_\ell)}\widehat f(\mu-\nu_\ell)\,Q_\mu A^\alpha \widetilde\Pi_\ell \\
&\quad
+\sum_{j,\mu} e^{s(\lambda_j-\mu)}\widehat f(\lambda_j-\mu)\,\Pi_j A^\alpha Q_\mu
-\sum_{\ell,\mu} e^{s(\nu_\ell-\mu)}\widehat f(\nu_\ell-\mu)\,\widetilde\Pi_\ell A^\alpha Q_\mu \\
&\quad
+\sum_{j,k} e^{s(\lambda_j-\lambda_k)}\widehat f(\lambda_j-\lambda_k)\,\Pi_j A^\alpha \Pi_k
-\sum_{\ell,\ell'} e^{s(\nu_\ell-\nu_{\ell'})}\widehat f(\nu_\ell-\nu_{\ell'})\,\widetilde\Pi_\ell A^\alpha \widetilde\Pi_{\ell'}.
\end{aligned}
\]
The \(Q\)-\(Q\) block cancels completely, because \(H\) and \(H_0\) have the same spectral projections on \(Q\mathcal H\). Therefore
\[
Q\mathcal R_s^\alpha Q=0.
\]
It remains to show that \(\mathcal R_s^\alpha\) is finite rank. Since \(P\mathcal H\subset D(A^\alpha)\), the operator \(A^\alpha P\) is well-defined on all of \(\mathcal H\). Because \(P\mathcal H\) is finite-dimensional, the range of \(A^\alpha P\) is contained in the finite-dimensional space \(A^\alpha(P\mathcal H)\). Hence \(A^\alpha P\) is finite rank. Since \(\Pi_j,\widetilde\Pi_\ell\le P\), it follows that
\[
A^\alpha \Pi_j=(A^\alpha P)\Pi_j,
\qquad
A^\alpha \widetilde\Pi_\ell=(A^\alpha P)\widetilde\Pi_\ell
\]
are finite-rank operators. Consequently,
\[
Q_\mu A^\alpha \Pi_j,
\qquad
Q_\mu A^\alpha \widetilde\Pi_\ell
\]
are finite rank as well. Next, by assumption, the densely defined operator \(PA^\alpha\) extends to a bounded operator on \(\mathcal H\), which we again denote by \(PA^\alpha\). Since \(\Pi_j,\widetilde\Pi_\ell\le P\), we have on \(D(A^\alpha)\),
\[
\Pi_j A^\alpha=\Pi_j(PA^\alpha),
\qquad
\widetilde\Pi_\ell A^\alpha=\widetilde\Pi_\ell(PA^\alpha).
\]
Hence \(\Pi_j A^\alpha\) and \(\widetilde\Pi_\ell A^\alpha\) extend continuously to all of \(\mathcal H\). Their ranges are contained in
\[
\Pi_j\mathcal H\subset P\mathcal H,
\qquad
\widetilde\Pi_\ell\mathcal H\subset P\mathcal H,
\]
which are finite-dimensional. Thus \(\Pi_j A^\alpha\) and \(\widetilde\Pi_\ell A^\alpha\) are finite rank. Therefore
\[
\Pi_j A^\alpha Q_\mu,\qquad
\widetilde\Pi_\ell A^\alpha Q_\mu,\qquad
\Pi_j A^\alpha \Pi_k,\qquad
\widetilde\Pi_\ell A^\alpha \widetilde\Pi_{\ell'}
\]
are all finite rank. Since there are only finitely many indices \(j,k,\ell,\ell'\), the last two lines involve only finitely many finite-rank terms. For the mixed terms, note that
\[
\sum_j \Pi_j=P,
\qquad
\sum_\ell \widetilde\Pi_\ell=P.
\]
Hence
\[
\sum_{\mu,j} e^{s(\mu-\lambda_j)}\widehat f(\mu-\lambda_j)\,Q_\mu A^\alpha \Pi_j
=
\Big(\sum_{\mu} e^{s\mu} \widehat f(\mu-\lambda_1)Q_\mu,\ldots\Big)
\]
is more cleanly viewed term-wise as a sum of operators, each factoring through the finite-dimensional space \(P\mathcal H\). Indeed, for each fixed \(j\),
\[
\sum_\mu e^{s(\mu-\lambda_j)}\widehat f(\mu-\lambda_j)\,Q_\mu A^\alpha \Pi_j
=
\Big(\sum_\mu e^{s(\mu-\lambda_j)}\widehat f(\mu-\lambda_j)\,Q_\mu\Big)A^\alpha \Pi_j.
\]
By well-definedness of \(L_s^\alpha(H)\), this operator is well-defined, and since \(A^\alpha \Pi_j\) is finite rank, it is finite rank. Summing over the finitely many \(j\) shows that
\[
\sum_{\mu,j} e^{s(\mu-\lambda_j)}\widehat f(\mu-\lambda_j)\,Q_\mu A^\alpha \Pi_j
\]
is finite rank. The same argument applies to
\[
\sum_{\mu,\ell} e^{s(\mu-\nu_\ell)}\widehat f(\mu-\nu_\ell)\,Q_\mu A^\alpha \widetilde\Pi_\ell,
\]
and similarly, using that \(\Pi_j A^\alpha\) and \(\widetilde\Pi_\ell A^\alpha\) are finite rank, to the terms
\[
\sum_{j,\mu} e^{s(\lambda_j-\mu)}\widehat f(\lambda_j-\mu)\,\Pi_j A^\alpha Q_\mu,
\qquad
\sum_{\ell,\mu} e^{s(\nu_\ell-\mu)}\widehat f(\nu_\ell-\mu)\,\widetilde\Pi_\ell A^\alpha Q_\mu.
\]
Therefore, every group of terms in \(\mathcal R_s^\alpha\) is finite rank, and hence \(\mathcal R_s^\alpha\) itself is finite rank. Finally, suppose that the conjugated operators \(e^{sH}L^\alpha(H)e^{-sH}\) and \(e^{sH_0}L^\alpha(H_0)e^{-sH_0}\) are well-defined and agree with the above spectral expansions. Then, using
\[
e^{sH}P_E(H)=e^{sE}P_E(H),
\qquad
P_{E'}(H)e^{-sH}=e^{-sE'}P_{E'}(H),
\]
we obtain
\[
e^{sH}L^\alpha(H)e^{-sH}
=
\sum_{E,E'} e^{s(E-E')}\widehat f(E-E')\,P_E(H)A^\alpha P_{E'}(H)
=
L_s^\alpha(H),
\]
and similarly
\[
e^{sH_0}L^\alpha(H_0)e^{-sH_0}=L_s^\alpha(H_0).
\]
The stated identity, therefore, follows from the first part of the proof. 
\end{proof}


\noindent Next, we establish that the finite-rank stability of Lemma \ref{lem:conjugated-L-finite-rank-remainder} propagates to the quadratic expressions appearing in the generator $L_{H_{n,M}}$ (cf.~\eqref{LHnmgenerator}). This is crucial since the generator depends on terms such as \((L^\alpha)^\dagger L^\alpha\) and \(L^\alpha (L^\alpha)^\dagger\).



\begin{corollary}[Finite-rank stability of quadratic expressions]
\label{cor:quadratic-finite-rank}
Under the assumptions of Lemma~\ref{lem:conjugated-L-finite-rank-remainder}, set \(L:=L_s^\alpha(H)\), \(L_0:=L_s^\alpha(H_0)\), and \(\mathcal R:=\mathcal R_s^\alpha\), so that \(L=L_0+\mathcal R\) with \(\mathcal R\) finite rank. Assume there exists a dense subspace \(\mathcal D\subset\mathcal H\) with \(\mathcal D\subset D(L)\cap D(L_0)\), \(Lx=L_0x+\mathcal R x\) for all \(x\in\mathcal D\), and such that all expressions below are well-defined on \(\mathcal D\). If moreover \(\operatorname{Ran}(\mathcal R)\subset D(L^\dagger)\cap D(L_0^\dagger)\), then on \(\mathcal D\)
\[
L^\dagger L-L_0^\dagger L_0=L_0^\dagger \mathcal R+\mathcal R^\dagger L_0+\mathcal R^\dagger \mathcal R,
\]
and \(K:=L^\dagger L-L_0^\dagger L_0\) is finite rank on \(\mathcal D\). Likewise, if \(\operatorname{Ran}(\mathcal R^\dagger)\subset D(L)\cap D(L_0)\), then on every dense subspace \(\mathcal D'\subset D(L^\dagger)\cap D(L_0^\dagger)\) such that \(L^\dagger x=L_0^\dagger x+\mathcal R^\dagger x\) for all \(x\in\mathcal D'\) and all expressions are well-defined,
\[
LL^\dagger-L_0L_0^\dagger=L_0\mathcal R^\dagger+\mathcal R\mathcal R^\dagger+\mathcal R L_0^\dagger,
\]
and \(\widetilde K:=LL^\dagger-L_0L_0^\dagger\) is finite rank on \(\mathcal D'\).
\end{corollary}


\begin{proof}
Since \(\mathcal R\) is finite rank, it is bounded, hence \(\mathcal R^\dagger\) is bounded. For \(x\in\mathcal D\) one has \(Lx=L_0x+\mathcal R x\), thus \(L^\dagger Lx=L^\dagger L_0x+L^\dagger \mathcal R x\), and subtracting \(L_0^\dagger L_0x\) gives \((L^\dagger L-L_0^\dagger L_0)x=(L^\dagger-L_0^\dagger)L_0x+L^\dagger \mathcal R x\). On \(D(L^\dagger)\cap D(L_0^\dagger)\) one has \(L^\dagger=L_0^\dagger+\mathcal R^\dagger\), and since \(\operatorname{Ran}(\mathcal R)\subset D(L^\dagger)\cap D(L_0^\dagger)\), this yields \(L^\dagger L_0x=L_0^\dagger L_0x+\mathcal R^\dagger L_0x\) and \(L^\dagger \mathcal R x=L_0^\dagger \mathcal R x+\mathcal R^\dagger \mathcal R x\), hence \(L^\dagger L-L_0^\dagger L_0=L_0^\dagger \mathcal R+\mathcal R^\dagger L_0+\mathcal R^\dagger \mathcal R\) on \(\mathcal D\). Finite-rankness follows since \(\operatorname{Ran}(\mathcal R)\) is finite-dimensional and contained in \(D(L_0^\dagger)\), so \(L_0^\dagger \mathcal R\) is finite rank; \(\mathcal R^\dagger L_0\) is finite rank; and \(\mathcal R^\dagger \mathcal R\) is finite rank.

For the second identity, let \(x\in\mathcal D'\). Then \(L^\dagger x=L_0^\dagger x+\mathcal R^\dagger x\), hence \(LL^\dagger x=LL_0^\dagger x+L\mathcal R^\dagger x\), and subtracting \(L_0L_0^\dagger x\) gives \((LL^\dagger-L_0L_0^\dagger)x=(L-L_0)L_0^\dagger x+L\mathcal R^\dagger x=\mathcal R L_0^\dagger x+L\mathcal R^\dagger x\). Since \(\operatorname{Ran}(\mathcal R^\dagger)\subset D(L)\cap D(L_0)\), one has \(L\mathcal R^\dagger x=L_0\mathcal R^\dagger x+\mathcal R\mathcal R^\dagger x\), hence \(LL^\dagger-L_0L_0^\dagger=L_0\mathcal R^\dagger+\mathcal R\mathcal R^\dagger+\mathcal R L_0^\dagger\) on \(\mathcal D'\). Again, each term is of finite rank, so \(\widetilde K\) is of finite rank.
\end{proof}




\noindent We can now combine these two ingredients to study the full perturbed generator $L_{H_{n,M}}$: all modifications remain confined to finite-rank coefficients. We recall that conjugation by $e^{\pm\beta H/4}$ leads to the generator on $\mathscr{T}_2(\cH)$: on $\mathscr F:=\operatorname{span}\{|E\rangle\langle E'|:E,E'\in\operatorname{Sp}(H)\}$, denoting $\Gamma_\tau(X):=e^{\tau H}Xe^{-\tau H}$ and a finite set of bare jumps $\{A^\alpha\}_{\alpha\in\mathcal{A}}$,
\begin{align*}
L_{H}
=
-i(B_{+,H}\bullet -\bullet B_{-,H})
+\sum_{\alpha\in\mathcal A}
\Big(L^\alpha_{+,H}\,\bullet\,(L^\alpha_{-,H})^\dagger-\tfrac12 K^\alpha_{+,H}\bullet -\tfrac12 \bullet K^\alpha_{-,H}\Big),
\end{align*}
with 
\begin{align}
B_H&:=\frac{i}{2}\sum_{\alpha\in\mathcal A}\sum_{E,F,G\in\operatorname{Sp}(H)}
\tanh\!\Big(\frac{\beta(F-E)}{4}\Big)\,
\overline{\widehat f(G-F)}\,\widehat f(G-E)\,
P_{F}(A^\alpha)^\dagger P_GA^\alpha P_E\nonumber \\
&=\frac{i}{2}\sum_{\alpha\in\mathcal A}\sum_{E,F\in\operatorname{Sp}(H)}
\tanh\!\Big(\frac{\beta(F-E)}{4}\Big)\,\,
P_{F}L^\alpha(H)^\dagger L^\alpha(H) P_E\label{BHexpressionLa}
\end{align}
and
\begin{align}
L^\alpha_{\pm,H}:=\Gamma_{\pm\beta/4}(L^\alpha(H)),
\quad
K^\alpha_{\pm,H}:=\Gamma_{\pm\beta/4}\big((L^\alpha(H))^\dagger L^\alpha(H)\big),\quad B_{\pm,H}:=\Gamma_{\pm\beta/4}(B_H).\label{pmallneededops}
\end{align} 
Analogously, we denote the same operators associated with $H_0$ by $L_{\pm,0}^\alpha$ and $B_{\pm,0}$. Invoking Lemma \ref{lem:conjugated-L-finite-rank-remainder}, we further denote $L_{\pm,H}^\alpha=L_{\pm,0}^\alpha+\mathcal R_\pm^\alpha$ with \(\mathcal R_\pm^\alpha\) having finite rank.




\begin{proposition}[Finite-rank perturbation of the conjugated generator]
\label{prop:conjugated-generator-perturbation}
  With the notations of the above paragraph,
\[
L_{\operatorname{pert}}=L_{H}-L_{H_0}=-i\big((B_{+,H}-B_{+,0})\bullet -\bullet (B_{-,H}-B_{-,0})\big)+\sum_{\alpha\in\mathcal A}\big(\mathfrak G_\alpha -\tfrac12\mathfrak A_\alpha -\tfrac12\mathfrak B_\alpha\big),
\]
where
\[
\mathfrak G_\alpha=\mathcal R_+^\alpha \bullet  (L_{+,0}^\alpha)^\dagger+L_{+,0}^\alpha \bullet (\mathcal R_+^\alpha)^\dagger+\mathcal R_+^\alpha \bullet  (\mathcal R_+^\alpha)^\dagger,
\]
\[
\mathfrak A_\alpha=\big((L_{-,H}^\alpha)^\dagger L_{+,H}^\alpha-(L_{-,0}^\alpha)^\dagger L_{+,0}^\alpha\big)\bullet \qquad \text{ and }\qquad
\mathfrak B_\alpha(x)=\bullet \big(L_{-,H}^\alpha (L_{+,H}^\alpha)^\dagger-L_{-,0}^\alpha (L_{+,0}^\alpha)^\dagger\big).
\]
Moreover, on any common domain where all products are defined,
\begin{equation}
    \begin{split}
    \label{eq:anti-comm_terms}
(L_{-,H}^\alpha)^\dagger L_{+,H}^\alpha-(L_{-,0}^\alpha)^\dagger L_{+,0}^\alpha&=(L_{-,0}^\alpha)^\dagger \mathcal R_+^\alpha+(\mathcal R_-^\alpha)^\dagger L_{+,0}^\alpha+(\mathcal R_-^\alpha)^\dagger \mathcal R_+^\alpha,\\
L_{-,H}^\alpha (L_{+,H}^\alpha)^\dagger-L_{-,0}^\alpha (L_{+,0}^\alpha)^\dagger&=L_{-,0}^\alpha (\mathcal R_+^\alpha)^\dagger+\mathcal R_-^\alpha (L_{+,0}^\alpha)^\dagger+\mathcal R_-^\alpha (\mathcal R_+^\alpha)^\dagger.
    \end{split}
\end{equation}
In particular, all coefficient differences are finite rank under the same range assumptions as in Corollary~\ref{cor:quadratic-finite-rank}.
\end{proposition}

\begin{proof}
Insert \(\rho=e^{-\beta H/4}xe^{-\beta H/4}\) into \(\mathcal L_{\widehat f,H}\) and conjugate by \(e^{\beta H/4}\). The commutator gives \(e^{\beta H/4}[B_H,\rho]e^{\beta H/4}=B_{+,H}x-xB_{-,H}\). The gain term becomes \(e^{\beta H/4}L^\alpha(H)\rho (L^\alpha(H))^\dagger e^{\beta H/4}=L_{+,H}^\alpha x (L_{+,H}^\alpha)^\dagger\). For the anticommutator, note \((L_{-,H}^\alpha)^\dagger=e^{\beta H/4}(L^\alpha(H))^\dagger e^{-\beta H/4}\), hence
\begin{align*}
&(L_{-,H}^\alpha)^\dagger L_{+,H}^\alpha=e^{\beta H/4}(L^\alpha(H))^\dagger L^\alpha(H)e^{-\beta H/4}\\
&  L_{-,H}^\alpha (L_{+,H}^\alpha)^\dagger=e^{-\beta H/4}L^\alpha(H)(L^\alpha(H))^\dagger e^{\beta H/4}
\end{align*}
which yields the displayed form. The same identities hold for \(H_0\), and subtraction gives the difference formula. Using \(L_{+,H}^\alpha=L_{+,0}^\alpha+\mathcal R_+^\alpha\), expand \((L_{+,H}^\alpha)x(L_{+,H}^\alpha)^\dagger\) to obtain \(\mathfrak G_\alpha\). Similarly, combining \(L_{-,H}^\alpha=L_{-,0}^\alpha+\mathcal R_-^\alpha\) and \(L_{+,H}^\alpha=L_{+,0}^\alpha+\mathcal R_+^\alpha\) gives the expansions for \(\mathfrak A_\alpha\) and \(\mathfrak B_\alpha\). The explicit identities for the coefficient differences follow by direct expansion on a common domain where all products are defined. Since \(\mathcal R_\pm^\alpha\) are finite rank, all purely quadratic terms are finite rank, and the mixed terms are finite rank under the same domain and range assumptions as in Corollary~\ref{cor:quadratic-finite-rank}.
\end{proof}



\subsection{Spectral gap for finite-rank perturbation of quadratic Hamiltonians}
\label{sec:GapForFiniteRankPerturb}

\noindent Next, we make use of the structural result of Proposition \ref{prop:conjugated-generator-perturbation} in the special case of a Gaussian reference Hamiltonian, i.e.~$H_0\equiv H_{0,n}$, $H\equiv H_{n,M}=H_{0,n}+ W_{n,M}$, and the bare jumps are chosen as $A^\alpha:=a_{ij}^\sigma$, for $\alpha=(i,j,\sigma)\in[n]\times [d]\times \{\emptyset,\dagger\}$, to show the positivity of the spectral gap of $L_{H_{n,M}}$. The proof uses the combined fact that the unperturbed generator is explicit, and the perturbation inherits a very rigid structure.


\medskip

\noindent First, we can easily verify that the unperturbed filtered jumps introduced in \eqref{pmallneededops} for $H=H_{0,n}$ take the form, for any $\alpha=(i,k,\sigma)\in\mathcal{A}$,
\begin{align*}
L_{\pm,0}^{\alpha}=\left\{\begin{aligned}
&\widehat f(-2)\,e^{\mp\beta/4}a_{ik}&\sigma=\emptyset\\
&\widehat{f}(2)\,e^{\mp\beta/4}a^\dagger_{ik}&\sigma=\dagger
\end{aligned}\,\,\,\right. .
\end{align*}
Moreover, by the detailed balance condition \eqref{eq:KMSFilterFunction}, $\widehat f(-2)e^{-\beta/2}=\widehat f(2)e^{\beta/2}$, and thus, denoting $\nu_{-}:=\vert  \widehat{f}(-2)\vert^2 $, $\nu_{+}:=\vert  \widehat{f}(2)\vert^2$ and $n_{ik}=a_{ik}^\dagger a_{ik}$,
\begin{align}
\label{eq:HSgenerator-specialized}
L_{H_{0,n}}
=\sum_{i\in[n]}\sum_{k\in[d]}
-\Big(\frac{\nu_-+\nu_+}{2}(n_{ik} \bullet +\bullet n_{ik})+\nu_+\bullet \Big)
+\sqrt{\nu_+\nu_-}\big(a_{ik} \bullet  a_{ik}^\dagger+a_{ik}^\dagger \bullet a_{ik}\big).
\end{align}
Therefore, from the observation made in \eqref{ladderblockboundqOU}, we see that
\begin{align}
L_{H_{0,n}}\ge \sum_{j\in[n]}\sum_{k\in[d]} \frac{1}{2}(\nu_+ - \nu_-)^2 (n_{jk} \bullet + \bullet n_{jk}) \;-\; \nu_+(\nu_- - \nu_+)\, \bullet\equiv L_{\operatorname{LB},n}\,.
\end{align}
\noindent Next, we consider general non-positive perturbations of the Gaussian generator $L_{H_{0,n}}$ of the form
\[
L' = L_{H_{0,n}} + L_{\operatorname{pert}}\le 0,
\]
To prove the existence of a spectral gap for $L'$, we use the following lemma with the operator $L_{\mathrm{LB}}$ in \eqref{eq:ladderblock}, along with the many-particle operator $L_{\text{LB},n}=\oplus_{i\in[n]}\oplus_{k\in[d]} L_{\text{LB}}.$


\begin{lemma}
\label{lem:compact_resolv}
If $\Vert L_{\operatorname{pert}} (L_{\operatorname{LB},n}-\lambda )^{-1} \Vert <1$
for some $\lambda\in\mathbb{C}$, then $L'$ has a discrete spectrum and hence is gapped. In particular, it is enough to assume that $L_{\operatorname{pert}}$ is
 $L_{\mathrm{LB},n}$-bounded with relative bound $0$, i.e. $D(L_{\mathrm{LB},n})\subset D(L_{\mathrm{pert}})$
and for every $\varepsilon>0$ there exists $C_\varepsilon\ge 0$ such that
\[
\|L_{\mathrm{pert}}x\|_2
\le
\varepsilon \|L_{\mathrm{LB},n}x\|_2+C_\varepsilon \|x\|_2,
\qquad x\in D(L_{\mathrm{LB},n}).
\]



\end{lemma}

\begin{proof}
Since $L'\le 0$, it suffices to show that it has a compact resolvent. 
For this, we use that by \eqref{ladderblockboundqOU} $L_{H_{0,1}} \ge L_{\text{LB}}$ implies $L_{H_{0,n}}\ge L_{\text{LB},n}$. We then define $L_{\text{LB},n}':=L_{\text{LB},n} + L_{\text{pert}}$ and have the resolvent identity
\[
 (L_{\text{LB},n}'-\lambda)^{-1}
 =(L_{\text{LB},n}-\lambda )^{-1}
 (1 +L_{\operatorname{pert}} (L_{\text{LB},n}-\lambda )^{-1} )^{-1}
\]
if $\Vert L_{\operatorname{pert}} (L_{\operatorname{LB},n}-\lambda )^{-1} \Vert <1$ for some $\lambda.$ Thus $L_{\text{LB},n}'$ has compact resolvent as long as the resolvent of $L_{\operatorname{LB},n}$ is defined, since the resolvent of $L_{\operatorname{LB},n}$ is compact.
Below, we use the essential spectrum of a self-adjoint operator $(T,D(T))$, denoted $\operatorname{Sp}_{\text{ess}}(T)$, which is the set of real numbers $\lambda\in\mathbb{R}$ such that $T-\lambda I$ is not a Fredholm operator. We recall that an operator is Fredholm if it is closed with a finite dimensional kernel and cokernel. By the monotonicity of operators, we have
\[
\inf \operatorname{Sp}_{\text{ess}}(L')
\ge
\inf \operatorname{Sp}_{\text{ess}}(L'_{\operatorname{LB},n})=\infty.
\]
Thus, the spectrum of $L'$ is purely discrete. 






Next, to see the second condition, fix $\varepsilon>0$, and choose $C_\varepsilon\ge0$ such that
\[
\|L_{\mathrm{pert}}x\|
\le
\varepsilon \|L_{\mathrm{LB},n}x\|+C_\varepsilon \|x\|,
\qquad x\in D(L_{\mathrm{LB},n}).
\]
For $y\in \mathscr T_2(\mathcal H)$, let $x:=(L_{\mathrm{LB},n}-i\mu)^{-1}y.$
Since $L_{\mathrm{LB},n}$ is self-adjoint, we have $x\in D(L_{\mathrm{LB},n})$, and therefore
\[
\|L_{\mathrm{pert}}(L_{\mathrm{LB},n}-i\mu)^{-1}y\|
=
\|L_{\mathrm{pert}}x\|
\le
\varepsilon \|L_{\mathrm{LB},n}x\|+C_\varepsilon \|x\|.
\]
Thus
\[
\|L_{\mathrm{pert}}(L_{\mathrm{LB},n}-i\mu)^{-1}y\|
\le
\varepsilon \|L_{\mathrm{LB},n}(L_{\mathrm{LB},n}-i\mu)^{-1}y\|
+
C_\varepsilon \|(L_{\mathrm{LB},n}-i\mu)^{-1}y\|.
\]
Taking the supremum over $\|y\|=1$, we obtain
\[
\|L_{\mathrm{pert}}(L_{\mathrm{LB},n}-i\mu)^{-1}\|
\le
\varepsilon \|L_{\mathrm{LB},n}(L_{\mathrm{LB},n}-i\mu)^{-1}\|
+
C_\varepsilon \|(L_{\mathrm{LB},n}-i\mu)^{-1}\|.
\]
Since $L_{\mathrm{LB},n}$ is self-adjoint,
\[
\|(L_{\mathrm{LB},n}-i\mu)^{-1}\|\le \frac1{|\mu|},
\]
and
\[
\|L_{\mathrm{LB},n}(L_{\mathrm{LB},n}-i\mu)^{-1}\|
=
\sup_{\lambda\in\spec(L_{\mathrm{LB},n})}\frac{|\lambda|}{|\lambda-i\mu|}
\le 1.
\]
Hence
\[
\|L_{\mathrm{pert}}(L_{\mathrm{LB},n}-i\mu)^{-1}\|
\le
\varepsilon+\frac{C_\varepsilon}{|\mu|}.
\]
Choose $\varepsilon<1$. Then for $|\mu|$ sufficiently large we have
\[
\varepsilon+\frac{C_\varepsilon}{|\mu|}<1,
\]
which proves
\[
\|L_{\mathrm{pert}}(L_{\mathrm{LB},n}-i\mu)^{-1}\|<1.
\]
\end{proof}







% \noindent A more quantitative estimate of the gap can also be derived at the cost of imposing a threshold on the perturbation strength.
% \begin{lemma}
% \label{lem:OU_gap_perturb}
% Let $r:=\lambda_0/2$, where $\lambda_0 = \frac{\nu_--\nu_+}{2}$ is the gap of the generator $L_{N}$, and assume that 
% \[
%   \sup_{\vert \lambda \vert=r}\frac{\varepsilon(\lambda)}{1-\varepsilon(\lambda)}<\frac{\lambda_0}{2},
% \]
% with
% \[
% \varepsilon(\lambda):=\Vert L_{\operatorname{pert}} (L_{N}-\lambda )^{-1} \Vert,
% \]
% then the gap $\operatorname{gap}(L')$ of the full generator $L'$ satisfies
% \[
% \operatorname{gap}(L') \ge r.
% \]
% \end{lemma}

% \begin{proof}

% To compare spectral projections, we write

% \[
% (L'-\lambda)^{-1} -(L_{N}-\lambda )^{-1}  
% =
% (L_{N}-\lambda )^{-1}
% \left((1 +L_{\operatorname{pert}} (L_{N}-\lambda )^{-1} )^{-1}-1 \right).
% \]
% Using the Neumann series
% \[(L'-\lambda)^{-1} -(L_{N}-\lambda )^{-1}  
% =
% (L_{N}-\lambda )^{-1}
% \sum_{j=1}^{\infty}
% (-L_{\operatorname{pert}} (L_{N}-\lambda )^{-1})^j .
% \]
% Hence, recalling the Hausdorff distance between two sets
% \[
% d_{H}(A,B)
% :=
% \max\left\{
% \sup_{a\in A}\inf_{b\in B} |a-b|,
% \;
% \sup_{b\in B}\inf_{a\in A} |a-b|
% \right\},
% \]
% we have that 
% \[
%  \Vert (L'-\lambda)^{-1} -(L_{N}-\lambda )^{-1} \Vert
%  \le
%  d_H(\operatorname{Sp}(L_{N}),\{\lambda\})^{-1}
%  \frac{\varepsilon(\lambda)}{1-\varepsilon(\lambda)} .
% \]
% For a radius $r$, we have
% \[
%  \sup_{\vert \lambda \vert =r}
%  \Vert (L'-\lambda)^{-1} -(L_{N}-\lambda )^{-1} \Vert
%  \le
%  \sup_{\vert \lambda \vert =r}
%  d_H(\operatorname{Sp}(L_{N}),\{\lambda\})^{-1}
%  \frac{\varepsilon(\lambda)}{1-\varepsilon(\lambda)} .
% \]
% Thus, from \cite[(4.43)]{Kato}, the condition
% \[
%  \sup_{\vert \lambda \vert =r}
%  d_H(\operatorname{Sp}(L_{N}),\{\lambda\})^{-1}
%  \frac{\varepsilon(\lambda)}{1-\varepsilon(\lambda)}
%  <1
% \]
% implies
% \[
% \operatorname{rank}(1_{[0,r]}(L'))
% =
% \operatorname{rank}(1_{[0,r]}(L_{N})).
% \]
% We may choose $r:=\lambda_0/2$, where $\lambda_0$ is the gap of the OU generator, such that
% \[
% \operatorname{rank}(1_{[0,r]}(L_{N}))=1.
% \]
% Hence
% \[
%   \sup_{\vert \lambda \vert=r}\frac{\varepsilon(\lambda)}{1-\varepsilon(\lambda)}<\frac{\lambda_0}{2}
%   \quad\Longrightarrow\quad
%   \operatorname{gap}(L')>\frac{\lambda_0}{2}.
% \]
% \end{proof}


\noindent Next, we specialize our discussion to the case of $L'=L_{H_{n,M}}$. The following constitutes the main result of this section:

\begin{theorem}  
\label{corr:Gaussian_gap}
 The generator $L_{H_{n,M}}$ exhibits discrete spectrum and thus $\operatorname{gap}(L_{H_{n,M}})>0$.
\end{theorem}


\begin{proof}
We start by specializing the structural result of Proposition \ref{prop:conjugated-generator-perturbation} to the present Gaussian unperturbed setting: we recall that \(H_0\equiv H_{0,n}=nd+2\sum_{i\in[n]}\sum_{k\in[d]} n_{ik}\) and \(H\equiv H_{n,M}=H_{0,n}+W_{n,M}\), where \(W_{n,M}=\Pi_M W_n\Pi_M\) is bounded and self-adjoint for $\Pi_M:=P_M^{\otimes n}$ with \([\Pi_M,H_{0,n}]=0\). 
Moreover,
\[
L_{H_{n,M}}-L_{H_{0,n}}=-i\big((B_{+,H_{n,M}}-B_{+,0})\bullet -\bullet (B_{-,H_{n,M}}-B_{-,0})\big)+\sum_{\alpha\in\mathcal A}\big(\mathfrak G_\alpha -\tfrac12\mathfrak A_\alpha -\tfrac12\mathfrak B_\alpha\big),
\]
where
\begin{align*}
&\mathfrak G_\alpha=\mathcal R_+^\alpha \bullet  (L_{+,0}^\alpha)^\dagger+L_{+,0}^\alpha \bullet (\mathcal R_+^\alpha)^\dagger+\mathcal R_+^\alpha \bullet  (\mathcal R_+^\alpha)^\dagger\equiv \mathfrak{G}_\alpha^{(1)}+\mathfrak{G}_\alpha^{(2)}+\mathfrak{G}_\alpha^{(3)}\\
&\mathfrak A_\alpha=\big((L_{-,H_{n,M}}^\alpha)^\dagger L_{+,H_{n,M}}^\alpha-(L_{-,0}^\alpha)^\dagger L_{+,0}^\alpha\big)\bullet \\
&\mathfrak B_\alpha=\bullet \big(L_{-,H_{n,M}}^\alpha (L_{+,H_{n,M}}^\alpha)^\dagger-L_{-,0}^\alpha (L_{+,0}^\alpha)^\dagger\big).
\end{align*}
Moreover, on any common domain where all products are defined,
\[
A_\alpha:=(L_{-,H_{n,M}}^\alpha)^\dagger L_{+,H_{n,M}}^\alpha-(L_{-,0}^\alpha)^\dagger L_{+,0}^\alpha=(L_{-,0}^\alpha)^\dagger \mathcal R_+^\alpha+(\mathcal R_-^\alpha)^\dagger L_{+,0}^\alpha+(\mathcal R_-^\alpha)^\dagger \mathcal R_+^\alpha,
\]
\[
B_\alpha:=L_{-,H_{n,M}}^\alpha (L_{+,H_{n,M}}^\alpha)^\dagger-L_{-,0}^\alpha (L_{+,0}^\alpha)^\dagger=L_{-,0}^\alpha (\mathcal R_+^\alpha)^\dagger+\mathcal R_-^\alpha (L_{+,0}^\alpha)^\dagger+\mathcal R_-^\alpha (\mathcal R_+^\alpha)^\dagger.
\]
Next, we verify the relative boundedness condition of Lemma \ref{lem:compact_resolv}. We first argue that for each $\alpha$, $\mathfrak{A}_\alpha$ and $\mathfrak{B}_\alpha$ are both relatively bounded with respect to $L_{\operatorname{LB}}$ with relative bound $0$: indeed, since they correspond to left and right multiplication operators with respect to a finite rank operator by Corollary~\ref{cor:quadratic-finite-rank}, we get directly that the anticommutator parts $A_\alpha$ and $B_\alpha$ are all finite rank, and thus for all $x\in\mathscr{T}_2(\mathcal{H}_n)$,
\begin{align*}
\|\mathfrak{A}_\alpha x\|_2\le \|A_\alpha\|\,\|x\|_2\qquad \text{ and }\qquad \|\mathfrak{B}_\alpha x\|_2\le \|x\|_2\,\|B_\alpha\|\,.
\end{align*}
The same type of bound also holds for the last summand $\mathfrak{G}_\alpha^{(3)}$ of $\mathfrak{G}_\alpha$. The same argument also holds for the coherent term by \eqref{BHexpressionLa}.

In contrast to this, the perturbation of the jump operator is still unbounded. The corresponding unbounded contributions $\mathfrak{G}_\alpha^{(1)}$ and $\mathfrak{G}_\alpha^{(2)}$ are of the form $Z_1xZ_2,$ where either $Z_1,Z_2$ is always a finite rank operator and the other one is proportional to $a_{ik}$ or $a_{ik}^{\dagger}$. Fortunately, terms of the latter type are still zero bounded with respect to $L_{\operatorname{LB},n}$, so that Lemma \ref{lem:compact_resolv} can be applied. To show this, we make use of the operator 
\begin{align*}
\mathcal{N}:=\sum_{i\in[n]}\sum_{k\in[d]} n_{ik} \bullet + \bullet n_{ik}\equiv \sum_{i\in[n]}\sum_{k\in[d]}\mathcal{N}_{ik}.
\end{align*}
Then, consider for instance a term proportional to $a_{ik}x Z$ for some bounded operator $Z$ and $x\in D(\mathcal{N}_{ik}^{1/2})$. Then, 
\[
\|a_{ik}xZ\|_{2}\le \|Z\|\,\|a_{ij}x\|_{2}.
\]
Moreover,
\[
\|a_{ik}x\|_{2}^2
=\operatorname{Tr}\bigl((a_{ik}x)^\dagger(a_{ik}x)\bigr)
=\operatorname{Tr}(x^\dagger a_{ik}^\dagger a_{ik} x)
=\operatorname{Tr}(x^\dagger n_{ik} x).
\]
Since
\[
\langle x,\mathcal N_{ik} (x)\rangle
=\operatorname{Tr}\bigl(x^\dagger (n_{ik}x+xn_{ik})\bigr)
=\operatorname{Tr}(x^\dagger n_{ik}x)+\operatorname{Tr}(x^\dagger xn_{ik}),
\]
and both terms on the right are nonnegative, it follows that
\[
\operatorname{Tr}(x^\dagger n_{ik}x)\le \langle x,\mathcal N_{ik}(x)\rangle_{2}.
\]
Therefore
\[
\|a_{ik}xZ\|_{2}^2
\le \|Z\|^2\,\|a_{ik}x\|_{2}^2
=\|Z\|^2\,\langle x,\mathcal{N}_{ik}(x)\rangle
 \quad \Longrightarrow \quad \|a_{ik}xZ\|_2\le \|Z\|\,\|\mathcal{N}_{ik}^{1/2}(x)\|_2\,.
\]
Next, by the scalar inequality $\sqrt t\le \varepsilon t+\frac1{4\varepsilon}$, $t\ge0$ and functional calculus, since $\mathcal{N}_{ik}\ge 0$,
\[
\|\mathcal N_{ik}^{1/2}(x)\|_{2}\le \varepsilon \|\mathcal N_{ik} (x)\|_{2}+\frac1{4\varepsilon}\|x\|_{2},
\qquad x\in D(\mathcal N_{ik}).
\]
Hence
\[
\|a_{ik}xZ\|_{2}
\le \|Z\|\,\|\mathcal N_{ik}^{1/2}(x)\|_{2}
\le \varepsilon \|Z\|\,\|\mathcal N_{ik} (x)\|_{2}+\frac{\|Z\|}{4\varepsilon}\|x\|_{2}.
\]
Finally, since $\mathcal{N}_{ik}\le \mathcal{N}$, and since $L_{\operatorname{LB},n}= \frac{1}{2}(\nu_++\nu_-)^2\mathcal{N}-nd\nu_+(\nu_--\nu_+)\bullet $, we get constants $c_{ik}$ and $d_{ik}$ such that for any $\epsilon<1$,
\begin{align*}
\|a_{ik}xZ\|_2\le \varepsilon \,c_{ik}\|L_{\operatorname{LB},n}(x)\|_2\,+d_{ik}\Big(1+\frac{1}{\epsilon}\Big)\,\|x\|_2.
\end{align*}
The same control can be found for any of the terms $\mathfrak{G}_\alpha^{(1)}$ and $\mathfrak{G}_\alpha^{(2)}$. Thus, after summation and up to re-normalizing $\varepsilon$ accordingly, we find that the conditions of Lemma  \ref{lem:compact_resolv} are satisfied, and thus $L_{H_{n,M}}$ is gapped. 
 
\end{proof}


\subsection{Uniform gaps for weakly interacting gases}

In the weakly interacting regime, where \(B_n=\mathcal O(\varepsilon)\), obtained by choosing coupling strengths \(\lvert \alpha_{n,i,j}\rvert=\mathcal O(\varepsilon/n^2)\), the generator \(L_{\sigma_E,H_{nM}}\) has a positive spectral gap that is bounded uniformly in the particle number \(n\). Consequently, in this regime, the algorithm approximates the free energy densities efficiently and uniformly across all particle numbers. For our discussion of uniform gaps, we need a standard observation based on the min-max principle.
\begin{proposition}[Stability of low eigenvalues under form-small perturbations]
\label{prop:gap-form-small}
Let $A$ be a self-adjoint operator on a Hilbert space $\mathcal H$, bounded from below, with a compact resolvent. Let $C$ be a symmetric quadratic form with domain $Q(A)$ such that for some $\varepsilon\in(0,1)$, one has
\begin{equation}
\label{eq:form-small-eps}
|\langle u,Cu\rangle|
\le \varepsilon \langle u,(A+1)u\rangle ,
\qquad u\in Q(A).
\end{equation}
Then the form sum $A+C$ is well-defined, self-adjoint, bounded from below, and has compact resolvent. Moreover, if
\[
\lambda_1(A)\le \lambda_2(A)\le \cdots,\qquad
\lambda_1(A+C)\le \lambda_2(A+C)\le \cdots
\]
denote the eigenvalues counted with multiplicity, then for every $n\in \mathbb N$,
\begin{equation}
\label{eq:eigenvalue-comparison}
(1-\varepsilon)\lambda_n(A)-\varepsilon
\le \lambda_n(A+C)
\le (1+\varepsilon)\lambda_n(A)+\varepsilon.
\end{equation}
In particular, for the first spectral gap $\gap(A):=\lambda_2(A)-\lambda_1(A)$, $\gap(A+C):=\lambda_2(A+C)-\lambda_1(A+C)$, one has
\begin{equation}
\label{eq:gap-lower}
\gap(A)-\varepsilon\bigl(\lambda_1(A)+\lambda_2(A)+2\bigr)\le  \gap(A+C)\le \gap(A)+\varepsilon\bigl(\lambda_1(A)+\lambda_2(A)+2\bigr).
\end{equation}
\end{proposition}

\begin{proof}
By \eqref{eq:form-small-eps},
\begin{equation}
\label{eq:bound}
(1-\varepsilon)\langle u,Au\rangle-\varepsilon\|u\|^2
\le \langle u,(A+C)u\rangle
\le (1+\varepsilon)\langle u,Au\rangle+\varepsilon\|u\|^2,
\qquad u\in Q(A).
\end{equation}
Since $\varepsilon<1$, the form of $A+C$ is closed and lower bounded on $Q(A)$ by the KLMN theorem \cite{Schmdgen2012}, hence it defines a self-adjoint operator bounded from below. Because the form domain is the same as that of $A$, the compactness of the resolvent is preserved.
Now, apply the min-max principle. For any $n\in\mathbb N$,
\[
\lambda_n(A+C)
=
\inf_{\substack{\dim V=n\\V \subset Q(A)}}\ \sup_{\substack{u\in V\\ \|u\|=1}}
\langle u,(A+C)u\rangle.
\]
If we use the upper quadratic form bound \eqref{eq:bound}, we see that 
\[
\lambda_n(A+C)
\le
\inf_{\dim V=n}\ \sup_{\substack{u\in V\\ \|u\|=1}}
\bigl((1+\varepsilon)\langle u,Au\rangle+\varepsilon\bigr)
=
(1+\varepsilon)\lambda_n(A)+\varepsilon,
\]
while the lower quadratic form bound yields
\[
\lambda_n(A+C)
\ge
\inf_{\dim V=n}\ \sup_{\substack{u\in V\\ \|u\|=1}}
\bigl((1-\varepsilon)\langle u,Au\rangle-\varepsilon\bigr)
=
(1-\varepsilon)\lambda_n(A)-\varepsilon.
\]
This proves \eqref{eq:eigenvalue-comparison}.
To estimate the gap, combine the lower bound for $\lambda_2(A+C)$ with the upper bound for $\lambda_1(A+C)$:
\[
\gap(A+C)
=
\lambda_2(A+C)-\lambda_1(A+C)
\ge
\bigl((1-\varepsilon)\lambda_2(A)-\varepsilon\bigr)
-
\bigl((1+\varepsilon)\lambda_1(A)+\varepsilon\bigr).
\]
After rearranging,
\[
\gap(A+C)\ge \gap(A)-\varepsilon(\lambda_1(A)+\lambda_2(A)+2).
\]
 The upper bound follows analogously from the upper bound for $\lambda_2(A+C)$ and the lower bound for $\lambda_1(A+C)$.
\end{proof}



\noindent Now, by Theorem~\ref{thmfreenregyblabla}, the free energy densities 
\[
f_{\beta}(H_n):=\frac{F_{\beta}(H_n)}{n}
\qquad \text{ and }\qquad
f_{\beta}(H_{n,M}):=\frac{F_{\beta}(H_{n,M})}{n}
\]
satisfy, for $d\in\{2,3\}$, $\beta>0$, and interaction strengths obeying for some $\gamma\ge 1$ 
\begin{equation}
\label{eq:i_strengths}
|\alpha_{n,i,j}|\lesssim n^{-\gamma}\qquad(1\le i<j\le n),
\end{equation}
the estimate
\[
\bigl| f_{\beta}(H_n)-f_{\beta}(H_{n,M}) \bigr|
\le C_\beta\, n^{2-\gamma} M^{-\frac{1}{4d}}.
\]
In particular, if $\gamma\ge 2$, the right-hand side can be made arbitrarily small by choosing $M$ sufficiently large, independently of $n$. We may thus choose $\gamma=2$ in the sequel. It is meaningful to study free energy densities, as $F_{\beta}(H_{0,n})=n F_{\beta}(h),$ i.e. the densities are the unit-size objects in our scaling. 
For our argument, it is convenient to write 
\[H_{n,M}=H_{0,n} + W_{n,M}=H_{0,n}+\Pi_M \sum_{i<j} \alpha_{n,i,j} V_{ij} \Pi_M \text{ with } V_{ij}(x)=\begin{cases}
-\log \vert x_i-x_j \vert &\text{ if }d=2\\
    \vert x_i-x_j \vert^{-1} &\text{ if }d=3
\end{cases} \]
such that in this scaling $\Vert W_{n,M}\Vert$ is uniformly bounded in $n$.


\begin{proposition}[Locality yields an \(n^{-1}\) bound for the truncated remainder]
\label{prop:locality-remainder-1overn}

Fix a site \(i\in\{1,\dots,n\}\), and let \(A_i^\alpha\) be a local creation or annihilation operator acting only on the \(i\)-th tensor factor. For \(t\in\mathbb R\), define
\[
A_i^\alpha(t):=e^{itH_{0,n}}A_i^\alpha e^{-itH_{0,n}}.
\]
Assume that
\begin{equation}
\label{eq:uniform-local-commutator-bound}
\sup_{\substack{1\le i\neq j\le n\\ t\in\mathbb R}}
\big\|[\Pi_MV_{ij}\Pi_M,A_i^\alpha(t)]\big\|
\le C_M,
\end{equation}
for some constant \(C_M<\infty\) depending only on \(M\). Assume moreover that uniformly in j
\begin{equation}
\label{eq:mean-field-row-sum}
\sup_{1\le i\le n}\sum_{j\neq i}|\alpha_{n,i,j}|
\le \frac{\kappa}{n}
\end{equation}
for some \(\kappa>0\). Let \(f\in L^1(\mathbb R,(1+|t|)\,dt)\), we consider $L^\alpha_i(H):=\int_{\mathbb R} f(t)\,e^{itH}A_i^\alpha e^{-itH}\,dt$
for \(H=H_{0,n}\) and \(H=H_{n,M}\). Then there exists a constant \(C_{M,f}>0\), independent of \(n\) and \(i\), such that
\[
\big\|L_i^\alpha(H_{n,M})-L_i^\alpha(H_{0,n})\big\|
\le \frac{C_{M,f}}{n} \text{ with }C_{M,f}:=\kappa C_M \int_{\mathbb R}|t|\,|f(t)|\,dt.
\]

If, in addition, for some \(s\in\mathbb R\) the conjugated operators $L_{i,s}^\alpha(H):=e^{sH}L_i^\alpha(H)e^{-sH}$
are well-defined for \(H=H_{0,n},H_{n,M}\), and if
\begin{equation}
\label{eq:uniform-local-commutator-bound-conjugated}
\sup_{\substack{ 1\le i\neq j\le n\\ t\in\mathbb R}}
\big\|[\Pi_MV_{ij}\Pi_M,e^{itH_{0,n}}(e^{sH_{0,n}}A_i^\alpha e^{-sH_{0,n}})e^{-itH_{0,n}}]\big\|
\le C_{M,s},
\end{equation}
then
\[
\big\|L_{i,s}^\alpha(H_{n,M})-L_{i,s}^\alpha(H_{0,n})\big\|
\le \frac{C_{M,s,f}}{n},
\qquad
C_{M,s,f}:=\kappa C_{M,s}\int_{\mathbb R}|t|\,|f(t)|\,dt.
\]
\end{proposition}

\begin{proof}
We first prove the non-conjugated statement. By Duhamel's formula, for every \(t\in\mathbb R\),
\[
e^{itH_{n,M}}A_i^\alpha e^{-itH_{n,M}}-e^{itH_{0,n}}A_i^\alpha e^{-itH_{0,n}}
=
i\int_0^t e^{i(t-u)H_{n,M}}[W_{n,M},A_i^\alpha(u)]e^{-i(t-u)H_{n,M}}\,du
\]
for \(t\ge 0\), while for \(t\le 0\) the same identity holds with the integral taken over \([t,0]\). Hence, in either case,
\[
\big\|e^{itH_{n,M}}A_i^\alpha e^{-itH_{n,M}}-e^{itH_{0,n}}A_i^\alpha e^{-itH_{0,n}}\big\|
\le \int_0^{|t|}\big\|[W_{n,M},A_i^\alpha(\tau)]\big\|\,d\tau.
\]
It is therefore enough to estimate the commutator \([W_{n,M},A_i^\alpha(t)]\).

Since \(A_i^\alpha\), and thus \(A_i^\alpha(t)\), acts only on the \(i\)-th tensor factor and \(V_{jk}\) acts only on the \(j\)-th and \(k\)-th factors, we have
\[
[\Pi_MV_{jk}\Pi_M,A_i^\alpha(t)]=0
\qquad\text{whenever }i\notin\{j,k\}.
\]
Therefore, only the interaction terms touching site \(i\) survive:
\[
[W_{n,M},A_i^\alpha(t)]
=
\sum_{j\neq i}\alpha_{n,i,j}
\,[\Pi_MV_{ij}\Pi_M,A_i^\alpha(t)].
\]
Taking norms and using \eqref{eq:uniform-local-commutator-bound}, we obtain
\[
\big\|[W_{n,M},A_i^\alpha(t)]\big\|
\le
\sum_{j\neq i}|\alpha_{n,i,j}|\,
\big\|[\Pi_MV_{ij}\Pi_M,A_i^\alpha(t)]\big\|
\le
C_M\sum_{j\neq i}|\alpha_{n,i,j}|.
\]
By \eqref{eq:mean-field-row-sum},
\[
\big\|[W_{n,M},A_i^\alpha(t)]\big\|
\le \frac{\kappa C_M}{n},
\qquad t\in\mathbb R.
\]
Substituting this into the Duhamel estimate yields
\[
\big\|e^{itH_{n,M}}A_i^\alpha e^{-itH_{n,M}}-e^{itH_{0,n}}A_i^\alpha e^{-itH_{0,n}}\big\|
\le \frac{\kappa C_M}{n}\,|t|.
\]
Multiplying by \(|f(t)|\) and integrating over \(t\in\mathbb R\), we get
\[
\begin{aligned}
\big\|L_i^\alpha(H_{n,M})-L_i^\alpha(H_{0,n})\big\|
&\le
\int_{\mathbb R}|f(t)|\,
\big\|e^{itH_{n,M}}A_i^\alpha e^{-itH_{n,M}}-e^{itH_{0,n}}A_i^\alpha e^{-itH_{0,n}}\big\|\,dt\\
&\le
\frac{\kappa C_M}{n}\int_{\mathbb R}|t|\,|f(t)|\,dt.
\end{aligned}
\]
This proves the first claim.

For the conjugated statement, we proceed along the same lines and define
\[
B_{i,s}^\alpha:=e^{sH_{0,n}}A_i^\alpha e^{-sH_{0,n}}.
\]
Applying exactly the same argument to \(B_{i,s}^\alpha\) in place of \(A_i^\alpha\), and using the assumption \eqref{eq:uniform-local-commutator-bound-conjugated}, gives the analogous estimate
\[
\Big\|\int_{\mathbb R}f(t)\Big(e^{itH_{n,M}}B_{i,s}^\alpha e^{-itH_{n,M}}
-e^{itH_{0,n}}B_{i,s}^\alpha e^{-itH_{0,n}}\Big)\,dt\Big\|
\le \frac{\kappa C_{M,s}}{n}\int_{\mathbb R}|t|\,|f(t)|\,dt.
\]
Under the stated identification of the conjugated operators with these time-averaged expressions, this is precisely the desired estimate for the difference
\[
L_{i,s}^\alpha(H_{n,M})-L_{i,s}^\alpha(H_{0,n}).
\]
\end{proof}
In the next Lemma we directly verify the conditions of the previous result for our setting. 
\begin{lemma}[Verification of the local commutator bound for truncated Coulomb interactions]
\label{lem:local-commutator-harmonic-cutoff}
Let
\[
\mathcal H_n:=\bigotimes_{r=1}^n L^2(\mathbb R^d),\qquad
H_{0,n}:=\sum_{r=1}^n h_r,
\]
where each \(h_r\) is a d-dimensional harmonic oscillator, with $d \in \{2,3\}$ acting on the \(r\)-th tensor factor.
For \(1\le i<j\le n\), let
\[
V_{ij}(x):=\begin{cases}
-\log \vert x_i-x_j \vert &\text{ if }d=2\\
    \vert x_i-x_j \vert^{-1} &\text{ if }d=3.
\end{cases}
\]
Fix \(i\in\{1,\dots,n\}\), and let \(A_i^\alpha\) be a local creation or annihilation operator acting on the \(i\)-th tensor factor. Then, for every \(M\in\mathbb N\), there exists a constant \(C_M<\infty\), independent of \(n\), \(i\), \(j\), and \(t\), such that
\[
\sup_{\substack{1\le i\neq j\le n\\ t\in\mathbb R}}
\bigl\|[\Pi_MV_{ij}\Pi_M,A_i^\alpha(t)]\bigr\|
\le C_M.
\]
More precisely, one may take
\[\begin{split}
C_M&=
\Bigl\|
\Bigl[(P_M\otimes P_M)V_{ij}\,(P_M\otimes P_M),\,A^\alpha\otimes 1\Bigr]
\Bigr\|\le 2\,\bigl\|(P_M\otimes P_M)V_{ij}\,(P_M\otimes P_M)\bigr\|\,\|A^\alpha P_M\|\\
&\lesssim \bigl\|(P_M\otimes P_M)V_{ij}\,(P_M\otimes P_M)\bigr\|\ \sqrt{M+1}.
\end{split}
\]
\end{lemma}

\begin{proof}
Since \(H_{0,n}=\sum_{r=1}^n h_r\) and \(A_i^\alpha\) acts only on the \(i\)-th tensor factor, its free Heisenberg evolution is determined entirely by the one-particle oscillator on that factor:
\[
A_i^\alpha(t)=e^{itH_{0,n}}A_i^\alpha e^{-itH_{0,n}}=e^{ith_i}A_i^\alpha e^{-ith_i}.
\]
For the harmonic oscillator, the creation and annihilation operators are eigenoperators of the Heisenberg dynamics. Thus, there exists \(\theta_\alpha\in\mathbb R\) such that
\[
A_i^\alpha(t)=e^{i\theta_\alpha t}A_i^\alpha,
\]
and therefore $\bigl\|[\Pi_MV_{ij}\Pi_M,A_i^\alpha(t)]\bigr\|
=
\bigl\|[\Pi_MV_{ij}\Pi_M,A_i^\alpha]\bigr\|.$
It follows that it is enough to estimate the commutator at time \(t=0\).

We now use the tensor-product structure. Since \(V_{ij}\) acts only on the \(i\)-th and \(j\)-th variables and \(\Pi_M=P_M^{\otimes n}\), the compressed interaction \(\Pi_MV_{ij}\Pi_M\) acts trivially on all tensor factors except the \(i\)-th and \(j\)-th. Likewise, \(A_i^\alpha\) acts only on the \(i\)-th tensor factor. Hence, after identifying the \(i\)- and \(j\)-factors with two copies of \(L^2(\mathbb R^d)\), one may write
\[
\Pi_MV_{ij}\Pi_M
=
I_{\widehat{ij}}\otimes
\bigl((P_M\otimes P_M)V_{ij}\,(P_M\otimes P_M)\bigr),
\]
and
\[
A_i^\alpha
=
I_{\widehat{ij}}\otimes (A^\alpha\otimes 1),
\]
where \(I_{\widehat{ij}}\) denotes the identity, on $\operatorname{ran}(\Pi_M)$, on all tensor factors other than \(i\) and \(j\). Consequently,
\[
[\Pi_MV_{ij}\Pi_M,A_i^\alpha]
=
I_{\widehat{ij}}\otimes
\Bigl[(P_M\otimes P_M)V_{ij}\,(P_M\otimes P_M),\,A^\alpha\otimes 1\Bigr].
\]
Taking norms yields
\[
\bigl\|[\Pi_MV_{ij}\Pi_M,A_i^\alpha]\bigr\|
=
\Bigl\|
\Bigl[(P_M\otimes P_M)V_{ij}\,(P_M\otimes P_M),\,A^\alpha\otimes 1\Bigr]
\Bigr\|.
\]
The right-hand side depends only on the two-particle compressed operator and is therefore independent of \(n\), \(i\), and \(j\). This proves the claimed reduction.

It remains to show that this norm is finite. Set $E_M:=P_ML^2(\mathbb R^d).$
Since \(P_M\) is a spectral projection of the harmonic oscillator onto energies \(\le M\), the space \(E_M\) is finite-dimensional. Therefore, \(E_M\otimes E_M\) is finite-dimensional as well, and
\[
(P_M\otimes P_M)V_{ij}\,(P_M\otimes P_M)
\]
is a bounded operator on \(E_M\otimes E_M\). Moreover, \(A^\alpha P_M\) is bounded because it acts on the finite-dimensional space \(E_M\). It follows that the commutator
\[
\Bigl[(P_M\otimes P_M)V_{ij}\,(P_M\otimes P_M),\,A^\alpha\otimes 1\Bigr]
\]
is bounded on \(E_M\otimes E_M\), and hence has finite norm. Defining
\[
C_M:=
\Bigl\|
\Bigl[(P_M\otimes P_M)V_{ij}\,(P_M\otimes P_M),\,A^\alpha\otimes 1\Bigr]
\Bigr\|,
\]
we conclude that
\[
\sup_{\substack{1\le i\neq j\le n\\ t\in\mathbb R}}
\bigl\|[\Pi_MV_{ij}\Pi_M,A_i^\alpha(t)]\bigr\|
\le C_M.
\]

To see where the explicit bound comes from, we write
\[
X_M:=(P_M\otimes P_M)V_{ij}\,(P_M\otimes P_M).
\]
Then
\[
[X_M,A^\alpha\otimes 1]=X_M(A^\alpha\otimes 1)-(A^\alpha\otimes 1)X_M,
\]
so by the triangle inequality,
\[
\|[X_M,A^\alpha\otimes 1]\|
\le \|X_M(A^\alpha\otimes 1)\|+\|(A^\alpha\otimes 1)X_M\|
\le 2\|X_M\|\,\|A^a P_M\|.
\]
This proves
\[
C_M\le 2\,\bigl\|(P_M\otimes P_M)V_{ij}\,(P_M\otimes P_M)\bigr\|\,\|A^\alpha P_M\|.
\]

Finally, for a one-particle creation or annihilation operator \(A^\alpha\), the standard ladder-operator bounds on the oscillator basis imply
\[
\|A^\alpha P_M\|\lesssim \sqrt{M+1},
\]
which gives the last claim.
\end{proof}

We have thus established that each perturbation of the $j$-th Lindblad operator is a finite-rank operator with a norm bound $\mathcal O(\sum_{i\neq j}|\alpha_{n,i,j}|).$ This allows us to conclude that the spectral gap of the unperturbed Lindbladian associated with $H_{0,n}$ persists, as long as the sum of their norms does not exceed the size of the gap. 
\begin{theorem}
\label{thm:uniform_gap}
Let $\alpha_{n,i,j} \lesssim \varepsilon(M)/ n^{2}$ for $\varepsilon(M) >0$ sufficiently small and $f \in L^1(\mathbb R; (1+|t|) \ dt)$, then the Lindbladian $L_{H_{n,M}}$ associated with the Hamiltonian $H_{n,M}$ has a uniform gap in $n$.
\end{theorem}
\begin{proof}
When studying the Lindbladian associated with $H_{n,M}$, the perturbations to the anti-commutator and commutator terms in the Lindbladian associated with $H_{0,n}$ involve products of creation and annihilation operators with the finite rank operator difference of the respective Lindblad operators; see \eqref{eq:anti-comm_terms}. Since the norm of this difference for the $j$-th Lindblad operator is, by Proposition \ref{prop:locality-remainder-1overn} and Lemma \ref{lem:local-commutator-harmonic-cutoff}, bounded by $\mathcal O(\sum_{i\neq j}|\alpha_{n,i,j}|)=\mathcal O(\varepsilon(M)/n)$. Summing $n$ many such terms still leads only to a perturbation of size $\mathcal O(\varepsilon(M))$, which does not close the unit-size gap of the unperturbed Lindbladian by choosing $\varepsilon(M)$ sufficiently small. 


It is only the perturbation of the jump operator that involves unbounded perturbations, but we may proceed as in the proof of Theorem \ref{corr:Gaussian_gap}:

Denoting the perturbation obtained from the jump operators of the $j$-th particle Lindblad operator by $C_j$, we see that each of them satisfies the relative bound with respect to $A_j(x)=\frac{1}{2}(\nu_+-\nu_-)^2\{N_j,x\}$ for $\varepsilon',\delta=\mathcal O(\varepsilon/n)$
\begin{equation}
|\langle u,C_ju\rangle|
\le \varepsilon' \langle u,A_j u\rangle + \delta \|u\|^2. 
\end{equation}
This computation makes up the last part of the proof of Theorem \ref{corr:Gaussian_gap}.
Then, from \cite[Proof of Thm.~6.3]{OU}, we may use that $Ax=L_{\operatorname{LB}}^{\oplus_{i=1}^n}x+n \nu_+(\nu_--\nu_+)x \le L_{H_{0,n}} $, we have 
\[\sum_{i} |\langle u,C_i u\rangle|
\le \varepsilon' \langle u,L_{H_{0,n}} u\rangle + \left(\delta n +  \varepsilon' n\nu_+ | \nu_--\nu_+| \right) \|u\|^2 \]
such that for $\eta = \max\{\varepsilon',\delta n +  \varepsilon' n\nu_+ |\nu_--\nu_+| \} =\mathcal O(\varepsilon)$ we have 
\[\sum_{i} |\langle u,C_i u\rangle| \le \eta \langle u,(L_{H_{0,n}}+1) u\rangle= \mathcal O( \varepsilon \langle u,(L_{H_{0,n}}+1) u\rangle). \]

We thus see that we can apply Proposition \ref{prop:gap-form-small} to the self-adjoint Lindbladian on the space of Hilbert-Schmidt operators, as the unperturbed operator $L_{H_{0,n}}$ exhibits a uniform gap, as it is just a $n$-fold direct sum, to conclude the existence of a uniform gap for the perturbed one as well. 
\end{proof}




\section{Algorithmic implementations}
\subsection{Preparing Gibbs states of trapped Coulomb gases}
\label{sec:GibbscircuitCoulomb}
In the following, we apply the general framework of \cite[Section 4.3 \& 4.5]{BeckerRouzeSalzmannToAppearcmp} to provide a finite-dimensional circuit implementation for the Gibbs state of the trapped Coulomb gases considered in this paper, see Theorem~\ref{thm:GibbsPrepareCoulomb} below. In particular, we employ the Gibbs sampler generated by $\cL_{\sigma_E,H}$ with bare jumps $\{A^\alpha\}_{\alpha\in\cA} \equiv \{a_{i,j},a_{i,j}^\dagger\}_{i\in[n],j\in[d]}.$  As argued in \cite[Theorem 4.12]{BeckerRouzeSalzmannToAppearcmp} the dynamics $e^{t\cL_{\sigma_E,H}}$ can be well-approximated by the finite-dimensional dynamics whose generator can be obtained by replacing the unbounded bare jumps and Hamiltonian with their truncated, finite-dimensional counterparts
\begin{align}
\label{eq:DefTruncOperators}
    a^{\le \widetilde M}_{i,j} := P^i_{\widetilde M} a_{i,j} P^i_{\widetilde M},\quad\quad \left(a_{i,j}^{\le\widetilde M}\right)^\dagger = P^i_{\widetilde M}\,a_{i,j}^\dagger P^i_{\widetilde M} =: \left(a_{i,j}^\dagger\right)^{\le \widetilde M}, \quad\quad H_{\le \widetilde M} := \Pi_{\widetilde M} \,H\,\Pi_{\widetilde M}
\end{align}
for some truncation level $\widetilde M\in \N$ and where\footnote{, note that in the notation of \cite{BeckerRouzeSalzmannToAppearcmp}, we 
hence take for $i\in[n],$ $j\in[d]$ and multi-index $\alpha = (i,j)$ the projection $\pi^{\alpha}_{\widetilde M} \equiv P^{i}_{\widetilde M+1},$, which is hence independent of the specific label $j\in[d].$
}  $P^i_{\widetilde M} :=\1_{L^2(\R^{d(i-1)})}\otimes P_{\widetilde M}\otimes  \1_{L^2(\R^{d(n-i-1)})}$ and $P_{\widetilde M}$ denote the projection onto the $\widetilde M$-first eigenstates of $h:=-\Delta + |x|^2$ on $L^{2}(\R^{d})$ and $\Pi_{\widetilde M} := \left(P_{\widetilde M}\right)^{\otimes n}.$

Here, we consider Hamiltonians that are given as finite-rank perturbations of the operator $H_{0,n}$, i.e.
\begin{align*}
    H:= H_{0,n} + V,
\end{align*}
where $V$ is self-adjoint and satisfies $P^{\otimes n}_M VP^{\otimes n}_M$ for some $M\le \widetilde M.$ Furthermore, we consider a self-adjoint positive semidefinite operator $\NA$, which is used in \cite{BeckerRouzeSalzmannToAppearcmp} to regularize within the finite dimensional implementation scheme, a shift of the Hamiltonian, namely
\begin{align*}
    \NA := \frac{1}{2}\left(H- dn + \|V\|\right) = \Ntot + \frac{1}{2}\left(V + \|V\|\right)   \ge 0,
\end{align*}
where for the last equality we denoted the total number operator $\Ntot = \sum_{i=1}^{n}\sum_{j=1}^d a^\dagger_{i,j}a_{i,j}.$
For these choices,  we see in the following lemma that the conditions on the bare jumps and the Hamiltonian employed in the efficient implementation theorem \cite[Theorem~4.12]{BeckerRouzeSalzmannToAppearcmp} are satisfied.
\begin{lemma}
\label{lem:VerificationCondFiniteDimImpl}
Let $n,d\in\N$ and $H_{0,n}:=\sum_{i\in[n]}h_i$,
where $h_i = -\Delta_i +|x_i|^2$ and $\Delta$ stand for $d$-dimensional Laplacian and $|x_i|^2=\sum_{j\in [d]}x_{ij}^2$. Let further $P_M$ be the spectral projection onto the first $M$ lowest eigenstates of the single mode number operator $h := -\Delta + |x|^2$ on $L^2(\R^d)$. 
Then, define
\begin{align*}
H = H_{0,n} + V,\qquad\quad
N_{\mathcal A}= \frac{1}{2}\left(H - dn + \|V\|\right)\ge 0,
\end{align*}
where $V$ is some finite-rank self-adjoint operator satisfying $P^{\otimes n}_{M}VP^{\otimes n}_{M}= V.$ 
Then, for $\kappa\in(0,1]$, there exists a constant \(C_\kappa>0\)
such that for all \(\widetilde M\in\mathbb N\) satisfying $\widetilde M \ge M+2$, we have
\begin{align}
\label{eq:NormTruncatedJump}
\left\|a_{i,j}^{\le \widetilde{M}}\right\|\,,\,\left\|\left(a_{i,j}^{\le \widetilde{M}}\right)^{\dagger}\right\|\le \sqrt{\widetilde M+1},
\end{align}
\begin{align}
\label{eq:AbstractFiniteTruncBareJumps}
\left\|e^{(l-1)N^\kappa}\bigl(a_{i,j}-(a_{i,j})^{\le \widetilde{M}}\bigr)e^{-lN_{\mathcal A}^\kappa}\right\|\,,\, \left\|e^{(l-1)N^\kappa}\bigl(a_{i,j}^\dagger-(a_{i,j}^{\le \widetilde{M}})^\dagger \bigr)e^{-lN_{\mathcal A}^\kappa}\right\|
\le C_{\kappa} \sqrt{\widetilde M+1}\,e^{-\widetilde M^\kappa}
\end{align}
for $l=1,2$ and furthermore
\begin{align}
\label{eq:AnniCraNormTrunc}
\|e^{\NA^\kappa}a_{i,j}^{\le \widetilde{M}}e^{-N_{\mathcal A}^\kappa}\|
\,,\,
\|e^{-N_{\mathcal A}^\kappa}(a_{i,j}^{\le \widetilde{M}})^\dagger e^{N_{\mathcal A}^\kappa}\| &\le C_{\kappa} \sqrt{\widetilde M} e^{(nM)^\kappa + (\|V\|/2)^\kappa},\\ 
\label{eq:CABound}
\|a_{i,j} e^{-N_{\mathcal A}^\kappa}\|\, , \,
\|e^{-N_{\mathcal A}^\kappa}a_{i,j}\|\,,\, \|e^{N_{\mathcal A}^\kappa}a_{i,j} e^{-2N_{\mathcal A}^\kappa}\|\,,\,\|e^{-2N_{\mathcal A}^\kappa}a_{i,j} e^{N_{\mathcal A}^\kappa}\|  &\le C_\kappa\, e^{2((nM)^\kappa + (\|V\|/2)^\kappa)}
\end{align}
Moreover, we have
\begin{align}
\label{eq:HamiltonianLeakTrunc}
\|(H-H_{\le \widetilde{M}})e^{-N_{\mathcal A}^\kappa}\|
\le C_\kappa (2\widetilde M+dn)e^{-\widetilde M^\kappa},
\end{align}
and for \(k=2,4\) and all \(t\in\mathbb R\),
\begin{align}
\label{eq:ExpBound2}
e^{-itH}e^{kN_{\mathcal A}^\kappa}e^{itH}=e^{kN_{\mathcal A}^\kappa}.
\end{align}
\bin{In particular, one may take
\[
q(\widetilde M)=C_{\kappa} \sqrt{\widetilde M+1},
\qquad
q'(\widetilde M)=2 \sqrt{\widetilde M+1},
\qquad
p_n(|\mathcal A|,M)=C_{\kappa,n}(\widetilde M+1),
\qquad
r=0.
\]
}
\end{lemma}

\begin{proof}
Let \(\{|\mathbf n\rangle\}_{\mathbf n\in\mathbb N^{dn}}\) be the occupation number basis,
where
\[
\mathbf n=(n_1,\dots,n_{dn}),
\qquad
|\mathbf n|:=\sum_{j=1}^{dn} n_j.
\]
Then
\[
H_{0,n}|\mathbf n\rangle=|\mathbf n|\,|\mathbf n\rangle,
\]
and for each site
\[
a_{i,j}|\mathbf n\rangle
=
\begin{cases}
\sqrt{n_{i,j}}\,|\mathbf n-e_{i,j}\rangle,& n_{i,j}\ge 1,\\[1mm]
0,& n_{i,j}=0,
\end{cases}
\qquad
a_{i,j}^\dagger|\mathbf n\rangle
=
\sqrt{n_{i,j}+1}\,|\mathbf n+e_{i,j}\rangle.
\]
Since \(n_{i,j}\le |\mathbf n|\), all relevant coefficients are controlled by the total
occupation number.

We first verify \eqref{eq:NormTruncatedJump}. On \(\operatorname{Ran} P_{\widetilde M}^{\otimes n}\), one has
\(n_{i,j}\le \widetilde{M}\), hence
\[
\|P_{\widetilde M}^{\otimes n}a_{i,j} P_{\widetilde M}^{\otimes n}\|\le \sqrt {\widetilde M},
\qquad
\|P_{\widetilde M}^{\otimes n}a_{i,j}^\dagger P_{\widetilde M}^{\otimes n}\|\le \sqrt{\widetilde M+1}.
\]
Thus
\[
\|(a_{i,j})^{\le \widetilde{M}}\|\le \sqrt{\widetilde M+1},
\]
so \eqref{eq:NormTruncatedJump} holds. We continue with our proof focussing  on verifying the conditions for \(a_{i,j}\) as the ones for $a_{i,j}^\dagger$ follow similarly.


Next, we verify \eqref{eq:ExpBound2}. Since by definition the operator \(H\) commutes with \(N_{\mathcal A}\),
hence also with \(N_{\mathcal A}^\kappa\) and \(e^{kN_{\mathcal A}^\kappa}\). Therefore, for all
\(k,t\in\mathbb R\) we have
\[
e^{-itH}e^{kN_{\mathcal A}^\kappa}e^{itH}=e^{kN_{\mathcal A}^\kappa}.
\]


We next verify \eqref{eq:AbstractFiniteTruncBareJumps}. For this 
set \(R^{\widetilde M}_{i,j}:=a_{i,j}-P^{i}_{\widetilde M}a_{i,j} P^{i}_{\widetilde M}\) and note that since $\widetilde M\ge M$ we have $R^{\le \widetilde M}_{i}P^{\otimes n}_{M} = 0.$
Therefore, using the block diagonality \begin{align*}
\NA &= P^{\otimes n}_{M}\NA P^{\otimes n}_{M} + (1-P^{\otimes n}_{M})\NA (1-P^{\otimes n}_{M}) \\&= P^{\otimes n}_{M}\NA P^{\otimes n}_{M} +(1-P^{\otimes n}_{M}) (\Ntot+\|V\|/2) (1-P^{\otimes n}_{M})
\end{align*}
we see for $l=1,2$ that 
\begin{align}
    R^{\le \widetilde M}_{i,j} e^{-l\NA^{\kappa}} = R^{\le \widetilde M}_{i,j} (1-P^{\otimes n}_{M})e^{-l(\Ntot+\|V\|/2)^{\kappa}}(1-P^{\otimes n}_{M}) 
\end{align}
and therefore
\[
\|R^{\widetilde M}_{i,j} e^{-N_{\mathcal A}^\kappa}\|
\lesssim \sup_{r\ge \widetilde M+1}\sqrt r\,e^{-(r+\|V\|/2)^\kappa} \le \sup_{r\ge \widetilde M+1}\sqrt r\,e^{-r^\kappa} \le C_{\kappa} \sqrt{\widetilde M+1}\, e^{-\widetilde M^\kappa},
\]
where we used in the last inequality that \(r\mapsto \sqrt r\,e^{-r^\kappa}\) is eventually decreasing.  Similarly we obtain
\[
\|e^{\NA^\kappa}R^{\widetilde M}_{i,j} e^{-2N_{\mathcal A}^\kappa}\|
\lesssim \sup_{r\ge \widetilde M+1}\sqrt r\,e^{(r-1+\|V\|/2)^\kappa-2(r+\|V\|/2)^\kappa} \le \sup_{r\ge \widetilde M+1}\sqrt r\,e^{-r^\kappa} \le C_{\kappa} \sqrt{\widetilde M+1} e^{-\widetilde M^\kappa}.
\]
The argument for the $a_{i,j}^{\dagger}$ is analogous and proves the claim for \(l=1,2\).


We continue by showing \eqref{eq:AnniCraNormTrunc} and writing
\[\begin{split}
&e^{N_{\mathcal A}^\kappa}P_{\widetilde M}^{\otimes n}a_{i,j} P_{\widetilde M}^{\otimes n}e^{-N_{\mathcal A}^\kappa}= e^{\NA^{\kappa}}P_{M}^{\otimes n}a_{i,j}P_{M}^{\otimes n}e^{-\NA^{\kappa}} +e^{\NA^{\kappa}}(P_{\widetilde M}^{\otimes n}-P_{M}^{\otimes n})a_{i,j} P_{M}^{\otimes n}e^{-\NA^{\kappa}} \\&\qquad\quad+ e^{\NA^{\kappa}}P^{\otimes n}_{M}a_{i,j} (P^{\otimes n}_{\widetilde M}-P_{M}^{\otimes n})e^{-\NA^{\kappa}} + e^{\NA^{\kappa}}(P^{\otimes n}_{\widetilde M}-P^{\otimes n}_{M})a_{i,j} (P^{\otimes n}_{\widetilde M}-P_{M}^{\otimes n})e^{-\NA^{\kappa}}\\&\quad= e^{\NA^{\kappa}}P_{M}^{\otimes n}a_{i,j}P_{M}^{\otimes n}e^{-\NA^{\kappa}} +(P_{\widetilde M}^{\otimes n}-P_{M}^{\otimes n})e^{(\Ntot + \|V\|/2)^{\kappa}}a_{i,j} P_{M}^{\otimes n}e^{-\NA^{\kappa}} \\&\qquad\quad+ e^{\NA^{\kappa}}P^{\otimes n}_{M}a_{i,j} e^{-(\Ntot + \|V\|/2)^{\kappa}}(P^{\otimes n}_{\widetilde M}-P_{M}^{\otimes n}) \\&\qquad\quad+ (P^{\otimes n}_{\widetilde M}-P^{\otimes n}_{M})e^{(\Ntot + \|V\|/2)^{\kappa}} a_{i,j} e^{-(\Ntot + \|V\|/2)^{\kappa}}  (P^{\otimes n}_{\widetilde M}-P_{M}^{\otimes n}),
\end{split}
\]
where for the last equality we have used that $\NA(P^{\otimes n}_{\widetilde M} - P^{\otimes n}_{M}) = (\Ntot + \|V\|/2)(P^{\otimes n}_{\widetilde M} - P^{\otimes n}_{M})= (P^{\otimes n}_{\widetilde M} - P^{\otimes n}_{M})(\Ntot + \|V\|/2).$ We bound the norm of each term individually using
\begin{align*}
P^{\otimes n}_{M}\NA P^{\otimes n}_{M} \le nM +\|V\|
\end{align*}
which gives 
\begin{align*}
&\left\|e^{\NA^{\kappa}}P_{M}^{\otimes n}a_{i,j}P_{M}^{\otimes n}e^{-\NA^{\kappa}}\right\| \le \left\|e^{\NA^{\kappa}}P_{M}^{\otimes n}\right\|\left\|P_{M}^{\otimes n}a_{i,j}P_{M}^{\otimes n}\right\| \le \sqrt{M} \,e^{(nM)^\kappa + (\|V\|/2)^\kappa},\\
&\left\|(P_{\widetilde M}^{\otimes n}-P_{M}^{\otimes n})e^{(\Ntot + \|V\|/2)^{\kappa}}a_{i,j} P_{M}^{\otimes n}e^{-\NA^{\kappa}}\right\| \le \sqrt{M}   \| e^{(\Ntot + \|V\|/2)^{\kappa}}P^{\otimes n}_{M - e_{i,j}}\| \le \sqrt{M} e^{(nM)^\kappa + (\|V\|/2)^{\kappa}},\\&\left\|
e^{\NA^{\kappa}}P^{\otimes n}_{M}a_{i,j} e^{-(\Ntot + \|V\|/2)^{\kappa}}(P^{\otimes n}_{\widetilde M}-P_{M}^{\otimes n})\right\| \le \sqrt{M}\, e^{(nM)^\kappa  +(\|V\|/2)^\kappa},\\ &\left\| e^{(\Ntot + \|V\|/2)^{\kappa}} a_{i,j} e^{-(\Ntot + \|V\|/2)^{\kappa}}  (P^{\otimes n}_{\widetilde M}-P_{M}^{\otimes n})\right\| =\left\| e^{(\Ntot + \|V\|/2)^{\kappa} -(\Ntot + \|V\|/2+1)^{\kappa}} a_{i,j}   (P^{\otimes n}_{\widetilde M}-P_{M}^{\otimes n})\right\|\\&\qquad\lesssim \sqrt{\widetilde M},
\end{align*}
where we denoted $P^{\otimes n}_{M -e_{i,j}} = P^{\otimes d(i-1) +j-1}_{M}\otimes P_{M -1}\otimes P^{\otimes dn -d(i-1) -j}_{M}$ and consistently used that $(x+y)^\kappa \le x^\kappa +y^\kappa$ for $x,y\ge 0.$

We continue with proving \eqref{eq:CABound}, which follows using a similar argument: We write $$a_{i,j}e^{-\NA^\kappa} =a_{i,j}P^{\otimes n}_{M}e^{-\NA^\kappa} + a_{i,j}(1-P^{\otimes n}_{M})e^{-\NA^\kappa} = a_{i,j}P^{\otimes n}_{M}e^{-\NA^\kappa} + a_{i,j}e^{-(\Ntot + \|V\|/2)^\kappa}(1-P^{\otimes n}_{M})$$ and bound each term individually as
\begin{align*}
 \left\| a_{i,j}P^{\otimes n}_{M}e^{-\NA^\kappa} \right\|  \le \sqrt{M},\qquad \left\|a_{i,j}e^{-(\Ntot + \|V\|/2)^\kappa}(1-P^{\otimes n}_{M})\right\| \le C_\kappa.
\end{align*}
Moreover, 
we write
\begin{align*}
e^{\NA^\kappa}a_{i,j}e^{-2\NA^\kappa} &= e^{\NA^\kappa}P^{\otimes n}_{M}a_{i,j}P^{\otimes n}_{M}e^{-2\NA^\kappa} + (1-P^{\otimes n}_{M}) e^{(\Ntot + \|V\|/2)^\kappa}a_{i,j}e^{-2\NA^\kappa}P^{\otimes n}_{M} \\&\qquad +e^{\NA^\kappa}P^{\otimes n}_{M}a_{i,j}(1-P^{\otimes n}_{M})e^{-2(\Ntot+ \|V\|/2)^\kappa}\\&\qquad  +(1-P^{\otimes n}_{M})e^{(\Ntot + \|V\|/2)^\kappa}a_{i,j}e^{-2(\Ntot + \|V\|/2)^\kappa}(1-P^{\otimes n}_{M})  
\end{align*}
and bound the operator norm of each of the four terms individually as 
\begin{align*}
&\|e^{\NA^\kappa}P^{\otimes n}_{M}a_{i,j}P^{\otimes n}_{M}e^{-2\NA^\kappa}\| \le \sqrt{M} e^{(nM)^\kappa + (\|V\|/2)^\kappa} 
\\ &\|e^{(\Ntot + \|V\|/2)^\kappa}a_{i,j}e^{-2\NA^\kappa}\| \lesssim e^{2((nM)^\kappa +(\|V\|/2)^\kappa)} \|e^{(H_{0,n} +\|V\|)^\kappa}a_{i,j}e^{-2(\Ntot + \|V\|/2)^\kappa}\| \le C_{\kappa}e^{2((nM)^\kappa +(\|V\|/2)^\kappa)}\\&\|e^{\NA^\kappa}P^{\otimes n}_{M}a_{i,j}(1-P^{\otimes n}_{M})e^{-2(H_{0,n}+\|V\|)^\kappa}\| \le C_\kappa e^{(nM)^\kappa + (\|V\|/2)^\kappa} \\&\|(1-P^{\otimes n}_{M})e^{(H_{0,n} +\|V\|)^\kappa}a_{i,j}e^{-2(\Ntot + \|V\|/2)^\kappa}(1-P^{\otimes n}_{M}) \| \le C_\kappa
\end{align*}
where for the second line we have used
\begin{align*} \left\|e^{2(\Ntot + \|V\|/2)^\kappa}e^{-2\NA^\kappa}\right\| &\le \left\|e^{2(\Ntot + \|V\|/2)^\kappa}P^{\otimes n}_{M}e^{-2\NA^\kappa}\right\| + \left\|e^{2(\Ntot + \|V\|/2)^\kappa}e^{-2(\Ntot + \|V\|/2)^\kappa}(1-P^{\otimes n}_{M})\right\| \\&  \le e^{2((nM)^\kappa + (\|V\|/2)^\kappa)} +1. 
\end{align*}
The remaining two terms in \eqref{eq:CABound} follow similarly.

Lastly we check \eqref{eq:HamiltonianLeakTrunc}. Since $\widetilde M \ge M$ we have 
\[
H_{\le \widetilde{M}}=P_{\widetilde M}^{\otimes n}HP_{\widetilde M}^{\otimes n}=
H_{0,n}P_{\widetilde M}^{\otimes n} + V,
\]
and hence
\[
H-H_{\le \widetilde{M}}=H_{0,n}(1-P_{\widetilde M}^{\otimes n}) 
\]
Therefore, we have 
\[\begin{split}
(H-H_{\le \widetilde{M}})e^{-N_{\mathcal A}^\kappa}&=H_{0,n}(1-P_{\widetilde M}^{\otimes n})e^{-N_{\mathcal A}^\kappa}=H_{0,n}(1-P_{\widetilde M}^{\otimes n})e^{-(\Ntot+\|V\|/2)^\kappa}\\
&= H_{0,n}e^{-(\Ntot+\|V\|/2)^\kappa}(1-P_{\widetilde M}^{\otimes n})  = (2\Ntot +dn)e^{-(\Ntot+\|V\|/2)^\kappa}(1-P_{\widetilde M}^{\otimes n})
\end{split}
\]
By functional calculus we hence see
\[
\|(H-H_{\le \widetilde{M}})e^{-N_{\mathcal A}^\kappa}\|
\le  \sup_{\mathbf{n}\in\N^n_0,\,|\mathbf{n}|>\widetilde M}(2|\mathbf{n}|+dn) e^{-|\mathbf{n}|^\kappa} \le C_\kappa (2\widetilde M + dn)e^{-\widetilde M^\kappa}.
\]
\end{proof}
We continue with three supporting lemmas to estimate the scaling of certain constants appearing in the implementation scheme of \cite[Theorem~4.12 \& Corollary~4.33]{BeckerRouzeSalzmannToAppearcmp} and then state and prove the main result of this section, Theorem~\ref{thm:GibbsPrepareCoulomb}. 
\begin{lemma}
\label{lem:PartionFunctionBound}
Let $d,n\in\N$ and $H =  H_{0,n} +V$ for some bounded self-adjoint operator $V.$ Then the partition function for $H$ and $\beta >0$ satisfies
\begin{align*}
    e^{-\beta\|V\|} (2\sinh(\beta))^{-dn}\le\mathcal{Z}_\beta(H) \le e^{\beta\|V\|} (2\sinh(\beta))^{-dn}.
\end{align*}
\end{lemma}
\begin{proof}
 Note that as $V$ is bounded we have 
$H_{0,n} +\|V\|\ge H \ge H_{0,n} - \|V\|$ and, hence, by the min-max principle for eigenvalues 
\begin{align*}
    \lambda_j(H_{0,n}) + \|V\|\ge \lambda_j(H) \ge \lambda_j(H_{0,n}) - \|V\|
\end{align*}
where $\lambda_j(H)$ and $\lambda_j(H_{0,n})$ denote the $j^{th}$ smallest eigenvalue of $H$ and $H_{0,n}$ respectively. Therefore, we see
\begin{align*}
\mathcal{Z}_\beta(H) &= \Tr\left(e^{-\beta H}\right) \le e^{\beta\|V\|} \Tr\left(e^{\beta H_{0,n}}\right) ´= e^{\beta\|V\|}  \left(\sum_{k=0}^{\infty} e^{-\beta (2k+1)} \right)^{dn} =e^{\beta\|V\|} \left(\frac{e^{-\beta}}{1-e^{-2\beta}}\right)^{dn} \\&=e^{\beta\|V\|} (2\sinh(\beta))^{-dn}   
\end{align*}
and similarly 
\begin{align*}
\mathcal{Z}_\beta(H) \ge e^{-\beta\|V\|} (2\sinh(\beta))^{-dn}.
\end{align*}
\end{proof}
\begin{lemma}
\label{lem:cScalingFiniteRank}
Let $d,n\in\N$ and $H =  H_{0,n} +V$ for some finite-rank self-adjoint operator satisfying $V = P^{\otimes n}_{M}VP^{\otimes n}_{M}.$
Then for all normalised $\ket{\Psi} = P^{\otimes n}_{M}\ket{\Psi}$ we have
\begin{align*}
    \kb{\Psi} \le e^{\beta(2 \|V\|+nM)} (2\sinh(\beta))^{-dn} \ \sigma_\beta(H).
\end{align*}
In particular, the bound applies to  $\ket{\Psi} = \ket{0}^{\otimes dn},$ where $\ket{0}$ denotes the Fock state  satisfying $N\ket{0} =0$, where $N$ is the single mode number operator on $L^2(\R).$  
\end{lemma}
\begin{proof}
As $\kb{\Psi}$ is rank one, it is clear that the smallest constant $\mathfrak{c}$ such that 
\begin{align*}
\kb{\Psi} \le  \mathfrak{c}\, \sigma_\beta(H)  \qquad\text{is given by}  \qquad
\mathfrak{c} = \bra{\Psi}\sigma^{-1}_\beta(H)\ket{\Psi} = \mathcal{Z}_\beta(H) \bra{\Psi}e^{\beta H}\ket{\Psi}.
\end{align*}
We focus on bounding the term $\bra{\Psi}e^{\beta H}\ket{\Psi}:$
For that, note that as $[H_{0,n},P^{\otimes n}_{M}] = [V,P^{\otimes n}_{M}] = 0$ and $P^{\otimes n}_{M} \ket{\Psi} = \ket{\Psi}$ we have $e^{\beta H}\ket{\Psi} = e^{\beta P^{\otimes n}_{M}HP^{\otimes n}_{M}}\ket{\Psi}.$ Using $P^{\otimes n}_{M}H_{0,n}P^{\otimes n}_{M} \le nM$ and therefore $\|P^{\otimes n}_{M}HP^{\otimes n}_{M}\| \le nM +\|V\|$ we conclude
\begin{align*}
\bra{\Psi}e^{\beta H}\ket{\Psi} \le e^{\beta(nM + \|V\|)}.
\end{align*}
Combining this with the upper bound of the partition function in Lemma~\ref{lem:PartionFunctionBound} finishes the proof.
\end{proof}

\begin{lemma}
\label{lem:EGibbsBoundedperturbation}
Let $H= H_{0,n} + V$ for some bounded self-adjoint operator $V.$ Define the self-adjoint and positive semidefinite operator $\NA = \frac{1}{2}\left(H -dn+ \|V\|\right).$
Then, for all $\kappa\in[0,1)$, we have 
\begin{align}
E_{\operatorname{Gibbs}} := \Tr\left(e^{4\NA^\kappa} \,\sigma_\beta(H)\right) \le C_{\beta,\kappa} \,e^{\beta(2\|V\|-dn/2)}  \left(\frac{\sinh(\beta)}{\sinh(\beta/2)}\right)^{dn},
\end{align}
for some constant $C_{\beta,\kappa}$ depending only on $\beta$ and $\kappa.$
\end{lemma}
\begin{proof}
Note that for all $x\ge 0,$ we have
\[
4x^\kappa \le \beta x + C_\kappa\, \beta^{-\frac{\kappa}{1-\kappa}}
\]
for some $C_\kappa\ge0$ only depending on $\kappa.$ Therefore, denoting $C_{\beta,\kappa} := e^{C_\kappa \beta^{-\frac{\kappa}{1-\kappa}}}$
we have $$
e^{4\NA^\kappa}
\le
C_{\beta,\kappa}\,e^{\beta\NA}=C_{\beta,\kappa}\,e^{\beta(\|V\|-dn)/2}\,e^{\beta H/2},
$$
and so
\[
\Tr\!\bigl(e^{4\NA^\kappa}e^{-\beta H}\bigr)
\le
C_{\beta,\kappa}\,e^{\beta (\|V\|-dn)/2}\,
\Tr\bigl(e^{-\beta H/2}\bigr) \le C_{\beta,\kappa} \,e^{\beta (\|V\|-dn/2)}\left(2\sinh(\beta/2)\right)^{-dn},
 \]
where in the last inequality, we estimated the partition function using Lemma~\ref{lem:PartionFunctionBound}.
Therefore
\begin{align*}
\label{eq:E_GibbsProofSF}
\Tr\bigl(e^{4\NA^\kappa}\sigma_\beta(H)\bigr)
\le
C_{\beta,\kappa} \,e^{\beta(\|V\|-dn/2)} \frac{\left(2\sinh(\beta/2)\right)^{-dn}}{\Tr\left(e^{-\beta H}\right)} \le C_{\beta,\kappa} \,e^{\beta(2\|V\|-dn/2)}  \left(\frac{\sinh(\beta)}{\sinh(\beta/2)}\right)^{dn},
\end{align*}
where in the last inequality, we used the lower bound on the partition function provided by Lemma~\ref{lem:PartionFunctionBound}.
\end{proof}

\begin{theorem}
\label{thm:GibbsPrepareCoulomb}
Let $d\in\{2,3\}$ and $\eps,\beta,\sigma_E>0$  and consider 
\begin{align*}
H_{n,M}\equiv H_{0,n} + W_{n,M} \quad\text{with}\quad M = \operatorname{poly}(n,q(1/\eps)),
\end{align*} 
where $W_{n,M}$ is defined in Section~\ref{sec:HamiltonianTruncation} and where we assume that $\max_{i,j}|\alpha_{n,i,j}| \le \operatorname{poly}(n)$ and $q:\R_+\to \R_+$ denotes some non-decreasing function that is lower bounded as $q(x)\gtrsim \log(x)$.
Denote $\lambda_2 := \operatorname{gap}(L_{\sigma_E,H_{n,M}})>0,$, where the positivity of the gap follows from Theorem~\ref{mainthmpositivegap}.

Then we have that the Gibbs state $\sigma_\beta(H_{n,M})$ can be prepared within $\eps$-trace distance on a quantum computer with $\mathcal{O}\left(n\log (n\,q(1/\eps))\,\log\log(1/\lambda_2)))\right)$ qubits and circuit depth
\begin{align*}
  \widetilde{ \mathcal{O}}\left(\frac{1}{\lambda_2}\operatorname{poly}\left(n, q(1/\eps)\right)\right).
\end{align*}
Hence, for $\alpha^{\max}_n\lesssim 1$ and using Theorem~\ref{thm:trace-norm-perturbation-boundIntro} and $M = \Theta\left(\left(\frac{n^{5/2}}{\eps}\right)^{8d}\right),$ this provides a preparation procedure for $\sigma_\beta(H_n)$ with the same complexities as above corresponding to the specific choice $q(x) = x^{8d}.$ 
Here, the $\mathcal{O},$ $\widetilde{\mathcal{O}}$ and $\Omega$ notations hide constants independent of the displayed parameters, and $\widetilde{\mathcal{O}}$ additionally suppresses subdominant $\operatorname{poly}\log$ factors in $1/\lambda_2$.
\end{theorem}

\begin{proof}


We first use the fact that according to Lemma \ref{lem:product-cutoff-riesz-no-scaling} we have
\begin{align*} \vert W_{n,M}[\psi]\vert &= \vert W_{n}[\Pi_M \psi]\vert \lesssim \langle \Pi_M,(H_{0,n}+1)\Pi_M\rangle\sum_{1\le i<j\le n}|\alpha_{n,i,j}|\le 
\operatorname{poly}(n,M)\\
&=\operatorname{poly}(n,q(1/\eps))  
\end{align*}
and therefore, denoting the finite rank perturbation $V\equiv W_{n,M}$, we have that its operator norm is bounded as
\begin{align}
    \|V\| \le \operatorname{poly}(n,q(1/\eps)).
\end{align}

We prepare the Gibbs state $\sigma_\beta(H_{n,M})$ using a circuit implementation of the Gibbs sampler generated by $\cL_{\sigma_E,H_{n,M}}$, where the relevant complexity analysis has been carried out in \cite[Corollary 4.33]{BeckerRouzeSalzmannToAppearcmp}: 

We take the multi-mode vacuum state $\rho_{\operatorname{ini}}= \kb{0}^{\otimes dn}.$ as the initial state of the circuit. From Lemma~\ref{lem:cScalingFiniteRank} and Lemma~\ref{lem:EGibbsBoundedperturbation}, we see that the constants $\mathfrak{c}$ and $E_{\operatorname{Gibbs}}$ in  \cite[Corollary 4.33]{BeckerRouzeSalzmannToAppearcmp}
satisfy
\begin{align*}
    \mathfrak{c}\,,\, E_{\operatorname{Gibbs}} \le \exp(\operatorname{poly}(n,M,\|V\|)) = \exp(\operatorname{poly}(n,q(1/\eps))).
\end{align*}

The relevant assumptions in \cite[Theorem 4.12 and Corollary 4.33]{BeckerRouzeSalzmannToAppearcmp} on the bare jumps and the Hamiltonian have been verified in Lemma~\ref{lem:VerificationCondFiniteDimImpl}. 
To be precise, we need to note that in \cite[Theorem 4.12 and Corollary 4.33]{BeckerRouzeSalzmannToAppearcmp} the norms estimated in Lemma~\ref{lem:VerificationCondFiniteDimImpl} are only assumed to scale polynomially in $n.$ However, as becomes clear from revisiting \cite[Equation 4.38 \& Propositions 4.7 and 4.11]{BeckerRouzeSalzmannToAppearcmp}, exponential scaling in $n$ can also be dominated by the exponential decay $e^{-\widetilde M^\kappa}$ in the truncation parameter $\widetilde M$ found in Lemma~\ref{lem:VerificationCondFiniteDimImpl}. Hence, we can adjust the truncation level $\widetilde M$ compared to the one considered in \cite[Corollary 4.13]{BeckerRouzeSalzmannToAppearcmp} and take
\begin{align*}
    \widetilde M = \Theta\left(\operatorname{poly}\left(\log\left(\frac{ \exp(\operatorname{poly}(n,q(1/\eps)))\mathfrak{c} E_{\operatorname{Gibbs}} n}{\lambda_2\eps}\right)\right)\right) = \Theta\left(\operatorname{poly}\left(n,\log\left(1/\lambda_2\right),q(1/\eps)\right)\right),
\end{align*}
where in the first equality the $\exp(\operatorname{poly}(n,q(1/\eps)))$ includes the scalings of the norms considered in Lemma~\ref{lem:VerificationCondFiniteDimImpl}.




Using \cite[Remark 4.29 \&  4.30]{BeckerRouzeSalzmannToAppearcmp}, we see that the oracle access to block encodings of $a_i^{\le \widetilde M}/\sqrt{\widetilde M}$ and $\left(a_{i,j}^{\le \widetilde M}\right)^{\dagger}/\sqrt{\widetilde M}$, as well as Hamiltonian simulation $e^{it(H_{n,M})_{\le \widetilde M}}$ required in \cite[Corollary 4.33]{BeckerRouzeSalzmannToAppearcmp}, can both be obtained using a circuit depth of $\mathcal{O}\left(\operatorname{poly}(\log(\widetilde M),\log(1/\eps))\right)$  and \\$\mathcal{O}\left(|t|\operatorname{poly}(\widetilde M,\log(1/\eps))\right)$ respectively.

 We hence apply \cite[Corollary 4.33]{BeckerRouzeSalzmannToAppearcmp} to see that  $\sigma_\beta(H_{n,M})$ can be prepared within $\eps$-trace distance on a quantum computer with $\mathcal{O}(n\log(\widetilde M))=\mathcal{O}\left(n\log (n\,q(1/\eps))\,\log\log(1/\lambda_2)))\right)$ many qubits using a total Hamiltonian simulation time corresponding to the Hamiltonian $(H_{n,M})_{\le \widetilde M}$ of order
\begin{align*}
   \widetilde{ \mathcal{O}}\left(\frac{1}{\lambda_2}\operatorname{poly}\left(n, \log\left(\frac{\exp(\operatorname{poly}(n,q(1/\eps))\mathfrak{c}E_{\operatorname{Gibbs}}}{\eps}\right)\right)\right) =
   \widetilde{ \mathcal{O}}\left(\frac{1}{\lambda_2}\operatorname{poly}\left(n, q(1/\eps)\right)\right),
\end{align*} 
where we again adjusted the required Hamiltonian simulation time by including the \\$\exp(\operatorname{poly}(n,q(1/\eps)))$ in the first equality, in order to account for the scaling derived in Lemma~\ref{lem:VerificationCondFiniteDimImpl}.
Plugging this into the required estimate on the circuit depth for the Hamiltonian simulation of $(H_{n,M})_{\le \widetilde M}$ finishes the proof.
\end{proof}















\subsection{Estimating free energies}
\label{sec:eff_impl}


\noindent We close the paper with a direct application of our results to the estimation of the free energy in Coulomb-interacting systems. While the scheme is introduced here in a simplified framework, the main ideas carry over with little conceptual change to more interesting settings: we consider the Coulomb potential for $d\in\{2,3\}$, denote the Hamiltonians $H_{0,n}$, $H_n$, and $H_{n,M}$ as introduced in Section \ref{perturbationpartitionfunction}, and we denote  $\alpha^{\max}_n:=\max_{1\le i<j\le n}|\alpha_{n,i,j}|$. Our goal is to estimate the free energy $F_\beta(H_{n})$ at some given inverse temperature $\beta>0$. By Theorem \ref{thmfreenregyblabla}, we know that there exists a constant $C_\beta<\infty$, independent of $n,M\ge 1$, such that
\begin{align}
\label{eq:MScalingFreeEnergy}
|F_\beta(H_n)-F_{\beta}(H_{n,M})|
\le \varepsilon\qquad \text{ for }\qquad M=\Omega_\beta\left(\left(\frac{n^3\alpha^{\max}_n+n^{5}(\alpha^{\max}_n)^3}{\varepsilon}\right)^{4d}\right).
\end{align}
Therefore, it is enough to estimate the free energy $F_\beta(H_{n,M})$ of the truncated Hamiltonian $H_{n,M}$. Next, we consider the free energy difference
\begin{align*}
\Delta F_\beta:=F_\beta(H_{n,M})-F_\beta(H_{0,n}).
\end{align*}
Given the path $H(s):=(1-s)H_{0,n}+sH_{n,M}$, $s\in[0,1]$, we have
\begin{align}
\Delta F_\beta=\int_0^1 \frac{d}{ds} F_\beta(H(s))\,ds= \int_0^1\,\Tr\Big(\sigma_\beta(H(s))\Pi_M W_n \Pi_M \Big)\, ds.\label{integralpartitions}
\end{align}
Moreover, denoting $W_{n,M}:=\Pi_M W_n \Pi_M$ and using a standard Riemann-sum approximation,
 \begin{align}
 \label{eq:RiemannFreeEnergy}
\left|\int_0^1\!\Tr\Big(\sigma_\beta(H(s))W_{n,M} \Big)ds\!-\!\sum_{k=0}^{L-1}\frac{\Tr\Big(\sigma_\beta(H(k/L))W_{n,M} \Big)}{L}\right|&\!\lesssim \!\frac{1}{2L}\!\sup_{s\in [0,1]}\Big|\frac{d}{ds}\!\Tr\!\Big(\!\sigma_\beta(H(s))W_{n,M} \!\Big)\!\Big|
 \end{align}
In the next result, we further control the above bound. 
\begin{lemma}
\label{lem:riemann-derivative-bound}
Let
\[
f(s):=\Tr\!\big(\sigma_\beta(H(s))\,W_{n,M}\big),\qquad s\in[0,1].
\]
Then \(f\in C^1([0,1])\), \(f'(s)\le 0\) for all \(s\in[0,1]\), and there exists a constant \(C_\beta<\infty\), independent of \(n\), \(M\) and \(s\), such that
\begin{equation}
\label{eq:derivative-riemann-bound}
\sup_{s\in[0,1]}
\left|
\frac{d}{ds}\Tr\!\big(\sigma_\beta(H(s))\,W_{n,M}\big)
\right|
\lesssim
C_\beta\,\max_{i,j}|\alpha_{n,i,j}|^2\,n^4\,M^{2/d}.
\end{equation}
\end{lemma}

\begin{proof}
Set \(V:=W_{n,M}\). Since \(V\) is bounded, the map
$
s\longmapsto e^{-\beta(H_{0,n}+sV)}
$
is differentiable in trace norm, and therefore \(f\in C^1([0,1])\). By the Duhamel formula,
\[
\frac{d}{ds}e^{-\beta H(s)}
=
-\beta\int_0^1 e^{-(1-u)\beta H(s)}\,V\,e^{-u\beta H(s)}\,du.
\]
Writing \(Z(s):=\Tr(e^{-\beta H(s)})\), we obtain
\begin{align}
f'(s)
&=
\frac{d}{ds}\frac{\Tr(e^{-\beta H(s)}V)}{Z(s)}
\nonumber\\
&=
-\beta\int_0^1
\frac{\Tr\!\big(e^{-(1-u)\beta H(s)}V e^{-u\beta H(s)}V\big)}{Z(s)}\,du
+\beta\left(\frac{\Tr(e^{-\beta H(s)}V)}{Z(s)}\right)^2.
\label{eq:first-derivative-formula}
\end{align}
Equivalently,
\begin{equation}
\label{eq:centered-derivative-formula}
f'(s)
=
-\beta\int_0^1
\Tr\!\Big(
\sigma_\beta(H(s))^{1-u}\big(V-f(s)\big)\sigma_\beta(H(s))^{u}\big(V-f(s)\big)
\Big)\,du.
\end{equation}
Hence \(f'(s)\le 0\). Next, by \eqref{eq:first-derivative-formula} and Hölder's inequality in Schatten spaces,
\[
|f'(s)|
\le
\beta\int_0^1
\left|
\Tr\!\big(\sigma_\beta(H(s))^{1-u}V\sigma_\beta(H(s))^uV\big)
\right|\,du
\le
\beta\,\|V\|^2.
\]
It remains to estimate \(\|V\|\). Let \(\psi\in \mathrm{Ran}(\Pi_M)\). For each pair \(1\le i<j\le n\), the two-body form bound derived in the proof of Lemma \ref{lem:product-cutoff-riesz-no-scaling} gives
\[
\left|
\int w_2(x_i-x_j)|\psi(x)|^2\,dx
\right|
\lesssim
\langle \psi,(h_i+h_j+1)\psi\rangle.
\]
Since \(\psi\in \mathrm{Ran}(\Pi_M)\), one has \(h_k\le \lambda_M\) on the \(k\)-th factor; hence
\[
\langle \psi,(h_i+h_j+1)\psi\rangle
\le
(2\lambda_M+1)\|\psi\|^2
\lesssim
(\lambda_M+1)\|\psi\|^2.
\]
Summing over all pairs, we find $|\langle \psi,V\psi\rangle|
\lesssim
\max_{i,j}|\alpha_{n,i,j}|\binom{n}{2}(\lambda_M+1)\|\psi\|^2$ and therefore,
\begin{equation}
\label{eq:NormTruncatedCoulomb}
\|W_{n,M}\| =\|V\|
\lesssim
\max_{i,j}|\alpha_{n,i,j}|\,n(n-1)\,(\lambda_M+1).
\end{equation}
Combining this with the previous estimate yields
\[
|f'(s)|
\le
C_\beta\,\max_{i,j}|\alpha_{n,i,j}|^2\,n^2(n-1)^2\,(\lambda_M+1)^2,
\]
which, together with \(\lambda_M=\mathcal O(M^{1/d})\), proves \eqref{eq:derivative-riemann-bound}.
\end{proof}

\noindent Thus, combining \eqref{eq:RiemannFreeEnergy}, Lemma~\ref{lem:riemann-derivative-bound} and consider the scaling , we conclude
\begin{align}
\label{eq:RiemannConclusion}
\left|\int_0^1\!\Tr\Big(\sigma_\beta(H(s))W_{n,M} \Big)ds\!-\!\sum_{k=0}^{L-1}\frac{\Tr\Big(\sigma_\beta(H(k/L))W_{n,M} \Big)}{L}\right|&\!\lesssim \frac{1}{2L}C_\beta\,(\alpha^{\max}_n)^2 n^4\,M^{2/d}\le \varepsilon
\end{align}
for $L=\Omega_\beta\Big((\alpha^{\max}_n)^2n^4M^{2/d}\varepsilon^{-1} \Big)=\Omega_\beta\left(\left(n^{7/2}(\alpha^{\max}_n)^{5/4} + n^{11/2}(\alpha^{\max}_n)^{13/4}\right)^8\varepsilon^{-9}\right),$ where for the last equality we have used the scaling choice \eqref{eq:MScalingFreeEnergy}.

\begin{theorem}[Quantum algorithm for free energy estimation]
\label{thm:FreeEnergyCoulombQuantumAlgorithm}
The free energy $F_\beta(H_n)$ can be estimated with accuracy $\eps>0$ and a probability of failure bounded by $\delta>0$ on a quantum computer with $\mathcal{O}\left(n\log( n/\eps)\,\log\log(1/\lambda^{\min}_{2})\right)$ many qubits, with a total runtime of order 
\begin{align*}
    \widetilde{ \mathcal{O}}\left(\frac{1}{\lambda^{\min}_{2}}\log\left(1/\delta\right)\operatorname{poly}\left(n\,,\,1/\eps\right)\right)
\end{align*}
where $\lambda^{\min}_{2}:= \min_{s\in[0,1]} \lambda_2(s)>0$, $\lambda_2(s)\equiv\operatorname{gap}(L_{\sigma_E,H(s)})>0$, $H(s):=(1-s)H_{0,n}+sH_{n,M}$, with $M=\Theta\left(\left(\frac{n^3\alpha^{\max}_n+n^{5}(\alpha^{\max}_n)^3}{\varepsilon}\right)^{4d}\right),$ $\alpha^{\max}_n\le\operatorname{poly}(n)$ and where positivity of the spectral gap\footnote{Theorem~\ref{mainthmpositivegap} was only stated for $H_{n,M}.$ However, by replacing $\alpha_{n,i,j}$ with $s\alpha_{n,i,j},$, this shows positivity of the gap for all $s\in[0,1].$} follows by Theorem~\ref{mainthmpositivegap}.
 Moreover, the $\widetilde{\mathcal{O}}$ notation hides constants independent of the displayed parameters and additionally suppresses subdominant $\operatorname{poly}\log$ factors in the leading order.
\end{theorem}
\begin{proof}

By Theorem \ref{thmfreenregyblabla}, we know we can achieve 
\[
|F_\beta(H_n)-F_{\beta}(H_{n,M})|
\le \varepsilon/4.
\]
with \begin{align}
\label{eq:MScalingProooof}
    M = \Theta\left(\left(\frac{n^3\alpha^{\max}_n+n^{5}(\alpha^{\max}_n)^3}{\varepsilon}\right)^{4d}\right)\le \operatorname{poly}(n,1/\eps).
\end{align}
Therefore, it is enough to estimate the free energy $F_\beta(H_{n,M})$ of the truncated Hamiltonian $H_{n,M}$. Note, further, that we can analytically compute $F_\beta(H_{0,n})$ as 
\begin{align*}
\Tr\left(e^{-\beta H_{0,n}}\right) = \left(\sum_{k=0}^\infty e^{-\beta(2k+1)}\right)^{dn}=\left(2\sinh(\beta)\right)^{-dn}\quad\text{and}\quad F_\beta(H_{0,n}) = \frac{dn}{\beta}\,\log\left(2\sinh(\beta)\right).
\end{align*}
Hence, in order to estimate $F_\beta(H_{n,M}) = \Delta F_\beta + F_\beta(H_{0,n}),$ it remains to estimate the difference of free energies $\Delta F_\beta:$

We write $W_{n,M} = \sum_{1\le i<j\le n} \alpha_{n,i,j}\Pi_MW_{i,j} \Pi_M$ and note that by Lemma~\ref{lem:product-cutoff-riesz-no-scaling}, we have $\|W_{n,M}\| \le \operatorname{poly}(n,M)$, and by an analogous argument for $i,j$
fixed $\|\Pi_MW_{i,j}\Pi_M\| \lesssim \operatorname{poly}(n,M).$

From this, and Hoeffding's inequality, we see that for all $k=0,\cdots,L-1$, we can estimate the expectation value $\Tr(\sigma_\beta(k/L)\Pi_MW_{i,j}\Pi_M)$ with accuracy $\eps/(4n^2)$ and probability of failure $\delta>0$ using a number of samples of $\sigma_\beta(H(k/L))$ and measurements in the eigenbasis of $\Pi_MW_{i,j}\Pi_M$ of order
 \begin{align}
 \label{eq:SampleComplexitySigma1}
\mathcal{O}\left(\,\frac{n^4\|\Pi_MW_{i,j}\Pi_M\|^2}{\eps^2}\log\left(\frac{1}{\delta}\right)\,\right) = \mathcal{O}\left(\,\frac{\operatorname{poly}(n,M)}{\eps^2}\log\left(\frac{1}{\delta}\right)\,\right).
 \end{align}
Therefore, using the union bound, we can estimate 
\begin{align}
\label{eq:Riemann}
   \frac{1}{L}\sum_{k=0}^{L-1}\Tr(\sigma_\beta(H(k/L)) W_{n,M}) =  \frac{1}{L}\sum_{k=0}^{L-1}\sum_{1\le i<j\le n}\alpha_{n,i,j}\Tr(\sigma_\beta(H(k/L)) \Pi_M W_{i,j})
\end{align} 
 with accuracy $\eps/4$ and probability of failure bounded by $\delta$ using 
 \begin{align}
\label{eq:SampleComplexitySigmaL}
      \mathcal{O}\left(\,\frac{L \operatorname{poly}(n,M)}{\eps^2}\log\left(\frac{L}{\delta}\right)\,\right).
 \end{align}
 many samples of the Gibbs states and measurements in the eigenbasis of $\Pi_MW_{i,j}\Pi_M$ for different $i,j$. Using \eqref{integralpartitions} and \eqref{eq:RiemannConclusion}, we can estimate $\Delta F_\beta$ by \eqref{eq:Riemann} with accuracy $\eps/4$ by taking 
 \begin{align}
 \label{eq:ScalingLandM}
     L = \operatorname{poly}(n,1/\eps) 
 \end{align}
 Lastly, using Theorem~\ref{thm:GibbsPrepareCoulomb} for varying\footnote{That is for $k=0,\cdots,L-1$ we pick $\alpha^{(k)}_{n,i,j} = \frac{k}{L}\alpha_{n,i,j}$ to prepare $\sigma_\beta(H(k/L)).$} $\alpha_{n,i,j}$, we can prepare for each $k=0,\cdots, L-1$ the Gibbs state $\sigma_\beta(H(k/L))$  with accuracy $\eps/(4\|W_{n,M}\|)$ in trace distance on a quantum computer with \\ $\mathcal{O}\left(n\log( n/\eps)\,\log\log(1/\lambda_2(k/L))\right)\le\mathcal{O}\left(n\log( n/\eps)\,\log\log(1/(\lambda^{\min}_2))\right) $ many qubits using a circuit depth of order
 \begin{align*}
\widetilde{ \mathcal{O}}\left(\frac{1}{\lambda_2(k/L)}\operatorname{poly}\left(n, 1/\eps\right)\right) \le \widetilde{ \mathcal{O}}\left(\frac{1}{\lambda^{\min}_2}\operatorname{poly}\left(n, 1/\eps\right)\right).
 \end{align*}
 Multiplying this with the required numbers of samples of the Gibbs states in \eqref{eq:SampleComplexitySigmaL} and inserting the scaling of $M$ and $L$ in \eqref{eq:MScalingProooof} and \eqref{eq:ScalingLandM} yields the result.
Note that, since \(W_{i,j}\) is a two-particle interaction, \(\Pi_M W_{i,j}\Pi_M\) acts nontrivially only on two \(M\)-dimensional registers. Therefore, the required measurements in the eigenbasis of \(\Pi_M W_{i,j}\Pi_M\) can be realized by an additional quantum circuit consisting solely of 2-qubit gates and having depth \(\mathcal{O}(M^2)=\mathcal{O}(\operatorname{poly}(n,1/\eps))\) (see e.g. \cite[Section~4.5.1]{Nielsen_Chuang_2010}), followed by computational basis measurements. Hence, the measurement step requires circuit depth of the same order as the previous steps.



\end{proof}










\smallsection{Acknowledgement}
SB would also like to acknowledge support from the SNF Grant PZ00P2\_216019. 
CR is supported by France 2030 under the
French National Research Agency award number ''ANR-22-EXES-0013''. RS acknowledges support by the European Research Council (ERC Grant Agreement No.~948139) and from the Excellence Cluster Matter and Light for Quantum Computing (ML4Q-2).



\bibliographystyle{amsplain}
\bibliography{ref}



\end{document}
